% Options for packages loaded elsewhere
\PassOptionsToPackage{unicode}{hyperref}
\PassOptionsToPackage{hyphens}{url}
%
\documentclass[
  12pt
]{book}
\usepackage{amsmath,amssymb}
\usepackage[left=2cm,right=2cm]{geometry}
\usepackage{iftex}
\ifPDFTeX
  \usepackage[T1]{fontenc}
  \usepackage[utf8]{inputenc}
  \usepackage{textcomp} % provide euro and other symbols
\else % if luatex or xetex
  \usepackage{unicode-math} % this also loads fontspec
  \defaultfontfeatures{Scale=MatchLowercase}
  \defaultfontfeatures[\rmfamily]{Ligatures=TeX,Scale=1}
\fi
\usepackage{lmodern}
\ifPDFTeX\else
  \IfFontExistsTF{Bitstream Vera Sans}{
    \setmainfont{Bitstream Vera Sans}[Scale=MatchLowercase]
  }{}
  \IfFontExistsTF{Cascadia Code}{
    \setmonofont{Cascadia Code}[
      Scale=MatchLowercase,
      BoldFont={Cascadia Code:style=Bold},
      ItalicFont={Cascadia Code:style=Italic},
      BoldItalicFont={Cascadia Code:style=Bold Italic}
    ]
  }{}
\fi
% Use upquote if available, for straight quotes in verbatim environments
\IfFileExists{upquote.sty}{\usepackage{upquote}}{}
\IfFileExists{microtype.sty}{% use microtype if available
  \usepackage[]{microtype}
  \UseMicrotypeSet[protrusion]{basicmath} % disable protrusion for tt fonts
}{}
\makeatletter
\@ifundefined{KOMAClassName}{% if non-KOMA class
  \IfFileExists{parskip.sty}{%
    \usepackage{parskip}
  }{% else
    \setlength{\parindent}{0pt}
    \setlength{\parskip}{6pt plus 2pt minus 1pt}}
}{% if KOMA class
  \KOMAoptions{parskip=half}}
\makeatother
\usepackage{xcolor}
\usepackage{color}
\usepackage{fancyvrb}
\newcommand{\VerbBar}{|}
\newcommand{\VERB}{\Verb[commandchars=\\\{\}]}
\DefineVerbatimEnvironment{Highlighting}{Verbatim}{commandchars=\\\{\}}
% Add ',fontsize=\small' for more characters per line
\newenvironment{Shaded}{}{}
\newcommand{\AlertTok}[1]{\textcolor[rgb]{1.00,0.00,0.00}{\textbf{#1}}}
\newcommand{\AnnotationTok}[1]{\textcolor[rgb]{0.38,0.63,0.69}{\textbf{\textit{#1}}}}
\newcommand{\AttributeTok}[1]{\textcolor[rgb]{0.49,0.56,0.16}{#1}}
\newcommand{\BaseNTok}[1]{\textcolor[rgb]{0.25,0.63,0.44}{#1}}
\newcommand{\BuiltInTok}[1]{\textcolor[rgb]{0.00,0.50,0.00}{#1}}
\newcommand{\CharTok}[1]{\textcolor[rgb]{0.25,0.44,0.63}{#1}}
\newcommand{\CommentTok}[1]{\textcolor[rgb]{0.38,0.63,0.69}{\textit{#1}}}
\newcommand{\CommentVarTok}[1]{\textcolor[rgb]{0.38,0.63,0.69}{\textbf{\textit{#1}}}}
\newcommand{\ConstantTok}[1]{\textcolor[rgb]{0.53,0.00,0.00}{#1}}
\newcommand{\ControlFlowTok}[1]{\textcolor[rgb]{0.00,0.44,0.13}{\textbf{#1}}}
\newcommand{\DataTypeTok}[1]{\textcolor[rgb]{0.56,0.13,0.00}{#1}}
\newcommand{\DecValTok}[1]{\textcolor[rgb]{0.25,0.63,0.44}{#1}}
\newcommand{\DocumentationTok}[1]{\textcolor[rgb]{0.73,0.13,0.13}{\textit{#1}}}
\newcommand{\ErrorTok}[1]{\textcolor[rgb]{1.00,0.00,0.00}{\textbf{#1}}}
\newcommand{\ExtensionTok}[1]{#1}
\newcommand{\FloatTok}[1]{\textcolor[rgb]{0.25,0.63,0.44}{#1}}
\newcommand{\FunctionTok}[1]{\textcolor[rgb]{0.02,0.16,0.49}{#1}}
\newcommand{\ImportTok}[1]{\textcolor[rgb]{0.00,0.50,0.00}{\textbf{#1}}}
\newcommand{\InformationTok}[1]{\textcolor[rgb]{0.38,0.63,0.69}{\textbf{\textit{#1}}}}
\newcommand{\KeywordTok}[1]{\textcolor[rgb]{0.00,0.44,0.13}{\textbf{#1}}}
\newcommand{\NormalTok}[1]{#1}
\newcommand{\OperatorTok}[1]{\textcolor[rgb]{0.40,0.40,0.40}{#1}}
\newcommand{\OtherTok}[1]{\textcolor[rgb]{0.00,0.44,0.13}{#1}}
\newcommand{\PreprocessorTok}[1]{\textcolor[rgb]{0.74,0.48,0.00}{#1}}
\newcommand{\RegionMarkerTok}[1]{#1}
\newcommand{\SpecialCharTok}[1]{\textcolor[rgb]{0.25,0.44,0.63}{#1}}
\newcommand{\SpecialStringTok}[1]{\textcolor[rgb]{0.73,0.40,0.53}{#1}}
\newcommand{\StringTok}[1]{\textcolor[rgb]{0.25,0.44,0.63}{#1}}
\newcommand{\VariableTok}[1]{\textcolor[rgb]{0.10,0.09,0.49}{#1}}
\newcommand{\VerbatimStringTok}[1]{\textcolor[rgb]{0.25,0.44,0.63}{#1}}
\newcommand{\WarningTok}[1]{\textcolor[rgb]{0.38,0.63,0.69}{\textbf{\textit{#1}}}}
\usepackage{longtable,booktabs,array}
\usepackage{calc} % for calculating minipage widths
% Correct order of tables after \paragraph or \subparagraph
\usepackage{etoolbox}
\makeatletter
\patchcmd\longtable{\par}{\if@noskipsec\mbox{}\fi\par}{}{}
\makeatother
% Allow footnotes in longtable head/foot
\IfFileExists{footnotehyper.sty}{\usepackage{footnotehyper}}{\usepackage{footnote}}
\makesavenoteenv{longtable}
\setlength{\emergencystretch}{3em} % prevent overfull lines
\providecommand{\tightlist}{%
  \setlength{\itemsep}{0pt}\setlength{\parskip}{0pt}}
\setcounter{secnumdepth}{-\maxdimen} % remove section numbering
\ifLuaTeX
\usepackage[bidi=basic]{babel}
\else
\usepackage[bidi=default]{babel}
\fi
\babelprovide[main,import]{ngerman}
% get rid of language-specific shorthands (see #6817):
\let\LanguageShortHands\languageshorthands
\def\languageshorthands#1{}
\ifLuaTeX
  \usepackage{selnolig}  % disable illegal ligatures
\fi
\IfFileExists{bookmark.sty}{\usepackage{bookmark}}{\usepackage{hyperref}}
\IfFileExists{xurl.sty}{\usepackage{xurl}}{} % add URL line breaks if available
\urlstyle{same}
\hypersetup{
  pdftitle={OndselSolver für FreeCAD},
  pdflang={de},
  hidelinks,
  pdfcreator={LaTeX}}

\title{OndselSolver für FreeCAD}
\usepackage{etoolbox}
\makeatletter
\providecommand{\subtitle}[1]{% add subtitle to \maketitle
  \apptocmd{\@title}{\par {\large #1 \par}}{}{}
}
\makeatother
\subtitle{Ein lebendes Buch über den Assembly-Solver}
\author{}
\date{}

\begin{document}
\frontmatter
\maketitle

\renewcommand*\contentsname{Inhaltsverzeichnis}
{
\setcounter{tocdepth}{1}
\tableofcontents
}
\mainmatter
\hypertarget{kap-arbeitsprinzipien}{%
\chapter*{Arbeitsprinzipien dieses Buches}\label{kap-arbeitsprinzipien}}
\addcontentsline{toc}{chapter}{Arbeitsprinzipien dieses Buches}

\hypertarget{warum-dieses-buch-als-lebendes-projekt-aufgebaut-ist}{%
\section*{Warum dieses Buch als lebendes Projekt aufgebaut
ist}\label{warum-dieses-buch-als-lebendes-projekt-aufgebaut-ist}}
\addcontentsline{toc}{section}{Warum dieses Buch als lebendes Projekt
aufgebaut ist}

Ein Buch über einen aktiven Solver in einer aktiven CAD-Umgebung muss
mit dem Projekt wachsen. Das gilt besonders für FreeCAD und
OndselSolver, weil sich Architektur, Bugs, APIs und
Modellierungsprinzipien über Zeit verändern. Ein einmal fertiges Buch
würde schnell veralten.

Darum gilt hier:

\begin{itemize}
\tightlist
\item
  Das Buch wird nicht komplett neu geschrieben.
\item
  Kapitel werden ergänzt, korrigiert und vertieft.
\item
  Neue Erkenntnisse werden in die vorhandene Struktur eingefügt.
\item
  Jeder Abschnitt erhält eine sachliche, überprüfbare Grundlage.
\end{itemize}

\hypertarget{vorgehensweise-der-dokumentation}{%
\section*{Vorgehensweise der
Dokumentation}\label{vorgehensweise-der-dokumentation}}
\addcontentsline{toc}{section}{Vorgehensweise der Dokumentation}

Für dieses Projekt gilt ein einfacher Standard:

\begin{itemize}
\tightlist
\item
  Beschreibung des Sachverhalts auf verständlicher Ebene
\item
  Einordnung in die reale FreeCAD-/OndselSolver-Architektur
\item
  Beispiel oder kleine Reproduktion
\item
  Hinweise, was in diesem Kontext relevant ist
\item
  Verknüpfung mit den echten Dateien und Modulen im Sandbox-Projekt
\end{itemize}

\hypertarget{grundidee-von-ondselsolver}{%
\section*{Grundidee von OndselSolver}\label{grundidee-von-ondselsolver}}
\addcontentsline{toc}{section}{Grundidee von OndselSolver}

OndselSolver modelliert Mechanik als Beziehungssystem:

\begin{itemize}
\tightlist
\item
  Teile haben Positionen und Orientierungen.
\item
  Gelenke, Kontakte und Constraints beschreiben Beziehungen zwischen
  diesen Positionen.
\item
  Der Solver löst diese Beziehungen so, dass das System konsistent ist.
\end{itemize}

Das ist nicht dasselbe wie ein „Baukastensystem mit fertigen
3D-Schnitten``. Vielmehr ist der Solver ein numerischer Gleichungslöser,
der geometrische und kinematische Bedingungen erfüllt.

\hypertarget{wichtiges-verstuxe4ndnis-fuxfcr-leser}{%
\section*{Wichtiges Verständnis für
Leser}\label{wichtiges-verstuxe4ndnis-fuxfcr-leser}}
\addcontentsline{toc}{section}{Wichtiges Verständnis für Leser}

Ein neuer Leser muss die folgenden Grundbegriffe verstehen:

\begin{itemize}
\tightlist
\item
  Part: ein Bauteil oder ein physischer Körper
\item
  Joint: Verbindung zwischen Teilen
\item
  Constraint: geometrische oder kinematische Bedingung
\item
  Solve: Berechnung einer konsistenten Konfiguration
\item
  Residuum: Restfehler einer Gleichung oder Beziehung
\item
  Iteration: wiederholte Annäherung an eine Lösung
\end{itemize}

\hypertarget{wie-dieses-buch-mit-realcode-verknuxfcpft-wird}{%
\section*{Wie dieses Buch mit Realcode verknüpft
wird}\label{wie-dieses-buch-mit-realcode-verknuxfcpft-wird}}
\addcontentsline{toc}{section}{Wie dieses Buch mit Realcode verknüpft
wird}

Die Praxis ist wichtiger als bloße Theorie. Dieses Buch bezieht sich
daher auf die im Sandbox-Projekt vorhandenen Quellstellen, insbesondere:

\begin{itemize}
\tightlist
\item
  FreeCAD-Assembly-Struktur
\item
  OndselSolver-Quellcode im dritten Party Bereich
\item
  Aufgaben und Bugreports aus dem Sandbox
\end{itemize}

Das bedeutet: Wenn ein Kapitel etwas über die Kommunikation zwischen
FreeCAD und OndselSolver sagt, dann verweist es auf die Architektur der
Baugruppenmodul-Implementierung und auf die typische Datenflusslogik,
nicht nur auf allgemeine Theorie.

\hypertarget{regel-fuxfcr-erguxe4nzungen}{%
\section*{Regel für Ergänzungen}\label{regel-fuxfcr-erguxe4nzungen}}
\addcontentsline{toc}{section}{Regel für Ergänzungen}

Wenn ein neues Kapitel hinzugefügt wird, gilt:

\begin{itemize}
\tightlist
\item
  bestehende Kapitel bleiben unverändert, soweit nicht klar ob
  vervollständigung notwendig ist
\item
  neue Erkenntnisse werden an das Ende oder in eine passende Ergänzung
  eingefügt
\item
  bei Fehlern wird der Text korrigiert, aber nicht in einem „Neuaufbau``
  der gesamten Datei verworfen
\end{itemize}

\hypertarget{ziel-dieses-buches}{%
\section*{Ziel dieses Buches}\label{ziel-dieses-buches}}
\addcontentsline{toc}{section}{Ziel dieses Buches}

Das Buch soll das Verständnis schaffen, nicht nur einfach Fakten
aufzuzählen. Ein Lernender soll am Ende verstehen:

\begin{itemize}
\tightlist
\item
  Was ein Solver macht
\item
  Wie ein CAD-System ihm Daten liefert
\item
  Wie das Ergebnis zurück in die 3D-Geometrie übernommen wird
\item
  Warum FreeCAD und OndselSolver in einem gemeinsamen Architekturmodell
  verbunden sind
\end{itemize}

Dieses Buch ist somit ein langfristiger Lern- und Arbeitsdokument.

\hypertarget{grundlagen-freecad-und-asmt}{%
\chapter{\texorpdfstring{{Teil I --} Grundlagen, FreeCAD und
ASMT}{Teil I -- Grundlagen, FreeCAD und ASMT}}\label{grundlagen-freecad-und-asmt}}

\hypertarget{kap-was-ist-ondselsolver}{%
\chapter{\texorpdfstring{{1} Was ist
OndselSolver?}{1 Was ist OndselSolver?}}\label{kap-was-ist-ondselsolver}}

\hypertarget{ein-kurzer-satz}{%
\section{\texorpdfstring{{1.1} Ein kurzer
Satz}{1.1 Ein kurzer Satz}}\label{ein-kurzer-satz}}

OndselSolver ist ein Mehrkörper- und Assembliesolver. Er berechnet, wie
Teile in einer Konstruktion zueinander liegen müssen, wenn bestimmte
geometrische Beziehungen, Gelenke, Abstände oder Winkel eingehalten
werden sollen.

In einem CAD-Kontext ist das besonders wichtig für Baugruppen,
Mechanismen und Kinematik. FreeCAD kann Teile und Beziehungen
modellieren; der Solver findet dann eine konsistente Lage der Teile,
wenn diese Beziehungen erfüllt sein müssen.

\hypertarget{warum-ein-solver-uxfcberhaupt-nuxf6tig-ist}{%
\section{\texorpdfstring{{1.2} Warum ein Solver überhaupt nötig
ist}{1.2 Warum ein Solver überhaupt nötig ist}}\label{warum-ein-solver-uxfcberhaupt-nuxf6tig-ist}}

Ohne Solver wäre eine Baugruppe nur eine Menge von Objekten mit
Positionen. Der Benutzer kann sie manuell bewegen, aber das System kennt
die Beziehungen zwischen den Teilen nicht in numerisch umsetzbarer Form.

Ein Solver bringt genau diese logische Schicht ein:

\begin{itemize}
\tightlist
\item
  „Teil A ist mit Teil B über eine Drehachse verbunden``
\item
  „Teil C muss auf der Ebene von Teil D liegen``
\item
  „Die Distanz zwischen den Punkten muss 45 mm betragen``
\item
  „Die Drehung von Teil X muss mit Translation von Teil Y gekoppelt
  sein``
\end{itemize}

Der Solver nimmt diese Bedingungen, modelliert sie als Gleichungen und
sucht nach einer konsistenten Lösung.

\hypertarget{ondselsolver-in-der-freecad-welt}{%
\section{\texorpdfstring{{1.3} OndselSolver in der
FreeCAD-Welt}{1.3 OndselSolver in der FreeCAD-Welt}}\label{ondselsolver-in-der-freecad-welt}}

In der FreeCAD-Assembly-Architektur wirkt OndselSolver als die
Berechnungsebene hinter der Baugruppenlogik. Der Benutzersicht bleibt
dabei meist die 3D-Baugruppe und die Joint-Definitionen erhalten. Unter
der Oberfläche übernimmt der Solver die Berechnung der Positionen und
Orientierungen.

Das ist entscheidend: FreeCAD stellt die Daten bereit, OndselSolver
verarbeitet sie. Danach kann FreeCAD die berechneten Platzierungen
wieder in die 3D-Geometrie einsetzen.

\hypertarget{typische-aufgaben-eines-solvers}{%
\section{\texorpdfstring{{1.4} Typische Aufgaben eines
Solvers}{1.4 Typische Aufgaben eines Solvers}}\label{typische-aufgaben-eines-solvers}}

Ein Assembly-Solver muss typischerweise diese Dinge leisten:

\begin{itemize}
\tightlist
\item
  bekannte Ausgangslage prüfen
\item
  Freiheitsgrade erkennen
\item
  Constraints und Gelenke auflösen
\item
  fehlende oder redundante Bedingungen entdecken
\item
  zu einer konsistenten Lage konvergieren
\item
  Ergebnisse in Form neuer Platzierungen zurückgeben
\end{itemize}

\hypertarget{kinematik-vs.-dynamik}{%
\section{\texorpdfstring{{1.5} Kinematik
vs.~Dynamik}{1.5 Kinematik vs.~Dynamik}}\label{kinematik-vs.-dynamik}}

Ein wichtiger Punkt: Ein Solver für Baugruppen ist nicht automatisch nur
eine „Bewegungssimulation`` im Sinne einer vollständigen Dynamik mit
Kräften, Massen und Zeitintegration.

Die Grundidee ist oft:

\begin{itemize}
\tightlist
\item
  Kinematik: Wie liegen die Teile zueinander, wenn die Beziehungen
  erfüllt werden?
\item
  Dynamik: Wie bewegen sie sich über Zeit unter Kräften, Trägheit und
  Bewegungsgesetzen?
\end{itemize}

OndselSolver kann in diesem Zusammenhang als Mehrkörper-Mechanik-Engine
verstanden werden: Sie arbeitet mit geometrischen und mechanischen
Constraints, nicht nur mit bloßer Geometrie-Objektplatzierung.

\hypertarget{warum-das-fuxfcr-freecad-relevant-ist}{%
\section{\texorpdfstring{{1.6} Warum das für FreeCAD relevant
ist}{1.6 Warum das für FreeCAD relevant ist}}\label{warum-das-fuxfcr-freecad-relevant-ist}}

Für FreeCAD ist das besonders wichtig, weil Baugruppen oft aus mehreren
Teilen bestehen, die eine funktionale Beziehung zueinander besitzen.
Beispiele:

\begin{itemize}
\tightlist
\item
  Schwenkarm und Gelenk
\item
  Lager und Ausrichtung
\item
  Korrekte Position von Bohrungen oder Schnitten
\item
  Kinematische Kette mit mehreren beweglichen Gliedern
\end{itemize}

Der Solver macht aus diesen Beziehungen eine technische, berechenbare
Struktur. Das ist der Kern dessen, was in FreeCAD als Assembly-Lösung
bezeichnet wird.

\hypertarget{ein-einfaches-alltagsschema}{%
\section{\texorpdfstring{{1.7} Ein einfaches
Alltagsschema}{1.7 Ein einfaches Alltagsschema}}\label{ein-einfaches-alltagsschema}}

Wenn ein Benutzer eine Baugruppe zusammenbaut, dann beschreibt er meist
intuitiv:

\begin{itemize}
\tightlist
\item
  „Diese zwei Teile sollen fest verbunden sein.``
\item
  „Diese Linie und diese Achse sollen aufeinander liegen.``
\item
  „Dieser Winkel soll 90° betragen.``
\end{itemize}

Der Solver verwandelt diese Aussagen in formalisierte Gleichungen. Dann
löst er sie im Sinne eines Systems aus miteinander verbundenen
Bedingungen.

\hypertarget{die-zentrale-idee-in-einem-satz}{%
\section{\texorpdfstring{{1.8} Die zentrale Idee in einem
Satz}{1.8 Die zentrale Idee in einem Satz}}\label{die-zentrale-idee-in-einem-satz}}

OndselSolver ist die Recheneinheit, die in FreeCAD aus geometrischen
Relationen eine konsistente mechanische Konfiguration erzeugt.

Damit ist es nicht nur ein Spezialwerkzeug, sondern ein grundlegender
Baustein, der das Denken in „Objekten`` in das Denken in „Beziehungen``
übersetzt.

\hypertarget{wichtige-lernfrage}{%
\section{\texorpdfstring{{1.9} Wichtige
Lernfrage}{1.9 Wichtige Lernfrage}}\label{wichtige-lernfrage}}

Der Leser sollte sich im Verlauf des Buches immer fragen:

\begin{itemize}
\tightlist
\item
  Welche Beziehung beschreibt dieser Joint oder Constraint?
\item
  Welches Teil hat welchen Freiheitsgrad?
\item
  Welche Variable wird vom Solver gelöst?
\item
  Was ist die Ausgangsbedingung und welche Resultate sind akzeptabel?
\end{itemize}

Diese Fragen sind der Schlüssel zum Verständnis von OndselSolver.

\hypertarget{kap-wie-funktioniert-ondselsolver}{%
\chapter{\texorpdfstring{{2} Wie funktioniert
OndselSolver?}{2 Wie funktioniert OndselSolver?}}\label{kap-wie-funktioniert-ondselsolver}}

\hypertarget{grundprinzip}{%
\section{\texorpdfstring{{2.1}
Grundprinzip}{2.1 Grundprinzip}}\label{grundprinzip}}

OndselSolver arbeitet mit einem klassischen Muster: Er nimmt ein
Mechanikmodell, übersetzt es in Gleichungen und löst dann eine oder
mehrere konsistente Konfigurationen. Das geschieht üblicherweise nicht
in einem Schritt „per Hand``, sondern iterativ.

Die wichtigsten Schritte sind:

\begin{enumerate}
\def\labelenumi{\arabic{enumi}.}
\tightlist
\item
  Modellierung der Teile und Beziehungen
\item
  Erzeugung der Gleichungssysteme aus den Constraints
\item
  Berechnung eines Näherungswertes
\item
  Iterative Verfeinerung bis zur Lösung
\item
  Rückgabe der berechneten Platzierungen
\end{enumerate}

\hypertarget{teile-gelenke-und-constraints}{%
\section{\texorpdfstring{{2.2} Teile, Gelenke und
Constraints}{2.2 Teile, Gelenke und Constraints}}\label{teile-gelenke-und-constraints}}

Ein Mehrkörpersystem besteht aus mehreren Teilen. Sie bewegen sich
relativ zueinander, aber nicht frei. Ihre Bewegung wird durch
Verbindungen eingeschränkt.

Beispiele:

\begin{itemize}
\tightlist
\item
  Festes Gelenk
\item
  Drehgelenk
\item
  Schiebegelenk
\item
  Planare Beziehung
\item
  Punkt-auf-Linie
\item
  Winkel- oder Abstandsbeschränkung
\end{itemize}

Diese Beziehungen erzeugen sogenannte Constraints. Ein Constraint ist
eine mathematische Bedingung, mit der die mögliche Bewegung der Teile
eingeschränkt wird.

\hypertarget{freiheitsgrade}{%
\section{\texorpdfstring{{2.3}
Freiheitsgrade}{2.3 Freiheitsgrade}}\label{freiheitsgrade}}

Ein Teil im Raum hat normalerweise sechs Freiheitsgrade:

\begin{itemize}
\tightlist
\item
  drei Translationen
\item
  drei Rotationen
\end{itemize}

Ein Joint oder Constraint nimmt davon einige weg. Ein festes Gelenk
entfernt praktisch alle Freiheitsgrade, ein Drehgelenk nur eine
Rotation, ein Schiebegelenk nur eine Translation. So entsteht ein
System, das bei der korrekten Kombination stabil und lösbar wird.

\hypertarget{warum-iteration-nuxf6tig-ist}{%
\section{\texorpdfstring{{2.4} Warum Iteration nötig
ist}{2.4 Warum Iteration nötig ist}}\label{warum-iteration-nuxf6tig-ist}}

Die meisten realen mechanischen Beziehungen sind nicht linear. Das
bedeutet:

\begin{itemize}
\tightlist
\item
  Die Positionsänderung eines Teils beeinflusst andere Bedingungen
  gleichzeitig.
\item
  Die Lösung kann nicht immer geschlossen per Formel bestimmt werden.
\item
  Man muss mit einem Näherungsprozess starten und sich schrittweise
  verbessern.
\end{itemize}

Das ist der Kern der numerischen Lösung: Start mit einer geschätzten
Konfiguration, berechne Fehler, korrigiere Positionen, wiederhole bis
der Fehler klein genug ist.

\hypertarget{newton-verfahren-als-einfaches-denkmodell}{%
\section{\texorpdfstring{{2.5} Newton-Verfahren als einfaches
Denkmodell}{2.5 Newton-Verfahren als einfaches Denkmodell}}\label{newton-verfahren-als-einfaches-denkmodell}}

Ein sehr verbreiteter Gedanke in Solvern ist das Newton-Verfahren oder
ähnliche iterative Verfahren:

\begin{itemize}
\tightlist
\item
  Definiere eine Fehlermenge
\item
  Berechne die Ableitung des Fehlers
\item
  Korrigiere die Variablen in Richtung kleinerer Fehler
\item
  Wiederhole, bis der Fehler unter einem Toleranzwert liegt
\end{itemize}

In einfachen Worten:

\begin{itemize}
\tightlist
\item
  Der Solver schaut, wie groß der Fehler des aktuellen Zustands ist.
\item
  Dann schätzt er, in welche Richtung sich die beweglichen Größen
  verändern müssen.
\item
  Danach wiederholt er das, bis die Lösung „gut genug`` ist.
\end{itemize}

\hypertarget{wichtiger-unterschied-zwischen-geometrisch-korrekt-und-physikalisch-vollstuxe4ndig}{%
\section{\texorpdfstring{{2.6} Wichtiger Unterschied zwischen
„geometrisch korrekt`` und „physikalisch
vollständig``}{2.6 Wichtiger Unterschied zwischen „geometrisch korrekt`` und „physikalisch vollständig``}}\label{wichtiger-unterschied-zwischen-geometrisch-korrekt-und-physikalisch-vollstuxe4ndig}}

Ein Solver kann eine konsistente geometrische Lösung finden, ohne dass
das gesamte mechanische System mit Massen, Kräften und Simulationsdauer
vollständig modelliert ist. Das ist bei Assemblies oft genau der
gewünschte Einsatz.

Der Fokus liegt dann auf:

\begin{itemize}
\tightlist
\item
  Stabilität der Konfiguration
\item
  Konsistenz der Beziehungen
\item
  Erfüllung der Constraints
\end{itemize}

Nicht immer auf physikalische Dynamik im engeren Sinne.

\hypertarget{so-sieht-ein-typischer-ablauf-aus}{%
\section{\texorpdfstring{{2.7} So sieht ein typischer Ablauf
aus}{2.7 So sieht ein typischer Ablauf aus}}\label{so-sieht-ein-typischer-ablauf-aus}}

Ein typischer Solve-Lauf läuft in etwa so:

\begin{enumerate}
\def\labelenumi{\arabic{enumi}.}
\tightlist
\item
  Daten aus FreeCAD werden in Solver-Objekte übersetzt.
\item
  Jedes Teil bekommt eine Position und Orientierung.
\item
  Jede Verbindung wird als Constraint oder Joint modelliert.
\item
  Der Solver baut das Gleichungssystem auf.
\item
  Der Solver startet mit einer Näherung.
\item
  Fehler werden berechnet.
\item
  Iterationen verbessern die Lösung.
\item
  Das Ergebnis wird als neue Platzierungen zurückgegeben.
\end{enumerate}

\hypertarget{warum-dieses-modell-fuxfcr-freecad-relevant-ist}{%
\section{\texorpdfstring{{2.8} Warum dieses Modell für FreeCAD relevant
ist}{2.8 Warum dieses Modell für FreeCAD relevant ist}}\label{warum-dieses-modell-fuxfcr-freecad-relevant-ist}}

In FreeCAD macht das Modell genau folgendes:

\begin{itemize}
\tightlist
\item
  Die Benutzeroberfläche beschreibt nur die Beziehungen.
\item
  Der Solver führt die numerische Berechnung durch.
\item
  FreeCAD übernimmt das Ergebnis und zeigt es visuell an.
\end{itemize}

Damit bleibt die Bedienung verständlich, während der mathematische Kern
unabhängig arbeitet.

\hypertarget{ein-kleines-c-beispiel-konzept-statt-compile-ready-api}{%
\section{\texorpdfstring{{2.9} Ein kleines C++-Beispiel: Konzept statt
Compile-Ready-API}{2.9 Ein kleines C++-Beispiel: Konzept statt Compile-Ready-API}}\label{ein-kleines-c-beispiel-konzept-statt-compile-ready-api}}

Das folgende Beispiel ist bewusst einfach und dient dem Lernverständnis.
Es zeigt den denkbaren Ablauf, wie eine Beziehungsliste in einen
Solver-Kontext übersetzt werden kann.

\begin{Shaded}
\begin{Highlighting}[]
\PreprocessorTok{\#include }\ImportTok{\textless{}vector\textgreater{}}
\PreprocessorTok{\#include }\ImportTok{\textless{}string\textgreater{}}
\PreprocessorTok{\#include }\ImportTok{\textless{}cmath\textgreater{}}

\KeywordTok{struct}\NormalTok{ Part }\OperatorTok{\{}
    \BuiltInTok{std::}\NormalTok{string name}\OperatorTok{;}
    \DataTypeTok{double}\NormalTok{ x}\OperatorTok{;}
    \DataTypeTok{double}\NormalTok{ y}\OperatorTok{;}
    \DataTypeTok{double}\NormalTok{ z}\OperatorTok{;}
\OperatorTok{\};}

\KeywordTok{struct}\NormalTok{ DistanceConstraint }\OperatorTok{\{}
    \BuiltInTok{std::}\NormalTok{string partA}\OperatorTok{;}
    \BuiltInTok{std::}\NormalTok{string partB}\OperatorTok{;}
    \DataTypeTok{double}\NormalTok{ targetDistance}\OperatorTok{;}
\OperatorTok{\};}

\DataTypeTok{int}\NormalTok{ main}\OperatorTok{()} \OperatorTok{\{}
    \BuiltInTok{std::}\NormalTok{vector}\OperatorTok{\textless{}}\NormalTok{Part}\OperatorTok{\textgreater{}}\NormalTok{ parts }\OperatorTok{=} \OperatorTok{\{}
        \OperatorTok{\{}\StringTok{"A"}\OperatorTok{,} \FloatTok{0.0}\OperatorTok{,} \FloatTok{0.0}\OperatorTok{,} \FloatTok{0.0}\OperatorTok{\},}
        \OperatorTok{\{}\StringTok{"B"}\OperatorTok{,} \FloatTok{3.0}\OperatorTok{,} \FloatTok{0.0}\OperatorTok{,} \FloatTok{0.0}\OperatorTok{\}}
    \OperatorTok{\};}

\NormalTok{    DistanceConstraint c}\OperatorTok{\{}\StringTok{"A"}\OperatorTok{,} \StringTok{"B"}\OperatorTok{,} \FloatTok{5.0}\OperatorTok{\};}

    \CommentTok{// Einfache Fehlerfunktion: Abstand muss targetDistance sein.}
    \KeywordTok{auto}\NormalTok{ residual }\OperatorTok{=} \OperatorTok{[\&](}\AttributeTok{const}\NormalTok{ Part}\OperatorTok{\&}\NormalTok{ a}\OperatorTok{,} \AttributeTok{const}\NormalTok{ Part}\OperatorTok{\&}\NormalTok{ b}\OperatorTok{)} \OperatorTok{\{}
        \DataTypeTok{double}\NormalTok{ dx }\OperatorTok{=}\NormalTok{ b}\OperatorTok{.}\NormalTok{x }\OperatorTok{{-}}\NormalTok{ a}\OperatorTok{.}\NormalTok{x}\OperatorTok{;}
        \DataTypeTok{double}\NormalTok{ dy }\OperatorTok{=}\NormalTok{ b}\OperatorTok{.}\NormalTok{y }\OperatorTok{{-}}\NormalTok{ a}\OperatorTok{.}\NormalTok{y}\OperatorTok{;}
        \DataTypeTok{double}\NormalTok{ dz }\OperatorTok{=}\NormalTok{ b}\OperatorTok{.}\NormalTok{z }\OperatorTok{{-}}\NormalTok{ a}\OperatorTok{.}\NormalTok{z}\OperatorTok{;}
        \DataTypeTok{double}\NormalTok{ dist }\OperatorTok{=} \BuiltInTok{std::}\NormalTok{sqrt}\OperatorTok{(}\NormalTok{dx }\OperatorTok{*}\NormalTok{ dx }\OperatorTok{+}\NormalTok{ dy }\OperatorTok{*}\NormalTok{ dy }\OperatorTok{+}\NormalTok{ dz }\OperatorTok{*}\NormalTok{ dz}\OperatorTok{);}
        \ControlFlowTok{return}\NormalTok{ dist }\OperatorTok{{-}}\NormalTok{ c}\OperatorTok{.}\NormalTok{targetDistance}\OperatorTok{;}
    \OperatorTok{\};}

    \DataTypeTok{double}\NormalTok{ error }\OperatorTok{=}\NormalTok{ residual}\OperatorTok{(}\NormalTok{parts}\OperatorTok{[}\DecValTok{0}\OperatorTok{],}\NormalTok{ parts}\OperatorTok{[}\DecValTok{1}\OperatorTok{]);}
    \CommentTok{// Fehler wird iterativ reduziert.}
    \CommentTok{// Der Solver korrigiert dann die Positionen von B oder A.}

    \ControlFlowTok{return} \DecValTok{0}\OperatorTok{;}
\OperatorTok{\}}
\end{Highlighting}
\end{Shaded}

Wichtig ist hier nicht die genaue Bibliothek, sondern das Muster:

\begin{itemize}
\tightlist
\item
  Teile mit Variablen
\item
  Constraint mit gewünschter Kennzahl
\item
  Fehlerfunktion
\item
  iterative Annäherung
\end{itemize}

\hypertarget{ein-kleines-python-beispiel-datenmodellierung-und-fehlerbewertung}{%
\section{\texorpdfstring{{2.10} Ein kleines Python-Beispiel:
Datenmodellierung und
Fehlerbewertung}{2.10 Ein kleines Python-Beispiel: Datenmodellierung und Fehlerbewertung}}\label{ein-kleines-python-beispiel-datenmodellierung-und-fehlerbewertung}}

\begin{Shaded}
\begin{Highlighting}[]
\ImportTok{import}\NormalTok{ math}

\NormalTok{parts }\OperatorTok{=}\NormalTok{ \{}
    \StringTok{"A"}\NormalTok{: \{}\StringTok{"x"}\NormalTok{: }\FloatTok{0.0}\NormalTok{, }\StringTok{"y"}\NormalTok{: }\FloatTok{0.0}\NormalTok{, }\StringTok{"z"}\NormalTok{: }\FloatTok{0.0}\NormalTok{\},}
    \StringTok{"B"}\NormalTok{: \{}\StringTok{"x"}\NormalTok{: }\FloatTok{3.0}\NormalTok{, }\StringTok{"y"}\NormalTok{: }\FloatTok{0.0}\NormalTok{, }\StringTok{"z"}\NormalTok{: }\FloatTok{0.0}\NormalTok{\},}
\NormalTok{\}}

\KeywordTok{def}\NormalTok{ distance\_error(part\_a, part\_b, target\_distance):}
\NormalTok{    dx }\OperatorTok{=}\NormalTok{ part\_b[}\StringTok{"x"}\NormalTok{] }\OperatorTok{{-}}\NormalTok{ part\_a[}\StringTok{"x"}\NormalTok{]}
\NormalTok{    dy }\OperatorTok{=}\NormalTok{ part\_b[}\StringTok{"y"}\NormalTok{] }\OperatorTok{{-}}\NormalTok{ part\_a[}\StringTok{"y"}\NormalTok{]}
\NormalTok{    dz }\OperatorTok{=}\NormalTok{ part\_b[}\StringTok{"z"}\NormalTok{] }\OperatorTok{{-}}\NormalTok{ part\_a[}\StringTok{"z"}\NormalTok{]}
\NormalTok{    actual }\OperatorTok{=}\NormalTok{ math.sqrt(dx }\OperatorTok{*}\NormalTok{ dx }\OperatorTok{+}\NormalTok{ dy }\OperatorTok{*}\NormalTok{ dy }\OperatorTok{+}\NormalTok{ dz }\OperatorTok{*}\NormalTok{ dz)}
    \ControlFlowTok{return}\NormalTok{ actual }\OperatorTok{{-}}\NormalTok{ target\_distance}

\NormalTok{error }\OperatorTok{=}\NormalTok{ distance\_error(parts[}\StringTok{"A"}\NormalTok{], parts[}\StringTok{"B"}\NormalTok{], }\FloatTok{5.0}\NormalTok{)}
\BuiltInTok{print}\NormalTok{(}\SpecialStringTok{f"Fehler vor Korrektur: }\SpecialCharTok{\{}\NormalTok{error}\SpecialCharTok{\}}\SpecialStringTok{"}\NormalTok{)}
\end{Highlighting}
\end{Shaded}

Das zeigt auch die zentrale Idee: Eine Beziehung ist in der Regel eine
Funktion, die einen Fehlerwert liefert. Je kleiner der Fehler, desto
näher ist der Zustand an einer gültigen Lösung.

\hypertarget{fazit}{%
\section{\texorpdfstring{{2.11} Fazit}{2.11 Fazit}}\label{fazit}}

OndselSolver funktioniert nicht durch „magische Intuition``, sondern
nach einem klaren logischen Muster:

\begin{itemize}
\tightlist
\item
  Beziehungen werden in mathematische Bedingungen übersetzt
\item
  Fehler werden bewertet
\item
  Positionen werden iterativ angepasst
\item
  Eine Lösung wird gefunden, wenn der Fehler klein genug ist
\end{itemize}

Das ist das zentrale Verständnis, das für die gesamte
FreeCAD-Assembly-Arbeit wichtig ist.

\hypertarget{kap-datenvorbereitung}{%
\chapter{\texorpdfstring{{3} Datenvorbereitung für den Solver und
Ergebnisinterpretation}{3 Datenvorbereitung für den Solver und Ergebnisinterpretation}}\label{kap-datenvorbereitung}}

\hypertarget{warum-datenvorbereitung-so-wichtig-ist}{%
\section{\texorpdfstring{{3.1} Warum Datenvorbereitung so wichtig
ist}{3.1 Warum Datenvorbereitung so wichtig ist}}\label{warum-datenvorbereitung-so-wichtig-ist}}

Ein Solver ist nur so gut wie die Daten, die ihm gegeben werden. Wenn
die Eingabedaten schlecht modelliert, inkonsistent, doppeldeutig oder
numerisch instabil sind, wird selbst ein guter Solver keine zuverlässige
Antwort liefern.

Für FreeCAD bedeutet das: Eine Baugruppe muss sinnvoll aufgebaut sein,
damit OndselSolver sie verarbeiten kann.

\hypertarget{typische-vorbereitungsaufgaben}{%
\section{\texorpdfstring{{3.2} Typische
Vorbereitungsaufgaben}{3.2 Typische Vorbereitungsaufgaben}}\label{typische-vorbereitungsaufgaben}}

Vor einem Solve sind häufig diese Schritte nötig:

\begin{itemize}
\tightlist
\item
  Teil- und Joint-Definitionen prüfen
\item
  Referenzen auf richtige Objekte auflösen
\item
  unerwünschte oder redundante Constraints entfernen
\item
  geerdete Teile und relative Bewegungen sauber trennen
\item
  Startwerte plausibel wählen
\item
  numerische Toleranzen prüfen
\end{itemize}

\hypertarget{gute-datenstruktur-fuxfcr-den-solver}{%
\section{\texorpdfstring{{3.3} Gute Datenstruktur für den
Solver}{3.3 Gute Datenstruktur für den Solver}}\label{gute-datenstruktur-fuxfcr-den-solver}}

Ein einfacher und verständlicher Stil ist dieser:

\begin{itemize}
\tightlist
\item
  Jedes Teil erhält eine eindeutige ID
\item
  Jedes Joint oder Constraint erhält einen Typ
\item
  Jedes Joint enthält seine Referenzen und gewünschte Parameter
\item
  Der Solver bekommt eine konsistente Liste dieser Objekte
\end{itemize}

In Python sieht das typischerweise so aus:

\begin{Shaded}
\begin{Highlighting}[]
\NormalTok{parts }\OperatorTok{=}\NormalTok{ \{}
    \StringTok{"base"}\NormalTok{: \{}\StringTok{"id"}\NormalTok{: }\StringTok{"base"}\NormalTok{, }\StringTok{"x"}\NormalTok{: }\FloatTok{0.0}\NormalTok{, }\StringTok{"y"}\NormalTok{: }\FloatTok{0.0}\NormalTok{, }\StringTok{"z"}\NormalTok{: }\FloatTok{0.0}\NormalTok{\},}
    \StringTok{"arm"}\NormalTok{: \{}\StringTok{"id"}\NormalTok{: }\StringTok{"arm"}\NormalTok{, }\StringTok{"x"}\NormalTok{: }\FloatTok{2.0}\NormalTok{, }\StringTok{"y"}\NormalTok{: }\FloatTok{0.0}\NormalTok{, }\StringTok{"z"}\NormalTok{: }\FloatTok{0.0}\NormalTok{\},}
\NormalTok{\}}

\NormalTok{joints }\OperatorTok{=}\NormalTok{ [}
\NormalTok{    \{}\StringTok{"type"}\NormalTok{: }\StringTok{"revolute"}\NormalTok{, }\StringTok{"part\_a"}\NormalTok{: }\StringTok{"base"}\NormalTok{, }\StringTok{"part\_b"}\NormalTok{: }\StringTok{"arm"}\NormalTok{, }\StringTok{"axis"}\NormalTok{: }\StringTok{"z"}\NormalTok{\},}
\NormalTok{]}
\end{Highlighting}
\end{Shaded}

Diese Struktur ist nicht „die echte API`` von OndselSolver, sondern ein
Lernmodell. Es zeigt, wie man Daten logisch gruppiert und dem Solver
bereitstellt.

\hypertarget{c-ansatz-fuxfcr-solver-daten}{%
\section{\texorpdfstring{{3.4} C++-Ansatz für
Solver-Daten}{3.4 C++-Ansatz für Solver-Daten}}\label{c-ansatz-fuxfcr-solver-daten}}

In C++ ist eine klare Struktur oft sehr hilfreich. Ein Beispiel:

\begin{Shaded}
\begin{Highlighting}[]
\PreprocessorTok{\#include }\ImportTok{\textless{}string\textgreater{}}
\PreprocessorTok{\#include }\ImportTok{\textless{}vector\textgreater{}}

\KeywordTok{struct}\NormalTok{ PartData }\OperatorTok{\{}
    \BuiltInTok{std::}\NormalTok{string id}\OperatorTok{;}
    \DataTypeTok{double}\NormalTok{ x}\OperatorTok{;}
    \DataTypeTok{double}\NormalTok{ y}\OperatorTok{;}
    \DataTypeTok{double}\NormalTok{ z}\OperatorTok{;}
\OperatorTok{\};}

\KeywordTok{struct}\NormalTok{ JointData }\OperatorTok{\{}
    \BuiltInTok{std::}\NormalTok{string type}\OperatorTok{;}
    \BuiltInTok{std::}\NormalTok{string partA}\OperatorTok{;}
    \BuiltInTok{std::}\NormalTok{string partB}\OperatorTok{;}
    \DataTypeTok{double}\NormalTok{ value}\OperatorTok{;}
\OperatorTok{\};}

\DataTypeTok{int}\NormalTok{ main}\OperatorTok{()} \OperatorTok{\{}
    \BuiltInTok{std::}\NormalTok{vector}\OperatorTok{\textless{}}\NormalTok{PartData}\OperatorTok{\textgreater{}}\NormalTok{ parts }\OperatorTok{=} \OperatorTok{\{}
        \OperatorTok{\{}\StringTok{"base"}\OperatorTok{,} \FloatTok{0.0}\OperatorTok{,} \FloatTok{0.0}\OperatorTok{,} \FloatTok{0.0}\OperatorTok{\},}
        \OperatorTok{\{}\StringTok{"arm"}\OperatorTok{,} \FloatTok{2.0}\OperatorTok{,} \FloatTok{0.0}\OperatorTok{,} \FloatTok{0.0}\OperatorTok{\}}
    \OperatorTok{\};}

    \BuiltInTok{std::}\NormalTok{vector}\OperatorTok{\textless{}}\NormalTok{JointData}\OperatorTok{\textgreater{}}\NormalTok{ joints }\OperatorTok{=} \OperatorTok{\{}
        \OperatorTok{\{}\StringTok{"revolute"}\OperatorTok{,} \StringTok{"base"}\OperatorTok{,} \StringTok{"arm"}\OperatorTok{,} \FloatTok{90.0}\OperatorTok{\}}
    \OperatorTok{\};}

    \CommentTok{// Die Daten sind jetzt bereit für die Systembildung.}
    \CommentTok{// Der Solver kann diese Informationen in Constraints transformieren.}
    \ControlFlowTok{return} \DecValTok{0}\OperatorTok{;}
\OperatorTok{\}}
\end{Highlighting}
\end{Shaded}

Der Kern ist hier nicht die genaue API, sondern das Prinzip: Die Daten
müssen sauber, strukturiert und konsistent sein.

\hypertarget{huxe4ufige-fehler-beim-vorbereiten-von-solver-daten}{%
\section{\texorpdfstring{{3.5} Häufige Fehler beim Vorbereiten von
Solver-Daten}{3.5 Häufige Fehler beim Vorbereiten von Solver-Daten}}\label{huxe4ufige-fehler-beim-vorbereiten-von-solver-daten}}

Diese Punkte treten besonders oft auf:

\begin{itemize}
\tightlist
\item
  Referenz auf falsches Teil
\item
  unvollständige Gelenk-Definition
\item
  inkonsistente Achsenausrichtung
\item
  Redundanz von Constraints
\item
  Geometrie mit zu wenig Informationen
\item
  schlechte Startwerte
\end{itemize}

Diese Faktoren führen dazu, dass der Solver entweder nicht konvergiert
oder eine Lösung findet, die zwar mathematisch denkbar ist, aber für die
reale Anwendung unbrauchbar bleibt.

\hypertarget{ergebnisinterpretation}{%
\section{\texorpdfstring{{3.6}
Ergebnisinterpretation}{3.6 Ergebnisinterpretation}}\label{ergebnisinterpretation}}

Wenn ein Solver eine Berechnung beendet, muss der Leser nicht nur
wissen, ob der Solve „funktioniert hat``, sondern auch was das Ergebnis
bedeutet.

Wichtige Fragen:

\begin{itemize}
\tightlist
\item
  Ist der Lösungsfehler klein genug?
\item
  Gibt es noch freie Freiheitsgrade?
\item
  Waren die Constraints wirklich erfüllt?
\item
  Hat der Solver eine gültige Lösung gefunden oder nur einen
  Konsistenzpunkt?
\end{itemize}

\hypertarget{wie-fehlerwerte-interpretiert-werden}{%
\section{\texorpdfstring{{3.7} Wie Fehlerwerte interpretiert
werden}{3.7 Wie Fehlerwerte interpretiert werden}}\label{wie-fehlerwerte-interpretiert-werden}}

Ein Solver liefert normalerweise keine „schöne`` Zahl, sondern eine
Residuen- oder Fehlerinformation. Das muss man lesen:

\begin{itemize}
\tightlist
\item
  Residuum nahe null: gute Übereinstimmung
\item
  großes Residuum: Constraint nicht erfüllt
\item
  keine Konvergenz: Ausgangslage oder Modell unpassend
\item
  multiple Lösungen: gleiche Beziehung kann mehrere geometrische
  Lösungen zulassen
\end{itemize}

\hypertarget{python-example-zur-ergebnispruxfcfung}{%
\section{\texorpdfstring{{3.8} Python-Example zur
Ergebnisprüfung}{3.8 Python-Example zur Ergebnisprüfung}}\label{python-example-zur-ergebnispruxfcfung}}

\begin{Shaded}
\begin{Highlighting}[]
\NormalTok{results }\OperatorTok{=}\NormalTok{ \{}
    \StringTok{"base"}\NormalTok{: \{}\StringTok{"x"}\NormalTok{: }\FloatTok{0.0}\NormalTok{, }\StringTok{"y"}\NormalTok{: }\FloatTok{0.0}\NormalTok{, }\StringTok{"z"}\NormalTok{: }\FloatTok{0.0}\NormalTok{\},}
    \StringTok{"arm"}\NormalTok{: \{}\StringTok{"x"}\NormalTok{: }\FloatTok{4.0}\NormalTok{, }\StringTok{"y"}\NormalTok{: }\FloatTok{0.0}\NormalTok{, }\StringTok{"z"}\NormalTok{: }\FloatTok{0.0}\NormalTok{\},}
\NormalTok{\}}

\NormalTok{required\_distance }\OperatorTok{=} \FloatTok{4.0}

\KeywordTok{def}\NormalTok{ check\_distance(part\_a, part\_b, expected):}
\NormalTok{    dx }\OperatorTok{=}\NormalTok{ part\_b[}\StringTok{"x"}\NormalTok{] }\OperatorTok{{-}}\NormalTok{ part\_a[}\StringTok{"x"}\NormalTok{]}
\NormalTok{    dy }\OperatorTok{=}\NormalTok{ part\_b[}\StringTok{"y"}\NormalTok{] }\OperatorTok{{-}}\NormalTok{ part\_a[}\StringTok{"y"}\NormalTok{]}
\NormalTok{    dz }\OperatorTok{=}\NormalTok{ part\_b[}\StringTok{"z"}\NormalTok{] }\OperatorTok{{-}}\NormalTok{ part\_a[}\StringTok{"z"}\NormalTok{]}
\NormalTok{    actual }\OperatorTok{=}\NormalTok{ (dx }\OperatorTok{*}\NormalTok{ dx }\OperatorTok{+}\NormalTok{ dy }\OperatorTok{*}\NormalTok{ dy }\OperatorTok{+}\NormalTok{ dz }\OperatorTok{*}\NormalTok{ dz) }\OperatorTok{**} \FloatTok{0.5}
\NormalTok{    residual }\OperatorTok{=} \BuiltInTok{abs}\NormalTok{(actual }\OperatorTok{{-}}\NormalTok{ expected)}
    \ControlFlowTok{return}\NormalTok{ residual}

\NormalTok{residual }\OperatorTok{=}\NormalTok{ check\_distance(results[}\StringTok{"base"}\NormalTok{], results[}\StringTok{"arm"}\NormalTok{], required\_distance)}
\BuiltInTok{print}\NormalTok{(}\SpecialStringTok{f"Residual der Abstandsbedingung: }\SpecialCharTok{\{}\NormalTok{residual}\SpecialCharTok{\}}\SpecialStringTok{"}\NormalTok{)}
\end{Highlighting}
\end{Shaded}

Das zeigt die zentrale Regel: Nicht „es gibt ein Ergebnis``, sondern
„das Ergebnis erfüllt die gewünschte Beziehung innerhalb der Toleranz``
ist entscheidend.

\hypertarget{die-praktische-bedeutung-in-freecad}{%
\section{\texorpdfstring{{3.9} Die praktische Bedeutung in
FreeCAD}{3.9 Die praktische Bedeutung in FreeCAD}}\label{die-praktische-bedeutung-in-freecad}}

In FreeCAD bedeutet Ergebnisinterpretation konkret:

\begin{itemize}
\tightlist
\item
  untersucht man die Position der gelösten Teile,
\item
  prüft man anschließend, ob die gezeichneten Beziehungen konsistent
  sind,
\item
  vergleicht man die von der UI gezeigten Werte mit den vom Solver
  berechneten Werten,
\item
  erkennt man früh, ob das Modell numerisch stabil ist.
\end{itemize}

\hypertarget{fazit-1}{%
\section{\texorpdfstring{{3.10} Fazit}{3.10 Fazit}}\label{fazit-1}}

Vorbereitung und Interpretation sind zwei Seiten derselben Aufgabe:

\begin{itemize}
\tightlist
\item
  Die Vorbereitungsphase macht aus einer Idee ein mathematisch gültiges
  System.
\item
  Die Interpretationsphase macht aus einer Lösung eine verständliche,
  nutzbare Baugruppenlage.
\end{itemize}

Ohne diese beiden Schritte wird ein Solver nicht als hilfreiches
Werkzeug wahrgenommen, sondern als „Black Box``, die manchmal
funktioniert und manchmal nicht.

\hypertarget{kap-kommunikation-mit-freecad}{%
\chapter{\texorpdfstring{{4} Wie kommuniziert OndselSolver mit
FreeCAD?}{4 Wie kommuniziert OndselSolver mit FreeCAD?}}\label{kap-kommunikation-mit-freecad}}

\hypertarget{das-grundmodell}{%
\section{\texorpdfstring{{4.1} Das
Grundmodell}{4.1 Das Grundmodell}}\label{das-grundmodell}}

Die Verbindung zwischen FreeCAD und OndselSolver ist kein „magischer
Direktzugriff``, sondern ein klarer Datenfluss:

\begin{enumerate}
\def\labelenumi{\arabic{enumi}.}
\tightlist
\item
  FreeCAD modelliert Objekte, Plazierungen und Beziehungen.
\item
  Die Assembly-Architektur sammelt diese Informationen.
\item
  Die Daten werden in eine Solver-geeignete Struktur übersetzt.
\item
  OndselSolver berechnet die Lösung.
\item
  FreeCAD übernimmt die berechneten Platzierungen erneut in die
  3D-Objekte.
\end{enumerate}

Das ist die zentrale Kommunikationsschicht.

\hypertarget{ein-wichtiger-gedanke-darstellung-und-solver-sind-getrennt}{%
\section{\texorpdfstring{{4.2} Ein wichtiger Gedanke: Darstellung und
Solver sind
getrennt}{4.2 Ein wichtiger Gedanke: Darstellung und Solver sind getrennt}}\label{ein-wichtiger-gedanke-darstellung-und-solver-sind-getrennt}}

Eine FreeCAD-Baugruppe hat eine sichtbare 3D-Darstellung und eine
Mechanik-/Solver-Logik. Diese beiden Ebenen sind nicht identisch.

\begin{itemize}
\tightlist
\item
  Die GUI zeigt, was der Nutzer sieht.
\item
  Die Baugruppenstruktur hält die Beziehungen fest.
\item
  Der Solver arbeitet mit den modellierten Bedingungen.
\item
  Die Resultate werden in die Objekte zurückgeschrieben.
\end{itemize}

Diese Trennung ist wichtig. Der Datenfluss ist nicht „ein Objekt ändert
sich direkt im Solver und alles ist automatisch perfekt``, sondern
„FreeCAD organisiert die Daten, der Solver verarbeitet sie, dann
schreibt FreeCAD die Ergebnisse zurück``.

\hypertarget{ein-typischer-ablauf-in-der-freecad-assembly-architektur}{%
\section{\texorpdfstring{{4.3} Ein typischer Ablauf in der
FreeCAD-Assembly-Architektur}{4.3 Ein typischer Ablauf in der FreeCAD-Assembly-Architektur}}\label{ein-typischer-ablauf-in-der-freecad-assembly-architektur}}

Ein realistischer Ablauf sieht so aus:

\begin{itemize}
\tightlist
\item
  Ein Objekt wird als Teil in einer Baugruppe erkannt.
\item
  Ein Joint oder Constraint wird als Beziehung hinzugefügt.
\item
  Die Baugruppenlogik sammelt alle relevanten Objekte und Beziehungen.
\item
  Ein \texttt{AssemblyObject} oder ein vergleichbares Objekt dient als
  Fassade für den Solver.
\item
  Der Solver erzeugt intern ein Modell aus Teilen, Körpern und
  Beziehungen.
\item
  Nach der Berechnung werden die neuen Platzierungen in die
  FreeCAD-Objekte übernommen.
\end{itemize}

\hypertarget{warum-die-kommunikationsschicht-entscheidend-ist}{%
\section{\texorpdfstring{{4.4} Warum die Kommunikationsschicht
entscheidend
ist}{4.4 Warum die Kommunikationsschicht entscheidend ist}}\label{warum-die-kommunikationsschicht-entscheidend-ist}}

Die größte Schwierigkeit ist oft nicht der mathematische Kern, sondern
die Übersetzung zwischen zwei Welten:

\begin{itemize}
\tightlist
\item
  FreeCAD spricht in Objekten, Referenzen, Links, Plazierungen und
  Joints.
\item
  OndselSolver spricht in Körpern, Variablen, Constraints und iterativen
  Lösungsschritten.
\end{itemize}

Damit liegt die eigentliche Herausforderung darin, genau diese Daten
korrekt zu übersetzen.

\hypertarget{die-rolle-von-referenzen}{%
\section{\texorpdfstring{{4.5} Die Rolle von
Referenzen}{4.5 Die Rolle von Referenzen}}\label{die-rolle-von-referenzen}}

In FreeCAD sind Beziehungen häufig als Referenzen gespeichert. Ein Joint
zeigt zum Beispiel auf zwei Objekte oder Subkomponenten und definiert
ihre relative Lage.

Der Solver braucht daraus eine kompakte, berechenbare Darstellung. Das
heißt:

\begin{itemize}
\tightlist
\item
  Objekt identifizieren
\item
  Beziehung bestimmen
\item
  Parameter auslesen
\item
  numerischen Zustand aufbauen
\end{itemize}

Das ist der Kern einer Projektion von CAD-Daten in Solver-Daten.

\hypertarget{das-problem-der-identituxe4t}{%
\section{\texorpdfstring{{4.6} Das Problem der
Identität}{4.6 Das Problem der Identität}}\label{das-problem-der-identituxe4t}}

Ein zentrales Thema in der FreeCAD-Assembly-Architektur ist Identität:

\begin{itemize}
\tightlist
\item
  Das sichtbare Objekt kann eine Spiegelung oder ein Link sein.
\item
  Das echte Objekt kann an einer anderen Stelle im Dokument liegen.
\item
  Der Solver will eine eindeutige, konsistente Teil-Identität.
\end{itemize}

Das erklärt, warum in komplexen Baugruppen sehr viel an Auflösung und
Abbildung geschieht. Die Kommunikationsschicht muss sicherstellen, dass
identische physische Teile nicht verwechselt werden.

\hypertarget{warum-assemblyobject-und-assemblylink-wichtig-sind}{%
\section{\texorpdfstring{{4.7} Warum \texttt{AssemblyObject} und
\texttt{AssemblyLink} wichtig
sind}{4.7 Warum AssemblyObject und AssemblyLink wichtig sind}}\label{warum-assemblyobject-und-assemblylink-wichtig-sind}}

Im FreeCAD-Assembly-Code sind diese Strukturen der zentrale
Übersetzungsmechanismus:

\begin{itemize}
\tightlist
\item
  \texttt{AssemblyObject} ist die eigentliche Baugruppe und die
  Solver-Fassade
\item
  \texttt{AssemblyLink} stellt eingebettete oder verlinkte Strukturen in
  der Baugruppe dar
\end{itemize}

Diese Objekte helfen dabei, die freie CAD-Welt und die Solver-Welt
sauber zu trennen und wieder zusammenzubinden.

\hypertarget{was-bei-der-ruxfcckgabe-an-freecad-passiert}{%
\section{\texorpdfstring{{4.8} Was bei der Rückgabe an FreeCAD
passiert}{4.8 Was bei der Rückgabe an FreeCAD passiert}}\label{was-bei-der-ruxfcckgabe-an-freecad-passiert}}

Nach der Berechnung ist ein wichtiger Schritt die Rückgabe der Lösung in
die CAD-Welt:

\begin{itemize}
\tightlist
\item
  neue Plazierungen werden gesetzt
\item
  relative Beziehungen werden aktualisiert
\item
  visuelle Darstellung wird neu gezeichnet
\item
  Statusinformationen werden ergänzt
\end{itemize}

Ohne diese Rückgabe wäre die Lösung zwar numerisch korrekt, aber für den
Benutzer nicht sichtbar.

\hypertarget{beispiel-datenfluss-in-verstuxe4ndlicher-form}{%
\section{\texorpdfstring{{4.9} Beispiel: Datenfluss in verständlicher
Form}{4.9 Beispiel: Datenfluss in verständlicher Form}}\label{beispiel-datenfluss-in-verstuxe4ndlicher-form}}

\begin{Shaded}
\begin{Highlighting}[]
\CommentTok{\# 1) FreeCAD{-}Klasse modelliert Benutzer{-}Input}
\NormalTok{assembly }\OperatorTok{=}\NormalTok{ \{}
    \StringTok{"parts"}\NormalTok{: [}\StringTok{"base"}\NormalTok{, }\StringTok{"arm"}\NormalTok{],}
    \StringTok{"joints"}\NormalTok{: [\{}\StringTok{"type"}\NormalTok{: }\StringTok{"revolute"}\NormalTok{, }\StringTok{"a"}\NormalTok{: }\StringTok{"base"}\NormalTok{, }\StringTok{"b"}\NormalTok{: }\StringTok{"arm"}\NormalTok{\}],}
\NormalTok{\}}

\CommentTok{\# 2) Übersetzung in Solver{-}Daten}
\NormalTok{solver\_model }\OperatorTok{=}\NormalTok{ \{}
    \StringTok{"bodies"}\NormalTok{: [}
\NormalTok{        \{}\StringTok{"id"}\NormalTok{: }\StringTok{"base"}\NormalTok{, }\StringTok{"position"}\NormalTok{: [}\DecValTok{0}\NormalTok{, }\DecValTok{0}\NormalTok{, }\DecValTok{0}\NormalTok{]\},}
\NormalTok{        \{}\StringTok{"id"}\NormalTok{: }\StringTok{"arm"}\NormalTok{, }\StringTok{"position"}\NormalTok{: [}\DecValTok{2}\NormalTok{, }\DecValTok{0}\NormalTok{, }\DecValTok{0}\NormalTok{]\},}
\NormalTok{    ],}
    \StringTok{"constraints"}\NormalTok{: [\{}\StringTok{"kind"}\NormalTok{: }\StringTok{"revolute"}\NormalTok{, }\StringTok{"body\_a"}\NormalTok{: }\StringTok{"base"}\NormalTok{, }\StringTok{"body\_b"}\NormalTok{: }\StringTok{"arm"}\NormalTok{\}],}
\NormalTok{\}}

\CommentTok{\# 3) Ergebnis aus dem Solver}
\NormalTok{result }\OperatorTok{=}\NormalTok{ \{}
    \StringTok{"arm"}\NormalTok{: \{}\StringTok{"position"}\NormalTok{: [}\DecValTok{2}\NormalTok{, }\DecValTok{0}\NormalTok{, }\DecValTok{0}\NormalTok{], }\StringTok{"rotation"}\NormalTok{: [}\DecValTok{0}\NormalTok{, }\DecValTok{0}\NormalTok{, }\FloatTok{0.5}\NormalTok{]\},}
\NormalTok{\}}

\CommentTok{\# 4) FreeCAD übernimmt die Platzierung}
\CommentTok{\# {-}\textgreater{} Darstellung wird erneuert}
\end{Highlighting}
\end{Shaded}

Das ist ein didaktisches Beispiel. In der echten Architektur erfolgt die
Übersetzung deutlich komplexer, aber das Muster bleibt dasselbe.

\hypertarget{wichtige-lernfrage-fuxfcr-den-leser}{%
\section{\texorpdfstring{{4.10} Wichtige Lernfrage für den
Leser}{4.10 Wichtige Lernfrage für den Leser}}\label{wichtige-lernfrage-fuxfcr-den-leser}}

Wenn ein Solver in FreeCAD nicht funktioniert, dann lohnt es sich immer,
diese Fragen zu beantworten:

\begin{itemize}
\tightlist
\item
  Welche Objekte wurden genau an den Solver übergeben?
\item
  Welche Beziehungen wurden modelliert?
\item
  Gab es eine fehlerhafte Referenz?
\item
  Wurden die Resultate richtig zurück in die Platzierungen übernommen?
\end{itemize}

Das sind die Fragen, die die Kommunikationsschicht zwischen FreeCAD und
OndselSolver verständlich machen.

\hypertarget{fazit-2}{%
\section{\texorpdfstring{{4.11} Fazit}{4.11 Fazit}}\label{fazit-2}}

OndselSolver und FreeCAD kommunizieren nicht über „eine gemeinsame
Realität``, sondern über eine gut definierte Daten- und
Übersetzungslogik.

Die harte Arbeit liegt nicht nur im Rechnen, sondern in der sauberen
Verknüpfung von:

\begin{itemize}
\tightlist
\item
  3D-Modellierung
\item
  semantischer Baugruppenlogik
\item
  geometrischen Beziehungen
\item
  Solver-Parametern
\item
  resultierenden Placements
\end{itemize}

Das ist der Kern, den ein Lernender verstehen muss, wenn er FreeCAD und
OndselSolver wirklich durchdringen will.

\hypertarget{kap-freecad-beispiele}{%
\chapter{\texorpdfstring{{5} FreeCAD-Beispiele mit Bibliotheken und
Solver-Integration}{5 FreeCAD-Beispiele mit Bibliotheken und Solver-Integration}}\label{kap-freecad-beispiele}}

\hypertarget{warum-dieses-kapitel-wichtig-ist}{%
\section{\texorpdfstring{{5.1} Warum dieses Kapitel wichtig
ist}{5.1 Warum dieses Kapitel wichtig ist}}\label{warum-dieses-kapitel-wichtig-ist}}

Ein Leser versteht OndselSolver erst wirklich, wenn er sieht, wie die
Bibliotheken in der echten FreeCAD-Welt zusammenspielen. In der Praxis
ist das nicht „ein Solver, der isoliert läuft``, sondern ein
Zusammenspiel aus mehreren Bibliotheken und Modulen.

In FreeCADs Assembly-Implementierung ist das sichtbar:

\begin{itemize}
\tightlist
\item
  \href{https://github.com/FreeCAD/FreeCAD/blob/main/src/Mod/Assembly/App/CMakeLists.txt}{Assembly/App/CMakeLists.txt}
  verbindet das Assembly-Modul mit \texttt{Part}, \texttt{PartDesign},
  \texttt{Spreadsheet}, \texttt{FreeCADApp} und \texttt{OndselSolver}
\item
  \href{https://github.com/FreeCAD/FreeCAD/blob/main/src/Mod/Assembly/CommandExportASMT.py}{CommandExportASMT.py}
  exportiert eine Baugruppe als \texttt{.asmt}
\item
  \href{https://github.com/FreeCAD/FreeCAD/blob/main/src/Mod/Assembly/App/AssemblyObject.cpp}{AssemblyObject.cpp}
  und
  \href{https://github.com/FreeCAD/FreeCAD/blob/main/src/Mod/Assembly/App/AssemblyObject.h}{AssemblyObject.h}
  kapseln die Baugruppenlogik und den Solver-Zugang
\end{itemize}

Das zeigt: Es gibt eine klare Trennung zwischen

\begin{itemize}
\tightlist
\item
  CAD-Bibliotheken und Objekten,
\item
  Assembly-Logik,
\item
  Solver-Integration,
\item
  GUI/Interaktion.
\end{itemize}

\hypertarget{die-wichtigsten-bibliotheken-im-freecad-assembly-kontext}{%
\section{\texorpdfstring{{5.2} Die wichtigsten Bibliotheken im
FreeCAD-Assembly-Kontext}{5.2 Die wichtigsten Bibliotheken im FreeCAD-Assembly-Kontext}}\label{die-wichtigsten-bibliotheken-im-freecad-assembly-kontext}}

\hypertarget{freecadapp}{%
\subsection{\texorpdfstring{{5.2.1}
FreeCADApp}{5.2.1 FreeCADApp}}\label{freecadapp}}

\texttt{FreeCADApp} ist die Kernbibliothek zur Verwaltung von
Dokumenten, Objekten, Eigenschaften und Parametern.

Sie beschreibt die Grundlage, auf der ein Part, ein Link, eine Baugruppe
oder ein Joint entsteht.

\hypertarget{part-partdesign}{%
\subsection{\texorpdfstring{{5.2.2} Part /
PartDesign}{5.2.2 Part / PartDesign}}\label{part-partdesign}}

Diese Bibliotheken liefern die geometrischen und parametrischen
Funktionen. Sie sind wichtig, wenn geometrische Formen, Plazierungen und
profilebasiertes Arbeiten in die Solver-Logik eingebunden werden.

\hypertarget{assembly-modul}{%
\subsection{\texorpdfstring{{5.2.3}
Assembly-Modul}{5.2.3 Assembly-Modul}}\label{assembly-modul}}

Das Assembly-Modul ist die konkrete Brücke zwischen:

\begin{itemize}
\tightlist
\item
  geometrischem CAD-Modell,
\item
  Joints/Constraints,
\item
  Solver-Objekten,
\item
  berechneten Resultaten.
\end{itemize}

Im realen Code ist das in
\href{https://github.com/FreeCAD/FreeCAD/blob/main/src/Mod/Assembly/App/CMakeLists.txt}{Assembly/App/CMakeLists.txt}
gut sichtbar:

\begin{Shaded}
\begin{Highlighting}[]
\KeywordTok{set}\NormalTok{(Assembly\_LIBS}
\NormalTok{    Part}
\NormalTok{    PartDesign}
\NormalTok{    Spreadsheet}
\NormalTok{    FreeCADApp}
\NormalTok{    OndselSolver}
\NormalTok{)}
\end{Highlighting}
\end{Shaded}

Das ist ein sehr wichtiges Beispiel für die Integrationslogik: Das
Assembly-Modul verwendet FreeCAD-Bibliotheken und zusätzlich die
Solver-Bibliothek.

\hypertarget{typischer-datenfluss-in-einem-freecad-setup}{%
\section{\texorpdfstring{{5.3} Typischer Datenfluss in einem
FreeCAD-Setup}{5.3 Typischer Datenfluss in einem FreeCAD-Setup}}\label{typischer-datenfluss-in-einem-freecad-setup}}

Ein realer Ablauf sieht in der Praxis typischerweise so aus:

\begin{enumerate}
\def\labelenumi{\arabic{enumi}.}
\tightlist
\item
  FreeCAD erzeugt ein Dokument und Objekte.
\item
  Assembly-Objekte sammeln die Teile und Joints.
\item
  Das Assembly-Modul ordnet sie für den Solver an.
\item
  OndselSolver löst die Beziehungssysteme.
\item
  Die Ergebnisse werden in die FreeCAD-Platzierungen zurück geschrieben.
\end{enumerate}

Das heißt: Der Solver ist kein selbstständiges CAD-System, sondern ein
Teil einer größeren Architektur.

\hypertarget{ein-kleines-python-beispiel-mit-freecad-und-solver-logik}{%
\section{\texorpdfstring{{5.4} Ein kleines Python-Beispiel mit FreeCAD
und
Solver-Logik}{5.4 Ein kleines Python-Beispiel mit FreeCAD und Solver-Logik}}\label{ein-kleines-python-beispiel-mit-freecad-und-solver-logik}}

Das folgende Beispiel ist bewusst einfach und zeigt das Grundmuster,
ohne die komplette FreeCAD-API zu reproduzieren.

\begin{Shaded}
\begin{Highlighting}[]
\ImportTok{import}\NormalTok{ FreeCAD }\ImportTok{as}\NormalTok{ App}

\CommentTok{\# Ein einfaches Dokument als Basis für eine Baugruppe}
\CommentTok{\# In der echten FreeCAD{-}Architektur würden hier Beispiele mit AssemblyObject,}
\CommentTok{\# Joint{-}Definitionen und Solver{-}Aufrufen folgen.}
\NormalTok{doc }\OperatorTok{=}\NormalTok{ App.newDocument(}\StringTok{"DemoAssembly"}\NormalTok{)}

\NormalTok{base }\OperatorTok{=}\NormalTok{ doc.addObject(}\StringTok{"Part::Box"}\NormalTok{, }\StringTok{"Base"}\NormalTok{)}
\NormalTok{arm }\OperatorTok{=}\NormalTok{ doc.addObject(}\StringTok{"Part::Box"}\NormalTok{, }\StringTok{"Arm"}\NormalTok{)}

\CommentTok{\# Repräsentation der geometrischen Daten}
\NormalTok{base.Placement.Base }\OperatorTok{=}\NormalTok{ (}\DecValTok{0}\NormalTok{, }\DecValTok{0}\NormalTok{, }\DecValTok{0}\NormalTok{)}
\NormalTok{arm.Placement.Base }\OperatorTok{=}\NormalTok{ (}\DecValTok{20}\NormalTok{, }\DecValTok{0}\NormalTok{, }\DecValTok{0}\NormalTok{)}

\CommentTok{\# Logik: Die Position eines Teils ist das, was der Solver in}
\CommentTok{\# einer konsistenten Baugruppe schließlich berechnet.}
\BuiltInTok{print}\NormalTok{(}\StringTok{"Base: "}\NormalTok{, base.Placement.Base)}
\BuiltInTok{print}\NormalTok{(}\StringTok{"Arm: "}\NormalTok{, arm.Placement.Base)}
\end{Highlighting}
\end{Shaded}

Was wichtig ist: Das Beispiel zeigt den Stil der Verbindung zwischen
CAD-Objekten und einer für den Solver relevanten Geometrie- oder
Platzierungsdatenstruktur.

\hypertarget{ein-reales-muster-aus-dem-projekt}{%
\section{\texorpdfstring{{5.5} Ein reales Muster aus dem
Projekt}{5.5 Ein reales Muster aus dem Projekt}}\label{ein-reales-muster-aus-dem-projekt}}

In der echten Implementation werden Assembly-Objekte und Solver-Objekte
strukturell verbunden. Das ist im Build-System sichtbar, in
\texttt{src/Mod/Assembly/App/CMakeLists.txt}:

\begin{Shaded}
\begin{Highlighting}[]
\KeywordTok{set}\NormalTok{(Assembly\_LIBS}
\NormalTok{    Part}
\NormalTok{    PartDesign}
\NormalTok{    Spreadsheet}
\NormalTok{    FreeCADApp}
\NormalTok{    OndselSolver}
\NormalTok{)}

\CommentTok{\# ... weiter unten in derselben Datei:}
\KeywordTok{add\_library}\NormalTok{(}\BaseNTok{Assembly} \OtherTok{SHARED} \DecValTok{$\{Assembly\_SRCS\}}\NormalTok{)}
\KeywordTok{target\_link\_libraries}\NormalTok{(}\BaseNTok{Assembly} \DecValTok{$\{Assembly\_LIBS\}}\NormalTok{)}
\end{Highlighting}
\end{Shaded}

OndselSolver steht also nicht direkt in der
\texttt{target\_link\_libraries}-Zeile, sondern in der Variable
\texttt{Assembly\_LIBS}. Diese Variable wird beim Linken eingesetzt, und
damit wird das Assembly-Modul gegen OndselSolver gelinkt.

Das ist die zentrale Stelle für den gesamten Abschnitt:

\begin{itemize}
\tightlist
\item
  \texttt{Assembly} ist ein FreeCAD-Modul
\item
  \texttt{OndselSolver} ist die Solver-Bibliothek
\item
  beide werden zu einer gemeinsamen Laufzeit-/Link-Umgebung kombiniert
\end{itemize}

Die Header des Solvers findet der Compiler über
\texttt{src/Mod/Assembly/CMakeLists.txt}. Dort wird entweder die
mitgelieferte Kopie in \texttt{src/3rdParty/OndselSolver} eingebunden,
oder, wenn \texttt{FREECAD\_USE\_EXTERNAL\_ONDSELSOLVER} gesetzt ist,
eine im System installierte Version.

\hypertarget{kopflose-ausfuxfchrung-headless-von-freecad}{%
\section{\texorpdfstring{{5.6} Kopflose Ausführung (Headless) von
FreeCAD}{5.6 Kopflose Ausführung (Headless) von FreeCAD}}\label{kopflose-ausfuxfchrung-headless-von-freecad}}

Für Automatisierung oder Server-/CI-Umgebungen ist ein entscheidender
Fall die sogenannte headless Ausführung: FreeCAD wird ohne GUI
gestartet. Das ist besonders hilfreich, wenn eine Baugruppe aus einer
Python-Datei berechnet oder im Hintergrund aufgelöst werden soll.

Typischer Ablauf:

\begin{enumerate}
\def\labelenumi{\arabic{enumi}.}
\tightlist
\item
  \texttt{freecad.App} wird mit dem Unterbefehl \texttt{freecadcmd}
  gestartet
\item
  eine Python-Datei wird geladen
\item
  die Baugruppe wird erzeugt oder geöffnet
\item
  die Solver-Logik wird ausgeführt
\item
  das Ergebnis wird in die Konsole oder in Dateien geschrieben
\end{enumerate}

Ein typisches Muster sieht so aus:

\begin{Shaded}
\begin{Highlighting}[]
\ExtensionTok{\textasciitilde{}/freecad/App\_weekly/freecad.App}\NormalTok{ freecadcmd }\AttributeTok{{-}c} \StringTok{"import FreeCAD as App; doc = App.newDocument(\textquotesingle{}HeadlessDemo\textquotesingle{}); print(doc.Name)"}
\end{Highlighting}
\end{Shaded}

Der Unterbefehl \texttt{freecadcmd} sorgt dafür, dass das AppImage
FreeCAD ohne GUI startet. Eine separate Installation von
\texttt{FreeCADCmd} ist nicht nötig.

Das ist der entscheidende Unterschied zu einer GUI-Version: Es gibt kein
sichtbares Fenster, aber dieselbe CAD-/Solver-Logik kann verwendet
werden.

\hypertarget{ein-konkretes-beispiel-fuxfcr-headless-nutzung}{%
\section{\texorpdfstring{{5.7} Ein konkretes Beispiel für headless
Nutzung}{5.7 Ein konkretes Beispiel für headless Nutzung}}\label{ein-konkretes-beispiel-fuxfcr-headless-nutzung}}

Im Buch liegt ein Beispiel unter:

\begin{itemize}
\tightlist
\item
  \texttt{examples/05-freecad-headless/}
\end{itemize}

Das Beispiel zeigt genau diesen Kopf-Los-Start und setzt ein
Python-Skript mit der FreeCAD-API in Bewegung. Das ist die bevorzugte
Ausführungsform für Automatisierung, Testläufe und Build-/CI-Szenarien.

\hypertarget{warum-das-fuxfcr-ondselsolver-relevant-ist}{%
\section{\texorpdfstring{{5.8} Warum das für OndselSolver relevant
ist}{5.8 Warum das für OndselSolver relevant ist}}\label{warum-das-fuxfcr-ondselsolver-relevant-ist}}

Die Solver-Integration wirkt nur dann wirklich sinnvoll, wenn das
CAD-System die Daten sauber liefert und die Solver-Ergebnisse sauber in
die CAD-Welt zurückschreibt. Genau diese Abbildung ist die eigentliche
Ingenieursarbeit hinter FreeCAD + OndselSolver.

Die wichtigsten Fragen lauten deshalb:

\begin{itemize}
\tightlist
\item
  Welche Objekte werden aus der CAD-Welt an den Solver gegeben?
\item
  Welche Joints oder Constraints wurden modelliert?
\item
  Wie werden die Berechnungsergebnisse in die Ausgabe zurück übertragen?
\item
  Wie funktioniert das in einer headless Umgebung ohne GUI?
\end{itemize}

Diese Fragen sind der Kern des Verständnisses.

\hypertarget{fazit-3}{%
\section{\texorpdfstring{{5.9} Fazit}{5.9 Fazit}}\label{fazit-3}}

FreeCAD ist in der Praxis kein „einfaches`` 3D-Programm, sondern eine
ingenieurtechnische Plattform mit mehreren Bibliotheken und klarer
Modulstruktur. Die Kombination aus FreeCADApp, Part, PartDesign,
Assembly und OndselSolver ist der eigentliche Kern der
Mechanik-/Assembliestruktur.

Daraus ergibt sich eine einfache Wahrheit:

Der Solver wird nicht einfach „dazugeschaltet``, sondern als Teil einer
vollständigen CAD-/Modellierungs-Architektur eingebettet.

\hypertarget{kap-formatbeschreibung-asmt}{%
\chapter{\texorpdfstring{{6} Formatbeschreibung
ASMT}{6 Formatbeschreibung ASMT}}\label{kap-formatbeschreibung-asmt}}

\hypertarget{was-ist-asmt}{%
\section{\texorpdfstring{{6.1} Was ist
ASMT?}{6.1 Was ist ASMT?}}\label{was-ist-asmt}}

ASMT ist ein Assemblies-/Solver-Format, das als Austausch- und
Zwischendarstellung zwischen einer FreeCAD-Äquivalenzmodellierung und
dem OndselSolver genutzt werden kann. Es ist kein rein „grafisches``
Format, sondern ein strukturiertes Datenformat für Mechanik- und
Kinematik-Modelle.

In diesem Projekt ist ASMT bereits als Vorstellung und
Export-Mechanismus sichtbar:

\begin{itemize}
\tightlist
\item
  \href{https://github.com/FreeCAD/FreeCAD/blob/main/src/Mod/Assembly/CommandExportASMT.py}{CommandExportASMT.py}
  enthält den Exportbefehl für \texttt{.asmt}-Dateien
\item
  \href{https://github.com/FreeCAD/FreeCAD/blob/main/src/Mod/Assembly/Init.py}{Init.py}
  registriert \texttt{.asmt} als Importtyp
\item
  \texttt{vendor/OndselSolver/OndselSolver/ASMTAssembly.cpp} und
  verwandte Klassen im Submodul modellieren das Solver-seitige Format
\end{itemize}

Damit ist ASMT ein sehr konkretes Bindeglied zwischen FreeCAD-Modell und
OndselSolver-Repräsentation.

\hypertarget{warum-ein-separates-format-sinnvoll-ist}{%
\section{\texorpdfstring{{6.2} Warum ein separates Format sinnvoll
ist}{6.2 Warum ein separates Format sinnvoll ist}}\label{warum-ein-separates-format-sinnvoll-ist}}

Ein 3D-Modell in FreeCAD enthält oft mehr als nur mechanische
Beziehungen. Es enthält:

\begin{itemize}
\tightlist
\item
  geometrische Bauteile,
\item
  Platzierung und Transformation,
\item
  Joints und Constraints,
\item
  Gruppen und Bezugssysteme,
\item
  zusätzliche Metadaten der Bearbeitung.
\end{itemize}

Der Solver benötigt aber meist eine klarere, reduzierte und
strukturierte Darstellung. Genau dafür ist ASMT nützlich: Es beschreibt
die Assembly-Informationen in einer Form, die für die Berechnung besser
geeignet ist.

\hypertarget{typische-inhalte-eines-asmt-modells}{%
\section{\texorpdfstring{{6.3} Typische Inhalte eines
ASMT-Modells}{6.3 Typische Inhalte eines ASMT-Modells}}\label{typische-inhalte-eines-asmt-modells}}

Ein ASMT-Modell enthält typischerweise eine Kombination aus:

\begin{itemize}
\tightlist
\item
  Teilen (\texttt{part}, \texttt{body}, \texttt{item}-ähnliche Objekte)
\item
  Verbindungen (\texttt{joint}, \texttt{constraint})
\item
  Bezugssystemen (\texttt{frame}, \texttt{marker},
  \texttt{coordinate\ system})
\item
  Platzierungen (\texttt{position}, \texttt{rotation},
  \texttt{orientation})
\item
  Informationen über die Beziehung zwischen den Teilen
\end{itemize}

Das Format soll also nicht nur Geometrie speichern, sondern auch die
Kinematik- und Constraints-Semantik.

\hypertarget{das-asmt-modell-aus-sicht-der-solver-architektur}{%
\section{\texorpdfstring{{6.4} Das ASMT-Modell aus Sicht der
Solver-Architektur}{6.4 Das ASMT-Modell aus Sicht der Solver-Architektur}}\label{das-asmt-modell-aus-sicht-der-solver-architektur}}

Im OndselSolver-Code sind Strukturen wie \texttt{ASMTAssembly},
\texttt{ASMTPart}, \texttt{ASMTJoint} oder ähnliche Klasen eine
unmittelbare signalisierte Sichtweise auf eine Assembly-Definition.
Diese Klassen sind nicht zufällig benannt: Sie modellieren die gleiche
Idee wie in FreeCAD, nur in einer formelleren, solverorientierten Weise.

Das bedeutet:

\begin{itemize}
\tightlist
\item
  Ein Assembly-Teil in FreeCAD soll im ASMT-Format als ein solverfähiger
  Teil erscheinen.
\item
  Ein Joint in FreeCAD wird als ASMT-Joint bzw. Constraint-Klasse ins
  Modell übertragen.
\item
  Die Positionsergebnisse des Solvers werden später wieder in die
  FreeCAD-Objekte zurückgegeben.
\end{itemize}

\hypertarget{ein-beispiel-fuxfcr-die-struktur-eines-asmt-dokuments}{%
\section{\texorpdfstring{{6.5} Ein Beispiel für die Struktur eines
ASMT-Dokuments}{6.5 Ein Beispiel für die Struktur eines ASMT-Dokuments}}\label{ein-beispiel-fuxfcr-die-struktur-eines-asmt-dokuments}}

ASMT ist in der Praxis für Menschen meist nicht komplett
„handgeschriebenes XML``, sondern ein lesbares, strukturiertes
Textformat. Die genaue Ausformung kann je nach Implementierung
variieren, aber die Idee bleibt konstant:

\begin{Shaded}
\begin{Highlighting}[]
\NormalTok{Assembly \{}
\NormalTok{  Name: ExampleAssembly}

\NormalTok{  Part \{}
\NormalTok{    Name: Base}
\NormalTok{    Position: 0 0 0}
\NormalTok{  \}}

\NormalTok{  Part \{}
\NormalTok{    Name: Arm}
\NormalTok{    Position: 10 0 0}
\NormalTok{  \}}

\NormalTok{  Joint \{}
\NormalTok{    Type: Revolute}
\NormalTok{    PartA: Base}
\NormalTok{    PartB: Arm}
\NormalTok{    Axis: Z}
\NormalTok{  \}}
\NormalTok{\}}
\end{Highlighting}
\end{Shaded}

Wichtig ist hier nicht die exakte Syntax eines bestimmten Standards,
sondern das Konzept:

\begin{itemize}
\tightlist
\item
  Teile werden benannt
\item
  Beziehungen werden als Joints oder Constraints beschrieben
\item
  die relative Lage wird durch Formulierung der Verbindung definiert
\end{itemize}

\hypertarget{wie-asmt-mit-freecad-zusammenhuxe4ngt}{%
\section{\texorpdfstring{{6.6} Wie ASMT mit FreeCAD
zusammenhängt}{6.6 Wie ASMT mit FreeCAD zusammenhängt}}\label{wie-asmt-mit-freecad-zusammenhuxe4ngt}}

ASMT ist die Brücke zwischen:

\begin{itemize}
\tightlist
\item
  dem verständlichen, graphischen FreeCAD-Modell,
\item
  und der numerischen, solverorientierten Darstellung.
\end{itemize}

FreeCAD erzeugt oder exportiert die Assembly-Daten in ein Format, das
der Solver semantisch verarbeiten kann. Danach kann der Solver die
Positions- und Orientierungswerte berechnen und in die FreeCAD-Objekte
zurückführen.

Diese Rolle ist in \texttt{CommandExportASMT.py} deutlich sichtbar: Es
ist eine explizite Export-Aktion, die auf eine existierende Assembly
angewendet wird.

\hypertarget{warum-asmt-fuxfcr-lernende-wichtig-ist}{%
\section{\texorpdfstring{{6.7} Warum ASMT für Lernende wichtig
ist}{6.7 Warum ASMT für Lernende wichtig ist}}\label{warum-asmt-fuxfcr-lernende-wichtig-ist}}

Ein Leser, der das Buch ernsthaft verstehen will, sollte ASMT als „eine
zugeschnittene, solverfreundliche Version der
CAD-Assembly-Informationen`` verstehen. Das hilft bei drei Dingen:

\begin{itemize}
\tightlist
\item
  die Datenstruktur zu verstehen,
\item
  die Rollen von Teilen, Gelenken und Constraints zu erkennen,
\item
  die Verbindung zwischen Benutzeroberfläche und numerischer Berechnung
  zu sehen.
\end{itemize}

\hypertarget{ein-wichtiger-realistischer-punkt}{%
\section{\texorpdfstring{{6.8} Ein wichtiger realistischer
Punkt}{6.8 Ein wichtiger realistischer Punkt}}\label{ein-wichtiger-realistischer-punkt}}

In diesem Projekt selbst ist ASMT noch als
Forschungs-/Interoperabilitätsformat sichtbar, nicht als vollständig
ausgereiftes, universell abgeschlossenes Dateiformat. Das zeigt die
\texttt{Init.py}-Kommentierung sehr klar: „The correct format for
assembly interoperability is a research topic. ASMT is a placeholder.``

Das heißt:

\begin{itemize}
\tightlist
\item
  ASMT ist ein sehr wichtiges, relevantes Konzept,
\item
  aber es ist noch keine „fertige, abschließende universelle
  Spezifikation`` in dem Sinne, wie ein vollständiger Dateistandard sein
  könnte.
\item
  Für Lernzwecke ist es trotzdem ein sehr nützliches Denkmodell.
\end{itemize}

\hypertarget{fazit-4}{%
\section{\texorpdfstring{{6.9} Fazit}{6.9 Fazit}}\label{fazit-4}}

ASMT ist die sprachliche und strukturelle Brücke zwischen einem
visuellen CAD-Modell und dem mathematischen Mehrkörpersystem des
Solvers. Es macht aus einer Baugruppe eine strukturierte, berechenbare
Definition aus Teilen, Beziehungen und Constraints.

Für den Lernenden bedeutet das:

\begin{itemize}
\tightlist
\item
  ASMT ist kein bloßes Nebenprodukt,
\item
  sondern ein Schlüsselkonzept für das Verständnis von FreeCAD +
  OndselSolver.
\end{itemize}

Das ist genau der Punkt, an dem der Übergang von „Objektmodellierung``
zu „solverfähiger Mechanik`` sichtbar wird.

\hypertarget{kap-asmt-flow}{%
\chapter{\texorpdfstring{{7} Praxisfluss: FreeCAD → ASMT/Assembly →
OndselSolver → Rückgabe in
FreeCAD}{7 Praxisfluss: FreeCAD → ASMT/Assembly → OndselSolver → Rückgabe in FreeCAD}}\label{kap-asmt-flow}}

\hypertarget{warum-dieses-kapitel-notwendig-ist}{%
\section{\texorpdfstring{{7.1} Warum dieses Kapitel notwendig
ist}{7.1 Warum dieses Kapitel notwendig ist}}\label{warum-dieses-kapitel-notwendig-ist}}

Das bisherige Buch erklärt isolierte Konzepte. Das hier ist der
praktische „End-to-End-Fluss``: Wie die echte Architektur ein
Baugruppenmodell in eine solverfähige Darstellung überführt und danach
die Lösung wieder zurück in die CAD-Welt schreibt.

Das ist der Kern, den ein Lernender wirklich verstehen muss.

\hypertarget{der-komplette-ablauf}{%
\section{\texorpdfstring{{7.2} Der komplette
Ablauf}{7.2 Der komplette Ablauf}}\label{der-komplette-ablauf}}

Ein realer Flow in einem FreeCAD-Assembly-System läuft im Grundsatz so:

\begin{enumerate}
\def\labelenumi{\arabic{enumi}.}
\tightlist
\item
  Der Nutzer modelliert Teile und Beziehungen in FreeCAD.
\item
  Die Assembly-Logik erkennt die relevanten Objekte und ihre
  Verbindungen.
\item
  Die Baugruppe wird in eine solverorientierte Repräsentation übersetzt.
\item
  Die Übersetzung kann als ASMT- oder Assembly-ähnliches Modell
  erfolgen.
\item
  OndselSolver löst die Beziehungssysteme.
\item
  Die berechneten Platzierungen werden in die FreeCAD-Objekte
  zurückgeschrieben.
\item
  Die Darstellung im 3D-Fenster wird aktualisiert.
\end{enumerate}

Das ist der reale Kreislauf zwischen CAD und Solver.

\hypertarget{beispiel-aus-der-projektstruktur}{%
\section{\texorpdfstring{{7.3} Beispiel aus der
Projektstruktur}{7.3 Beispiel aus der Projektstruktur}}\label{beispiel-aus-der-projektstruktur}}

In FreeCADs Assembly-Implementierung ist die Verbindung konkret im Code
sichtbar:

\begin{itemize}
\tightlist
\item
  \href{https://github.com/FreeCAD/FreeCAD/blob/main/src/Mod/Assembly/CommandExportASMT.py}{CommandExportASMT.py}
  exportiert eine Baugruppe als \texttt{.asmt}
\item
  \href{https://github.com/FreeCAD/FreeCAD/blob/main/src/Mod/Assembly/App/CMakeLists.txt}{Assembly/App/CMakeLists.txt}
  bindet \texttt{OndselSolver} in die Assembly-Library ein
\item
  \href{https://github.com/FreeCAD/FreeCAD/blob/main/src/Mod/Assembly/App/AssemblyObject.cpp}{AssemblyObject.cpp}
  verarbeitet die Solve- und Placement-Logik
\end{itemize}

Das zeigt bereits den vollständigen gedachten Ablauf:

\begin{itemize}
\tightlist
\item
  FreeCAD erzeugt das Modell
\item
  Assembly kapselt und übersetzt es
\item
  OndselSolver berechnet die Lösung
\item
  FreeCAD zeigt das Ergebnis an
\end{itemize}

\hypertarget{ein-konkretes-vereinfachtes-beispiel}{%
\section{\texorpdfstring{{7.4} Ein konkretes, vereinfachtes
Beispiel}{7.4 Ein konkretes, vereinfachtes Beispiel}}\label{ein-konkretes-vereinfachtes-beispiel}}

Das folgende Beispiel zeigt den Ablauf ohne die komplette FreeCAD-API
und ohne die reale Solver-Bibliothek zu kopieren. Es ist eine
didaktische Darstellung.

\begin{Shaded}
\begin{Highlighting}[]
\ImportTok{import}\NormalTok{ FreeCAD }\ImportTok{as}\NormalTok{ App}

\CommentTok{\# 1) FreeCAD{-}Modell erzeugen}
\CommentTok{\# Ein Dokument wird als Assembly{-}Kontext gestartet.}
\NormalTok{doc }\OperatorTok{=}\NormalTok{ App.newDocument(}\StringTok{"TandemDemo"}\NormalTok{)}

\NormalTok{base }\OperatorTok{=}\NormalTok{ doc.addObject(}\StringTok{"Part::Box"}\NormalTok{, }\StringTok{"Base"}\NormalTok{)}
\NormalTok{base.Length }\OperatorTok{=} \DecValTok{40}
\NormalTok{base.Width }\OperatorTok{=} \DecValTok{20}
\NormalTok{base.Height }\OperatorTok{=} \DecValTok{10}

\NormalTok{arm }\OperatorTok{=}\NormalTok{ doc.addObject(}\StringTok{"Part::Box"}\NormalTok{, }\StringTok{"Arm"}\NormalTok{)}
\NormalTok{arm.Length }\OperatorTok{=} \DecValTok{30}
\NormalTok{arm.Width }\OperatorTok{=} \DecValTok{10}
\NormalTok{arm.Height }\OperatorTok{=} \DecValTok{10}
\NormalTok{arm.Placement.Base }\OperatorTok{=}\NormalTok{ (}\DecValTok{20}\NormalTok{, }\DecValTok{0}\NormalTok{, }\DecValTok{0}\NormalTok{)}

\CommentTok{\# 2) In der echten Architektur würde hier ein Joint oder eine Constraint{-}Definition}
\CommentTok{\#    entstehen, z. B. eine Achsen{-} oder Abstandsbeschränkung.}
\CommentTok{\#    Der Assembly{-}Code sammelt diese Daten und gibt sie an den Solver weiter.}

\NormalTok{assembly\_data }\OperatorTok{=}\NormalTok{ \{}
    \StringTok{"parts"}\NormalTok{: [}\StringTok{"Base"}\NormalTok{, }\StringTok{"Arm"}\NormalTok{],}
    \StringTok{"constraints"}\NormalTok{: [}
\NormalTok{        \{}\StringTok{"type"}\NormalTok{: }\StringTok{"distance"}\NormalTok{, }\StringTok{"part\_a"}\NormalTok{: }\StringTok{"Base"}\NormalTok{, }\StringTok{"part\_b"}\NormalTok{: }\StringTok{"Arm"}\NormalTok{, }\StringTok{"target"}\NormalTok{: }\FloatTok{25.0}\NormalTok{\}}
\NormalTok{    ]}
\NormalTok{\}}

\BuiltInTok{print}\NormalTok{(}\StringTok{"Assembly{-}Daten erstellt:"}\NormalTok{)}
\BuiltInTok{print}\NormalTok{(assembly\_data)}

\CommentTok{\# 3) Solver{-}ähnliche Darstellung}
\NormalTok{solver\_input }\OperatorTok{=}\NormalTok{ \{}
    \StringTok{"bodies"}\NormalTok{: [}
\NormalTok{        \{}\StringTok{"id"}\NormalTok{: }\StringTok{"Base"}\NormalTok{, }\StringTok{"position"}\NormalTok{: [}\DecValTok{0}\NormalTok{, }\DecValTok{0}\NormalTok{, }\DecValTok{0}\NormalTok{]\},}
\NormalTok{        \{}\StringTok{"id"}\NormalTok{: }\StringTok{"Arm"}\NormalTok{, }\StringTok{"position"}\NormalTok{: [}\DecValTok{20}\NormalTok{, }\DecValTok{0}\NormalTok{, }\DecValTok{0}\NormalTok{]\},}
\NormalTok{    ],}
    \StringTok{"constraints"}\NormalTok{: [\{}\StringTok{"kind"}\NormalTok{: }\StringTok{"distance"}\NormalTok{, }\StringTok{"target"}\NormalTok{: }\FloatTok{25.0}\NormalTok{\}]}
\NormalTok{\}}

\BuiltInTok{print}\NormalTok{(}\StringTok{"Solver{-}Input vorbereitet:"}\NormalTok{)}
\BuiltInTok{print}\NormalTok{(solver\_input)}

\CommentTok{\# 4) Ergebnisnachbereitung im FreeCAD{-}Kontext}
\CommentTok{\# In der echten Pipeline würde der Solver neue Positionen berechnen und}
\CommentTok{\# diese Werte in die Plazierungen zurückschreiben.}
\NormalTok{arm.Placement.Base }\OperatorTok{=}\NormalTok{ (}\DecValTok{25}\NormalTok{, }\DecValTok{0}\NormalTok{, }\DecValTok{0}\NormalTok{)}

\BuiltInTok{print}\NormalTok{(}\StringTok{"Arm{-}Platzierung nach Solve:"}\NormalTok{, arm.Placement.Base)}
\end{Highlighting}
\end{Shaded}

Wichtig ist hier die Reihenfolge:

\begin{itemize}
\tightlist
\item
  FreeCAD modelliert die Geometrie
\item
  Assembly definiert die Beziehungen
\item
  Solver verarbeitet die Beziehung
\item
  Ergebnis wird wieder in das CAD-Modell geschrieben
\end{itemize}

\hypertarget{der-dreiteilige-blick-auf-den-ablauf}{%
\section{\texorpdfstring{{7.5} Der dreiteilige Blick auf den
Ablauf}{7.5 Der dreiteilige Blick auf den Ablauf}}\label{der-dreiteilige-blick-auf-den-ablauf}}

\hypertarget{cad-seite}{%
\subsection{\texorpdfstring{{7.5.1}
CAD-Seite}{7.5.1 CAD-Seite}}\label{cad-seite}}

Die CAD-Seite kennt die Geometrie, die Objekte und die Sichtbarkeit. Sie
ist die Benutzerschnittstelle und das Modell des Baugruppenkontexts.

\hypertarget{assembly-asmt-seite}{%
\subsection{\texorpdfstring{{7.5.2}
Assembly-/ASMT-Seite}{7.5.2 Assembly-/ASMT-Seite}}\label{assembly-asmt-seite}}

Diese Ebene formt die CAD-Daten in eine klare, solverorientierte
Struktur um. Dabei ist ASMT das sichtbarste Format, aber der eigentliche
Gedanke ist die gleiche Haltung: eine definierte, berechenbare
Repräsentation der Baugruppe.

\hypertarget{solver-seite}{%
\subsection{\texorpdfstring{{7.5.3}
Solver-Seite}{7.5.3 Solver-Seite}}\label{solver-seite}}

Das Mehrkörpersystem löst die Beziehungen, korrigiert die Positionen und
liefert eine konsistente Konfiguration zurück.

\hypertarget{warum-die-ruxfcckgabe-in-freecad-so-wichtig-ist}{%
\section{\texorpdfstring{{7.6} Warum die Rückgabe in FreeCAD so wichtig
ist}{7.6 Warum die Rückgabe in FreeCAD so wichtig ist}}\label{warum-die-ruxfcckgabe-in-freecad-so-wichtig-ist}}

Wenn die Lösung berechnet ist, muss sie wieder in die CAD-Darstellung
zurückfließen. Nur dann ist sie für den Benutzer sichtbar und nutzbar.
Genau das passiert in der echten Architektur durch die Update- und
Placement-Mechanik der Assembly-Objekte.

Ohne diese Rückübertragung wäre die Lösung nur eine numerische
Nebenwirkung, kein nutzbares CAD-Ergebnis.

\hypertarget{praktische-lernfrage}{%
\section{\texorpdfstring{{7.7} Praktische
Lernfrage}{7.7 Praktische Lernfrage}}\label{praktische-lernfrage}}

Ein Leser sollte bei jeder Lösung sich immer fragen:

\begin{itemize}
\tightlist
\item
  Wo beginnt das Modell?
\item
  Welche Daten werden zum Solver geschickt?
\item
  Wie funktioniert die Übersetzung in die Assembly-/ASMT-Struktur?
\item
  Was ist der Fehlerwert und wie wird er gelöst?
\item
  Wie werden die berechneten Werte wieder in FreeCAD gesetzt?
\end{itemize}

Diese Fragen machen den gesamten Ablauf verständlich.

\hypertarget{fazit-5}{%
\section{\texorpdfstring{{7.8} Fazit}{7.8 Fazit}}\label{fazit-5}}

Der echte Flow ist kein linearer Einzeiler, sondern ein geschlossener
Kreis:

FreeCAD → Assembly/ASMT → OndselSolver → Ergebnis zurück in FreeCAD

Das ist die wichtigste praktische Sichtweise auf OndselSolver in einer
CAD-Umgebung. Wenn man diesen Kreislauf verstanden hat, dann ist das
gesamte Buch in seiner Grundlogik verankert.

\hypertarget{kap-freecad-app}{%
\chapter{\texorpdfstring{{8} FreeCAD AppImage unter Ubuntu und
Dash-Icon}{8 FreeCAD AppImage unter Ubuntu und Dash-Icon}}\label{kap-freecad-app}}

\hypertarget{appimage-unter-ubuntu-einrichten}{%
\section{\texorpdfstring{{8.1} AppImage unter Ubuntu
einrichten}{8.1 AppImage unter Ubuntu einrichten}}\label{appimage-unter-ubuntu-einrichten}}

Ein AppImage wird nicht mit \texttt{apt} installiert. Es ist eine
ausführbare, portable Datei. Lege die offizielle
FreeCAD-Linux-x86\_64-AppImage in einem dauerhaften Ordner ab, zum
Beispiel als \texttt{\textasciitilde{}/freecad/App\_weekly/freecad.App}.
Mache sie ausführbar und starte sie:

\begin{Shaded}
\begin{Highlighting}[]
\FunctionTok{mkdir} \AttributeTok{{-}p} \StringTok{"}\VariableTok{$HOME}\StringTok{/freecad/App\_weekly"}
\FunctionTok{chmod}\NormalTok{ +x }\StringTok{"}\VariableTok{$HOME}\StringTok{/freecad/App\_weekly/freecad.App"}
\StringTok{"}\VariableTok{$HOME}\StringTok{/freecad/App\_weekly/freecad.App"}
\end{Highlighting}
\end{Shaded}

Für die wöchentlich aktualisierte AppImage liegt im Buch das Skript
\texttt{examples/05-freecad-headless/update-weekly-appimage.sh}. Es
benötigt die GitHub-CLI \texttt{gh}, lädt das offizielle Weekly-AppImage
samt SHA256-Prüfung und speichert es als
\texttt{\textasciitilde{}/freecad/App\_weekly/freecad.App}:

\begin{Shaded}
\begin{Highlighting}[]
\ExtensionTok{examples/05{-}freecad{-}headless/update{-}weekly{-}appimage.sh}
\end{Highlighting}
\end{Shaded}

Ein anderer Zielordner lässt sich über \texttt{FREECAD\_APP\_DIR}
angeben.

Falls Ubuntu meldet, dass \texttt{libfuse.so.2} fehlt, installiere das
passende FUSE-2-Kompatibilitätspaket:

\begin{Shaded}
\begin{Highlighting}[]
\CommentTok{\# Ubuntu 24.04 und neuer}
\FunctionTok{sudo}\NormalTok{ apt install libfuse2t64}

\CommentTok{\# Ältere Ubuntu{-}Versionen, falls libfuse2t64 nicht verfügbar ist}
\FunctionTok{sudo}\NormalTok{ apt install libfuse2}
\end{Highlighting}
\end{Shaded}

\hypertarget{dash-eintrag-und-icon-erstellen}{%
\section{\texorpdfstring{{8.2} Dash-Eintrag und Icon
erstellen}{8.2 Dash-Eintrag und Icon erstellen}}\label{dash-eintrag-und-icon-erstellen}}

Zuerst wird das FreeCAD-Icon aus dem AppImage in das Benutzerprofil
kopiert:

\begin{Shaded}
\begin{Highlighting}[]
\VariableTok{APPIMAGE}\OperatorTok{=}\StringTok{"}\VariableTok{$HOME}\StringTok{/freecad/App\_weekly/freecad.App"}
\VariableTok{ICON\_DIR}\OperatorTok{=}\StringTok{"}\VariableTok{$HOME}\StringTok{/.local/share/icons"}
\VariableTok{DESKTOP\_DIR}\OperatorTok{=}\StringTok{"}\VariableTok{$HOME}\StringTok{/.local/share/applications"}
\VariableTok{TMP\_DIR}\OperatorTok{=}\StringTok{"}\VariableTok{$(}\FunctionTok{mktemp} \AttributeTok{{-}d}\VariableTok{)}\StringTok{"}

\FunctionTok{mkdir} \AttributeTok{{-}p} \StringTok{"}\VariableTok{$ICON\_DIR}\StringTok{"} \StringTok{"}\VariableTok{$DESKTOP\_DIR}\StringTok{"}
\KeywordTok{(}
  \BuiltInTok{cd} \StringTok{"}\VariableTok{$TMP\_DIR}\StringTok{"}
  \StringTok{"}\VariableTok{$APPIMAGE}\StringTok{"} \AttributeTok{{-}{-}appimage{-}extract}
\KeywordTok{)}
\FunctionTok{install} \AttributeTok{{-}m}\NormalTok{ 644 }\StringTok{"}\VariableTok{$TMP\_DIR}\StringTok{/squashfs{-}root/.DirIcon"} \StringTok{"}\VariableTok{$ICON\_DIR}\StringTok{/freecad{-}weekly.svg"}
\FunctionTok{rm} \AttributeTok{{-}rf} \StringTok{"}\VariableTok{$TMP\_DIR}\StringTok{"}
\end{Highlighting}
\end{Shaded}

Danach den Desktop-Eintrag erstellen:

\begin{Shaded}
\begin{Highlighting}[]
\FunctionTok{cat} \OperatorTok{\textgreater{}} \StringTok{"}\VariableTok{$DESKTOP\_DIR}\StringTok{/freecad{-}weekly.desktop"} \OperatorTok{\textless{}\textless{}EOF}
\StringTok{[Desktop Entry]}
\StringTok{Type=Application}
\StringTok{Name=FreeCAD Weekly (App)}
\StringTok{Comment=FreeCAD Weekly AppImage}
\StringTok{Exec=}\VariableTok{$HOME}\StringTok{/freecad/App\_weekly/freecad.App \%F}
\StringTok{Icon=}\VariableTok{$ICON\_DIR}\StringTok{/freecad{-}weekly.svg}
\StringTok{Terminal=false}
\StringTok{Categories=Graphics;Engineering;Science;}
\StringTok{MimeType=application/x{-}extension{-}fcstd;}
\OperatorTok{EOF}

\FunctionTok{chmod}\NormalTok{ +x }\StringTok{"}\VariableTok{$DESKTOP\_DIR}\StringTok{/freecad{-}weekly.desktop"}
\end{Highlighting}
\end{Shaded}

Drücke die Super-Taste und suche nach „FreeCAD Weekly (App)``. Über das
Kontextmenü kann der Eintrag zu den Favoriten hinzugefügt werden. Falls
die Suche den neuen Eintrag nicht sofort zeigt, melde dich einmal ab und
wieder an.

\hypertarget{beispiele-starten}{%
\section{\texorpdfstring{{8.3} Beispiele
starten}{8.3 Beispiele starten}}\label{beispiele-starten}}

Die GUI-Beispiele verwenden
\texttt{\textasciitilde{}/freecad/App\_weekly/freecad.App}:

\begin{Shaded}
\begin{Highlighting}[]
\BuiltInTok{cd}\NormalTok{ examples/04{-}freecad{-}pipeline}
\ExtensionTok{./run{-}app.sh}

\BuiltInTok{cd}\NormalTok{ ../05{-}freecad{-}headless}
\ExtensionTok{./run.sh}
\end{Highlighting}
\end{Shaded}

Die Beispiele \texttt{01} bis \texttt{03} und \texttt{06} sind bewusst
eigenständige C++-Konsolenprogramme. Sie demonstrieren Rechenideen und
ASMT-ähnliche Datenstrukturen und benötigen kein FreeCAD-AppImage. Das
Verzeichnis \texttt{05-freecad-headless} bietet mit \texttt{run.sh}
einen GUI-Smoke-Test und mit \texttt{run-headless.sh} einen optionalen
automatisierten Lauf ohne GUI. Beide verwenden dasselbe
\texttt{freecad.App}; für den Lauf ohne GUI wird es mit dem Unterbefehl
\texttt{freecadcmd} gestartet.

Für ein AppImage an einem anderen Ort kann der Pfad beim Start
überschrieben werden:

\begin{Shaded}
\begin{Highlighting}[]
\VariableTok{FREECAD\_APP}\OperatorTok{=}\StringTok{"}\VariableTok{$HOME}\StringTok{/Downloads/freecad.App"} \ExtensionTok{./run{-}app.sh}
\end{Highlighting}
\end{Shaded}

\hypertarget{kap-asmt-beispiele}{%
\chapter{\texorpdfstring{{9} Konkrete OndselSolver-Beispiele mit
ASMT-Dateien}{9 Konkrete OndselSolver-Beispiele mit ASMT-Dateien}}\label{kap-asmt-beispiele}}

\hypertarget{vom-datenmodell-zum-echten-solve}{%
\section{\texorpdfstring{{9.1} Vom Datenmodell zum echten
Solve}{9.1 Vom Datenmodell zum echten Solve}}\label{vom-datenmodell-zum-echten-solve}}

Die vorigen C++-Demos erklären einzelne Ideen. Dieses Beispiel
verarbeitet dagegen echte Dateien aus OndselSolvers
\texttt{testapp}-Sammlung. Es ruft den Originalcode
\texttt{MbD::ASMTAssembly::assemblyFromFile()}, \texttt{solve()} und
\texttt{outputFile()} auf.

Das Beispiel liegt unter
\href{./examples/07-ondselsolver-real-asmt/README.md}{\texttt{examples/07-ondselsolver-real-asmt}}.
CMake baut OndselSolver aus dem auf einen bekannten Stand fixierten
Submodul \texttt{vendor/OndselSolver}.

\hypertarget{die-drei-konkreten-mechanismen}{%
\section{\texorpdfstring{{9.2} Die drei konkreten
Mechanismen}{9.2 Die drei konkreten Mechanismen}}\label{die-drei-konkreten-mechanismen}}

Die drei verwendeten Fixture-Dateien liegen direkt im Beispielordner und
werden beim CMake-Konfigurieren in den Beispiel-Build kopiert:

\begin{longtable}[]{@{}
  >{\raggedright\arraybackslash}p{(\columnwidth - 4\tabcolsep) * \real{0.3333}}
  >{\raggedright\arraybackslash}p{(\columnwidth - 4\tabcolsep) * \real{0.3333}}
  >{\raggedright\arraybackslash}p{(\columnwidth - 4\tabcolsep) * \real{0.3333}}@{}}
\toprule\noalign{}
\begin{minipage}[b]{\linewidth}\raggedright
Datei
\end{minipage} & \begin{minipage}[b]{\linewidth}\raggedright
Mechanismus
\end{minipage} & \begin{minipage}[b]{\linewidth}\raggedright
Relevante Beziehungen
\end{minipage} \\
\midrule\noalign{}
\endhead
\bottomrule\noalign{}
\endlastfoot
\texttt{\_\_cubes.asmt} & zwei Körper & Fixed Joint und Angle Joint \\
\texttt{fourbar.asmt} & Viergelenk & mehrere Gelenke zwischen drei
beweglichen Teilen und der Assembly \\
\texttt{piston.asmt} & Kurbel-Schieber & Gelenke einschließlich eines
zylindrischen Gelenks \\
\end{longtable}

Die Dateien sind echte Solver-Eingaben und keine nachträglich
vereinfachten Pseudodateien. Sie enthalten unter anderem Assembly- und
Part-Positionen, Rotationsmatrizen, RefPoints, Marker,
Masseneigenschaften, Joint-Referenzen und Simulationsparameter.

\hypertarget{eine-asmt-datei-lesen}{%
\section{\texorpdfstring{{9.3} Eine ASMT-Datei
lesen}{9.3 Eine ASMT-Datei lesen}}\label{eine-asmt-datei-lesen}}

Eine ASMT-Datei beginnt mit dem Formatkopf und der Root-Assembly. Danach
folgen eingerückte Abschnitte. Beispielhaft sieht man in
\texttt{\_\_cubes.asmt}:

\begin{Shaded}
\begin{Highlighting}[]
\NormalTok{OndselSolver}
\NormalTok{Assembly}
\NormalTok{    Notes}
\NormalTok{    Name}
\NormalTok{        OndselAssembly}
\NormalTok{    Position3D}
\NormalTok{        0 0 0}
\NormalTok{    RotationMatrix}
\NormalTok{        1 0 0}
\NormalTok{        0 1 0}
\NormalTok{        0 0 1}
\end{Highlighting}
\end{Shaded}

Weiter unten referenzieren Joints die Marker ihrer beiden Seiten. Ein
Fixed Joint fixiert den ersten Körper an der Assembly; das Angle Joint
beschreibt die Winkelbeziehung zum zweiten Körper. Die Marker sind
lokale Bezugssysteme, keine sichtbaren CAD-Körper.

Die Geometrie ist in diesen Solver-Fixtures nicht als vollständiger
CAD-B-Rep enthalten. Das Beispiel demonstriert daher den echten
Kinematik- und Constraint-Solve mit konkreten Körpern, Markern und
Gelenken, nicht die Darstellung von Volumenkörpern.

\hypertarget{bauen-luxf6sen-und-ergebnis-schreiben}{%
\section{\texorpdfstring{{9.4} Bauen, lösen und Ergebnis
schreiben}{9.4 Bauen, lösen und Ergebnis schreiben}}\label{bauen-luxf6sen-und-ergebnis-schreiben}}

Im Repository-Root:

\begin{Shaded}
\begin{Highlighting}[]
\ExtensionTok{./examples/07{-}ondselsolver{-}real{-}asmt/run.sh}
\ExtensionTok{./examples/07{-}ondselsolver{-}real{-}asmt/run.sh}\NormalTok{ fourbar.asmt}
\ExtensionTok{./examples/07{-}ondselsolver{-}real{-}asmt/run.sh}\NormalTok{ piston.asmt}
\end{Highlighting}
\end{Shaded}

Der erste Aufruf baut die Solver-Bibliothek und den kleinen Runner.
Danach wird die ausgewählte \texttt{.asmt}-Datei geladen,
\texttt{ASMTAssembly::solve()} führt den Solve der Anfangslage aus und
\texttt{outputFile()} serialisiert das Ergebnis erneut.

Die Ergebnisse liegen im jeweiligen Build-Verzeichnis:

\begin{itemize}
\tightlist
\item
  \texttt{build/\_\_cubes-solved.asmt}
\item
  \texttt{build/fourbar-solved.asmt}
\item
  \texttt{build/piston-solved.asmt}
\end{itemize}

Im Log erscheinen die Solver-Konvergenzwerte. Bei
\texttt{\_\_cubes.asmt} meldet der Solver außerdem eine entfernte
redundante Winkelbedingung; bei \texttt{fourbar} und \texttt{piston}
sind mehrere Iterationswerte sichtbar, die gegen die Solve-Toleranz
konvergieren.

\hypertarget{installierte-bibliothek-und-reproduzierbarer-build}{%
\section{\texorpdfstring{{9.5} Installierte Bibliothek und
reproduzierbarer
Build}{9.5 Installierte Bibliothek und reproduzierbarer Build}}\label{installierte-bibliothek-und-reproduzierbarer-build}}

In der aktuellen Sandbox-Installation stürzte ein separater Runner gegen
\texttt{\textasciitilde{}/freecad-sandbox/install/lib/libOndselSolver.so}
beim ASMT-Marker-Parsing mit einem Heap-Fehler ab. Derselbe Fixture-Lauf
war erfolgreich, sobald OndselSolver frisch aus der Workspace-Source
gebaut wurde. Deshalb baut dieses Beispiel die Source-Version mit und
linkt direkt dagegen.

Das ist für Lernende auch ein wichtiger Integrationspunkt: Bei
C++-Bibliotheken müssen Header, Shared Library und ABI aus zueinander
passenden Builds stammen. Ein erfolgreicher Link allein beweist noch
nicht, dass ein installierter Solver-Build korrekt zum aktuellen
Source-Stand passt.

\hypertarget{ondselsolver-programmierschnittstelle}{%
\chapter{\texorpdfstring{{Teil II --}
OndselSolver-Programmierschnittstelle}{Teil II -- OndselSolver-Programmierschnittstelle}}\label{ondselsolver-programmierschnittstelle}}

Dieser Teil beschreibt die fachliche ASMT-Anwendungsoberfläche und
ordnet die wichtigsten Funktionsfamilien des Solverkerns ein. Die vielen
internen Methoden der Numerik- und Matrixklassen sind keine pauschal
stabile Client-API; für konkrete Integrationen ist der Quellstand im
Submodul maßgeblich.

\hypertarget{kap-api-ueberblick}{%
\chapter{\texorpdfstring{{10} Programmierschnittstelle von
OndselSolver}{10 Programmierschnittstelle von OndselSolver}}\label{kap-api-ueberblick}}

\hypertarget{solver-verwenden-der-ablauf}{%
\section{\texorpdfstring{{10.1} Solver verwenden: der
Ablauf}{10.1 Solver verwenden: der Ablauf}}\label{solver-verwenden-der-ablauf}}

Der typische Einsatz hat fünf Schritte: Eingabemodell laden oder in C++
aufbauen, die Assembly prüfen, die passende Analyse starten, Ergebnisse
abfragen und optional als ASMT-Datei schreiben.

\begin{Shaded}
\begin{Highlighting}[]
\NormalTok{flowchart TD}
\NormalTok{    A["ASMT{-}Datei"] {-}{-}\textgreater{} B["ASMTAssembly::assemblyFromFile()"]}
\NormalTok{    C["C++{-}Modell"] {-}{-}\textgreater{} D["ASMTAssembly::With() und Objekte ergänzen"]}
\NormalTok{    B {-}{-}\textgreater{} E["Parts, Joints und Marker prüfen"]}
\NormalTok{    D {-}{-}\textgreater{} E}
\NormalTok{    E {-}{-}\textgreater{} F\{"Analyse wählen"\}}
\NormalTok{    F {-}{-}\textgreater{}|"Anfangskonfiguration"| G["solve(): tstart = tend = 0"]}
\NormalTok{    F {-}{-}\textgreater{}|"Eigene Laufparameter"| H["ASMTSimulationParameters setzen"]}
\NormalTok{    H {-}{-}\textgreater{} I["runKINEMATIC()"]}
\NormalTok{    G {-}{-}\textgreater{} J["Zustände und Frames prüfen"]}
\NormalTok{    I {-}{-}\textgreater{} J}
\NormalTok{    J {-}{-}\textgreater{} K["outputFile() oder Ergebnisse im Programm nutzen"]}
\end{Highlighting}
\end{Shaded}

\begin{enumerate}
\def\labelenumi{\arabic{enumi}.}
\tightlist
\item
  \textbf{Laden oder aufbauen:} \texttt{assemblyFromFile()} lädt ASMT;
  \texttt{With()} und \texttt{add*()} bilden ein Modell in C++.
\item
  \textbf{Prüfen:} Namen, Markerpfade und Joint-Anzahl kontrollieren,
  bevor gerechnet wird.
\item
  \textbf{Analyse wählen:} \texttt{solve()} löst nur die
  Anfangskonfiguration und setzt dabei die Simulationsparameter zurück.
  Für eigene Zeitparameter diese setzen und \texttt{runKINEMATIC()}
  aufrufen.
\item
  \textbf{Ergebnisse lesen:} \texttt{numberOfFrames()},
  \texttt{updateForFrame(index)} und räumliche Getter bedienen die
  Ergebnisserien.
\item
  \textbf{Ausgeben:} \texttt{outputFile()} schreibt die Assembly. Der
  aktuelle GCC-Serializer kann ungültige Joint-Typnamen erzeugen; diese
  Datei ist dann nicht zuverlässig wieder einlesbar.
\end{enumerate}

Der schnellste praktische Einstieg ist das echte Fixture-Beispiel:
\texttt{cd\ examples/07-ondselsolver-real-asmt\ \&\&\ ./run.sh\ fourbar.asmt}.
Die API-Beispiele für Laden, Inspektion und Modellobjekte stehen in
\href{./examples/08-ondselsolver-api/README.md}{\texttt{examples/08-ondselsolver-api}}.

\hypertarget{ziel-und-geltungsbereich}{%
\section{\texorpdfstring{{10.2} Ziel und
Geltungsbereich}{10.2 Ziel und Geltungsbereich}}\label{ziel-und-geltungsbereich}}

Part II beschreibt, wie ein C++-Programm OndselSolver aufruft und
Solver-Assemblies modelliert. Die wichtigste Integrationsschicht dieses
Quellstands ist die ASMT-Objektfamilie (\texttt{ASMTAssembly},
\texttt{ASMTPart}, Marker, Joints, Motions und Parameter). Der Aufruf
\texttt{ASMTAssembly::solve()} übersetzt diese Objekte in das interne
MbD-System und startet die Berechnung.

„Alle Funktionen`` muss hier präzise abgegrenzt werden: Der Quellbaum
enthält Hunderte Klassen und tausende Methoden, darunter
Matrixoperationen, symbolische Ableitungen, Constraint-Gleichungen und
Integrator-Schritte. Diese sind nicht automatisch eine stabile
Anwendungs-API. Die Kapitel behandeln daher vollständig die fachlichen
ASMT-Einstiegspunkte und ordnen den internen Solver nach
Funktionsfamilien ein. Parser-, Serialisierungs- und Diagnosehilfen
werden als Implementierungs-/Erweiterungspunkte markiert, nicht als
notwendige Anwendungsaufrufe.

Maßgebliche Referenz ist immer der eingebundene Quellstand unter
\texttt{vendor/OndselSolver/OndselSolver}. Der Parser akzeptiert nur
eine Teilmenge der im Verzeichnis vorhandenen Klassen. Eine vorhandene
C++-Klasse ist kein Beweis, dass derselbe Typ aus einer ASMT-Datei
geladen oder über \texttt{solve()} verwendet werden kann.

\hypertarget{schichtenmodell}{%
\section{\texorpdfstring{{10.3}
Schichtenmodell}{10.3 Schichtenmodell}}\label{schichtenmodell}}

\begin{longtable}[]{@{}
  >{\raggedright\arraybackslash}p{(\columnwidth - 4\tabcolsep) * \real{0.3333}}
  >{\raggedright\arraybackslash}p{(\columnwidth - 4\tabcolsep) * \real{0.3333}}
  >{\raggedright\arraybackslash}p{(\columnwidth - 4\tabcolsep) * \real{0.3333}}@{}}
\toprule\noalign{}
\begin{minipage}[b]{\linewidth}\raggedright
Ebene
\end{minipage} & \begin{minipage}[b]{\linewidth}\raggedright
Haupttypen
\end{minipage} & \begin{minipage}[b]{\linewidth}\raggedright
Aufgabe
\end{minipage} \\
\midrule\noalign{}
\endhead
\bottomrule\noalign{}
\endlastfoot
Eingabe/Ausgabe & \texttt{ASMTAssembly::assemblyFromFile},
\texttt{outputFile} & ASMT laden und serialisieren \\
Solvermodell & \texttt{ASMTAssembly}, \texttt{ASMTPart},
\texttt{ASMTMarker}, \texttt{ASMTJoint} & Körper, Lage, Referenzen und
Beziehungen darstellen \\
Analysevorgaben & \texttt{ASMTMotion}, \texttt{ASMTLimit},
\texttt{ASMTForceTorque}, \texttt{ASMTConstantGravity},
\texttt{ASMTSimulationParameters} & Bewegung und Analyseparameter
festlegen \\
Übersetzung & \texttt{createMbD}, \texttt{deleteMbD},
\texttt{updateFromMbD} & ASMT-Objekte in interne MbD-Objekte abbilden
und Ergebnisse zurückkopieren \\
Numerischer Kern & \texttt{System}, \texttt{Constraint},
Newton-Raphson-, Integrator- und Matrixklassen & Gleichungen aufstellen,
lösen und integrieren \\
\end{longtable}

Für gewöhnliche Programme ist die oberste ASMT-Schicht der sinnvollste
Einstieg. Die Übersetzungs- und Numerikmethoden sind hauptsächlich für
Solverentwicklung und gezielte Erweiterungen relevant.

\hypertarget{linken-und-namensraum}{%
\section{\texorpdfstring{{10.4} Linken und
Namensraum}{10.4 Linken und Namensraum}}\label{linken-und-namensraum}}

Das echte Buchbeispiel bindet \texttt{ASMTAssembly.h} ein und linkt
gegen die mit CMake gebaute Bibliothek \texttt{OndselSolver}:

\begin{Shaded}
\begin{Highlighting}[]
\PreprocessorTok{\#include }\ImportTok{"ASMTAssembly.h"}

\PreprocessorTok{\#include }\ImportTok{\textless{}iostream\textgreater{}}
\PreprocessorTok{\#include }\ImportTok{\textless{}stdexcept\textgreater{}}

\DataTypeTok{int}\NormalTok{ main}\OperatorTok{(}\DataTypeTok{int}\NormalTok{ argc}\OperatorTok{,} \DataTypeTok{char}\OperatorTok{*}\NormalTok{ argv}\OperatorTok{[])} \OperatorTok{\{}
    \AttributeTok{const} \BuiltInTok{std::}\NormalTok{string input }\OperatorTok{=}\NormalTok{ argc }\OperatorTok{\textgreater{}} \DecValTok{1} \OperatorTok{?}\NormalTok{ argv}\OperatorTok{[}\DecValTok{1}\OperatorTok{]} \OperatorTok{:} \StringTok{"\_\_cubes.asmt"}\OperatorTok{;}
    \AttributeTok{const} \BuiltInTok{std::}\NormalTok{string output }\OperatorTok{=}\NormalTok{ argc }\OperatorTok{\textgreater{}} \DecValTok{2} \OperatorTok{?}\NormalTok{ argv}\OperatorTok{[}\DecValTok{2}\OperatorTok{]} \OperatorTok{:} \StringTok{"\_\_cubes{-}solved.asmt"}\OperatorTok{;}

    \ControlFlowTok{try} \OperatorTok{\{}
        \KeywordTok{auto}\NormalTok{ assembly }\OperatorTok{=}\NormalTok{ MbD}\OperatorTok{::}\NormalTok{ASMTAssembly}\OperatorTok{::}\NormalTok{assemblyFromFile}\OperatorTok{(}\NormalTok{input}\OperatorTok{);}
\NormalTok{        assembly}\OperatorTok{{-}\textgreater{}}\NormalTok{solve}\OperatorTok{();}
\NormalTok{        assembly}\OperatorTok{{-}\textgreater{}}\NormalTok{outputFile}\OperatorTok{(}\NormalTok{output}\OperatorTok{);}
        \BuiltInTok{std::}\NormalTok{cout }\OperatorTok{\textless{}\textless{}} \StringTok{"Ergebnis geschrieben: "} \OperatorTok{\textless{}\textless{}}\NormalTok{ output }\OperatorTok{\textless{}\textless{}} \CharTok{\textquotesingle{}}\SpecialCharTok{\textbackslash{}n}\CharTok{\textquotesingle{}}\OperatorTok{;}
    \OperatorTok{\}} \ControlFlowTok{catch} \OperatorTok{(}\AttributeTok{const} \BuiltInTok{std::}\NormalTok{exception}\OperatorTok{\&}\NormalTok{ error}\OperatorTok{)} \OperatorTok{\{}
        \BuiltInTok{std::}\NormalTok{cerr }\OperatorTok{\textless{}\textless{}}\NormalTok{ error}\OperatorTok{.}\NormalTok{what}\OperatorTok{()} \OperatorTok{\textless{}\textless{}} \CharTok{\textquotesingle{}}\SpecialCharTok{\textbackslash{}n}\CharTok{\textquotesingle{}}\OperatorTok{;}
        \ControlFlowTok{return} \DecValTok{1}\OperatorTok{;}
    \OperatorTok{\}}
\OperatorTok{\}}
\end{Highlighting}
\end{Shaded}

\texttt{MbD} ist der C++-Namensraum. \texttt{assemblyFromFile()} liefert
einen \texttt{std::shared\_ptr\textless{}ASMTAssembly\textgreater{}};
das Beispiel hält die Assembly deshalb mit \texttt{auto} und greift über
\texttt{-\textgreater{}} auf Methoden zu. Der genaue, getestete Build
einschließlich Fixture-Pfad ist
\href{./examples/07-ondselsolver-real-asmt/README.md}{\texttt{examples/07-ondselsolver-real-asmt}}.

\hypertarget{erzeugen-oder-laden}{%
\section{\texorpdfstring{{10.5} Erzeugen oder
Laden}{10.5 Erzeugen oder Laden}}\label{erzeugen-oder-laden}}

Es gibt zwei typische Wege:

\begin{enumerate}
\def\labelenumi{\arabic{enumi}.}
\tightlist
\item
  \textbf{ASMT laden:} \texttt{ASMTAssembly::assemblyFromFile(filename)}
  ist der kürzeste und im Buchbeispiel geprüfte Weg. Der Parser erzeugt
  die Objekte aus den unterstützten ASMT-Blöcken.
\item
  \textbf{Objekte in C++ erzeugen:} ASMT-Klassen bieten häufig eine
  statische Fabrik \texttt{With()}. Anschließend werden ihre Daten
  gesetzt und Objekte über Assembly-Methoden hinzugefügt. Diese Route
  verlangt genaue Kenntnis der konkreten Klassenfelder und
  Markerreferenzen; nicht alle Datenklassen bieten komfortable Setter.
\end{enumerate}

\texttt{ASMTAssembly::With()} erzeugt eine leere Root-Assembly.
\texttt{setName()} stammt von \texttt{ASMTItem}; Lage und Orientierung
werden über \texttt{ASMTSpatialItem} gesetzt. Beim direkten Aufbau muss
die Modellierung vollständig sein, bevor \texttt{solve()} aufgerufen
wird: insbesondere Markerpfade, Joint-Seiten und Anfangslagen.

\hypertarget{lebenszyklus-als-faustregel}{%
\section{\texorpdfstring{{10.6} Lebenszyklus als
Faustregel}{10.6 Lebenszyklus als Faustregel}}\label{lebenszyklus-als-faustregel}}

\begin{Shaded}
\begin{Highlighting}[]
\NormalTok{ASMT{-}Datei oder C++{-}Objekte}
\NormalTok{    {-}\textgreater{} ASMTAssembly}
\NormalTok{    {-}\textgreater{} solve()}
\NormalTok{    {-}\textgreater{} updateFromMbD() innerhalb der Rückgabe{-}/Ausgabepipeline}
\NormalTok{    {-}\textgreater{} outputFile()}
\end{Highlighting}
\end{Shaded}

Die öffentlichen Arbeitsaufrufe für den typischen Batch-Fall sind
\texttt{assemblyFromFile()}, \texttt{solve()} und \texttt{outputFile()}.
\texttt{parseASMT()}, \texttt{createMbD()}, \texttt{updateFromMbD()} und
\texttt{storeOnLevel()} gehören zur Objektübersetzung beziehungsweise
Serialisierung. Sie sind nützlich, wenn man die Schnittstelle erweitert,
aber normalerweise keine zusätzlichen Schritte im Clientprogramm.

\hypertarget{was-nicht-versprochen-wird}{%
\section{\texorpdfstring{{10.7} Was nicht versprochen
wird}{10.7 Was nicht versprochen wird}}\label{was-nicht-versprochen-wird}}

OndselSolver ist in diesem Repository eine C++-Bibliothek mit vielen
öffentlich deklarierten Headermethoden; daraus folgt nicht automatisch
eine versionierte ABI oder ein stabiler Vertrag für jede Methode. Der
Parser nutzt im geprüften Stand exakte Abschnittsreihenfolgen und
teilweise Assertions. Fehlerhafte ASMT-Eingaben können daher anders
scheitern als mit einer umfassenden Schema-Validierung. Die konkrete
Unterstützung ist je Typ in den Folgekapiteln an Parser und Fixtures
gebunden.

\hypertarget{kap-assembly-koerper}{%
\chapter{\texorpdfstring{{11} Assembly, Körper und räumliche
Referenzen}{11 Assembly, Körper und räumliche Referenzen}}\label{kap-assembly-koerper}}

\hypertarget{asmtassembly-die-modellwurzel}{%
\section{\texorpdfstring{{11.1} \texttt{ASMTAssembly}: die
Modellwurzel}{11.1 ASMTAssembly: die Modellwurzel}}\label{asmtassembly-die-modellwurzel}}

\texttt{ASMTAssembly} erbt von \texttt{ASMTSpatialContainer}. Sie ist
daher gleichzeitig Root-Assembly, räumliches Objekt und Container für
Unterobjekte. Die wichtigsten aufrufbaren Methoden sind:

\begin{longtable}[]{@{}
  >{\raggedright\arraybackslash}p{(\columnwidth - 2\tabcolsep) * \real{0.5000}}
  >{\raggedright\arraybackslash}p{(\columnwidth - 2\tabcolsep) * \real{0.5000}}@{}}
\toprule\noalign{}
\begin{minipage}[b]{\linewidth}\raggedright
Funktion
\end{minipage} & \begin{minipage}[b]{\linewidth}\raggedright
Verwendung
\end{minipage} \\
\midrule\noalign{}
\endhead
\bottomrule\noalign{}
\endlastfoot
\texttt{With()} & leere Assembly erzeugen \\
\texttt{assemblyFromFile(filename)} & ASMT-Datei parsen und Assembly
zurückgeben \\
\texttt{solve()} & Assembly einmal lösen; Details in
\protect\hyperlink{kap-solve-ergebnisse}{Kapitel 14} \\
\texttt{outputFile(filename)} & Assembly samt verfügbarem Ergebnis
serialisieren \\
\texttt{addPart(part)}, \texttt{addJoint(joint)},
\texttt{addMotion(motion)}, \texttt{addLimit(limit)} & Objekte zum
Modell hinzufügen \\
\texttt{setConstantGravity(gravity)},
\texttt{setSimulationParameters(parameters)} & globale Analysevorgaben
setzen \\
\texttt{partNamed(name)}, \texttt{partPartialNamed(name)} & Part nach
vollständigem oder partiellem Namen suchen \\
\texttt{partAt(path)}, \texttt{markerAt(path)}, \texttt{jointAt(path)},
\texttt{motionAt(path)}, \texttt{forceTorqueAt(path)} & Objekt anhand
eines ASMT-Pfads auflösen \\
\texttt{markerMap()}, \texttt{connectorList()} & Marker- beziehungsweise
Connector-Übersichten abrufen \\
\texttt{numberOfFrames()}, \texttt{updateForFrame(index)} & verfügbare
Ergebnisframes abfragen und einen Frame aktivieren \\
\texttt{setNotes(text)}, \texttt{setFilename(text)},
\texttt{setDebug(bool)} & Metadaten und Diagnoseeinstellung setzen \\
\end{longtable}

\texttt{parts}, \texttt{joints}, \texttt{motions}, \texttt{limits} und
\texttt{forcesTorques} sind im Header als \texttt{shared\_ptr} auf
Container sichtbar. Das ist eine konkrete Implementierungseigenschaft.
Für neue Programme sind die Methoden \texttt{addPart()} und ihre
Gegenstücke vorzuziehen, weil sie die Owner-Verknüpfung pflegen.

\hypertarget{beispiel-vorhandene-assembly-inspizieren}{%
\subsection*{Beispiel: vorhandene Assembly
inspizieren}\label{beispiel-vorhandene-assembly-inspizieren}}
\addcontentsline{toc}{subsection}{Beispiel: vorhandene Assembly
inspizieren}

\begin{Shaded}
\begin{Highlighting}[]
\KeywordTok{auto}\NormalTok{ assembly }\OperatorTok{=}\NormalTok{ MbD}\OperatorTok{::}\NormalTok{ASMTAssembly}\OperatorTok{::}\NormalTok{assemblyFromFile}\OperatorTok{(}\StringTok{"fourbar.asmt"}\OperatorTok{);}
\BuiltInTok{std::}\NormalTok{cout }\OperatorTok{\textless{}\textless{}} \StringTok{"Parts: "} \OperatorTok{\textless{}\textless{}}\NormalTok{ assembly}\OperatorTok{{-}\textgreater{}}\NormalTok{parts}\OperatorTok{{-}\textgreater{}}\NormalTok{size}\OperatorTok{()} \OperatorTok{\textless{}\textless{}} \CharTok{\textquotesingle{}}\SpecialCharTok{\textbackslash{}n}\CharTok{\textquotesingle{}}\OperatorTok{;}
\ControlFlowTok{for} \OperatorTok{(}\AttributeTok{const} \KeywordTok{auto}\OperatorTok{\&}\NormalTok{ part }\OperatorTok{:} \OperatorTok{*}\NormalTok{assembly}\OperatorTok{{-}\textgreater{}}\NormalTok{parts}\OperatorTok{)} \OperatorTok{\{}
    \BuiltInTok{std::}\NormalTok{cout }\OperatorTok{\textless{}\textless{}} \StringTok{"Part: "} \OperatorTok{\textless{}\textless{}}\NormalTok{ part}\OperatorTok{{-}\textgreater{}}\NormalTok{name }\OperatorTok{\textless{}\textless{}} \CharTok{\textquotesingle{}}\SpecialCharTok{\textbackslash{}n}\CharTok{\textquotesingle{}}\OperatorTok{;}
\OperatorTok{\}}
\end{Highlighting}
\end{Shaded}

Die Namen sind fixtureabhängig. Achtung: \texttt{partNamed()} und
\texttt{partPartialNamed()} behandeln einen nicht gefundenen Namen im
aktuellen Implementierungsstand nicht sicher, sondern dereferenzieren
den Suchiterator. Verwende diese Methoden nur mit verifizierten Namen.
\texttt{partAt(path)} liefert dagegen bei Nichttreffer \texttt{nullptr}.

Ausführbar mit \texttt{./run.sh\ inspect} im Ordner
\href{./examples/08-ondselsolver-api/README.md}{\texttt{examples/08-ondselsolver-api}}.

\hypertarget{basisklasse-asmtitem}{%
\section{\texorpdfstring{{11.2} Basisklasse
\texttt{ASMTItem}}{11.2 Basisklasse ASMTItem}}\label{basisklasse-asmtitem}}

Die ASMT-Objekte teilen sich eine Basisklasse. Fachlich relevante
gemeinsame Funktionen:

\begin{itemize}
\tightlist
\item
  \texttt{setName()} und \texttt{classname()} verwalten beziehungsweise
  liefern Identität/Typ.
\item
  \texttt{root()}, \texttt{partOrAssembly()} und \texttt{part()}
  navigieren in der Objektstruktur.
\item
  \texttt{fullName(partialName)} bildet einen vollständigen Namen/Pfad.
\item
  \texttt{initialize()} richtet ein Objekt ein.
\item
  \texttt{createMbD(system,\ units)} und \texttt{deleteMbD()} erzeugen
  beziehungsweise entfernen die interne MbD-Repräsentation.
\item
  \texttt{updateFromMbD()}, \texttt{updateFromInputState()},
  \texttt{updateFromInitiallyAssembledState()} und
  \texttt{updateForFrame(index)} übertragen Zustände zwischen Eingabe,
  Solver und Ergebnisframe.
\item
  \texttt{compareResults(type)} und \texttt{outputResults(type)}
  vergleichen oder schreiben Analyseergebnisse.
\end{itemize}

Die \texttt{read*}- und \texttt{storeOnLevel*}-Familien sind
Parser-/Serializer-Hilfen. Dazu zählen etwa \texttt{readDouble},
\texttt{readBool}, \texttt{readColumnOfDoubles} sowie das Schreiben von
Namen, Zahlen, Arrays und Tabs. Sie sind für neue Anwendungslogik
normalerweise nicht erforderlich.

\hypertarget{position-und-orientierung-asmtspatialitem}{%
\section{\texorpdfstring{{11.3} Position und Orientierung:
\texttt{ASMTSpatialItem}}{11.3 Position und Orientierung: ASMTSpatialItem}}\label{position-und-orientierung-asmtspatialitem}}

\texttt{ASMTSpatialItem} fügt Lage und Orientierung hinzu. Es wird von
Assembly, Part, Marker und Referenzobjekten verwendet.

\begin{itemize}
\tightlist
\item
  \texttt{setPosition3D(x,\ y,\ z)} und \texttt{setPosition3D(vector)}
  setzen Translation.
\item
  \texttt{getPosition3D(x,\ y,\ z)} liest die Lage in skalare
  Referenzen; \texttt{getPosition3D(frameIndex)} liefert den Zustand
  eines Ergebnisframes.
\item
  \texttt{setRotationMatrix(9\ Werte)} beziehungsweise die
  Matrix-Überladung setzen Orientierung.
\item
  \texttt{getRotationMatrix(frameIndex)} fragt die Matrix eines
  Ergebnisframes ab.
\item
  \texttt{setQuarternions(q0,\ q1,\ q2,\ q3)} setzt eine Quaternion. Der
  im Header vorhandene Schreibfehler \texttt{Quarternions} gehört
  tatsächlich zum Methodennamen.
\item
  \texttt{getQuarternions(...)} liest sie aus.
\item
  \texttt{restorePosRot()} stellt gespeicherte alte Position und
  Rotation wieder her.
\end{itemize}

Positionen und Maße müssen zu den Einheitenskalen des Modells passen.
Rotationsmatrizen sind mathematisch als gültige orthonormale Matrizen zu
verstehen; der Parser validiert nicht jede Eingabe robust.

\hypertarget{asmtspatialcontainer-und-asmtpart}{%
\section{\texorpdfstring{{11.4} \texttt{ASMTSpatialContainer} und
\texttt{ASMTPart}}{11.4 ASMTSpatialContainer und ASMTPart}}\label{asmtspatialcontainer-und-asmtpart}}

Ein räumlicher Container besitzt Marker und Referenzen:

\begin{itemize}
\tightlist
\item
  \texttt{addMarker(marker)}, \texttt{markerList()} und
  \texttt{generateUniqueMarkerName()} verwalten Marker.
\item
  \texttt{addRefPoint(point)} ergänzt einen Referenzpunkt.
\item
  \texttt{setPrincipalMassMarker(marker)} verknüpft das
  Masse-/Trägheitsbezugssystem.
\item
  \texttt{setVelocity3D(...)} und \texttt{setOmega3D(...)} setzen
  lineare beziehungsweise Winkelgeschwindigkeit.
\item
  \texttt{getVelocity3D(i)}, \texttt{getOmega3D(i)},
  \texttt{getAcceleration3D(i)} und \texttt{getAlpha3D(i)} lesen
  Zeitreihenwerte für einen Frame.
\item
  \texttt{rOcmO()} und \texttt{qEp()} liefern abgeleitete
  Orts-/Orientierungsdaten für die Solverabbildung.
\end{itemize}

\texttt{ASMTPart::With()} erzeugt einen Körper. \texttt{ASMTPart}
ergänzt die Containerfunktion um Part-Zeitreihen, \texttt{isFixed},
Parser für \texttt{FeatureOrder}/\texttt{PrincipalMassMarker} und die
MbD-Erzeugung. Die Masseparameter gehören zum
\texttt{ASMTPrincipalMassMarker}; dessen Setter und Einheiten müssen
passend gewählt werden. ASMT-Partgeometrie ist nicht automatisch eine
FreeCAD-B-Rep.

\hypertarget{marker-und-referenzobjekte}{%
\section{\texorpdfstring{{11.5} Marker und
Referenzobjekte}{11.5 Marker und Referenzobjekte}}\label{marker-und-referenzobjekte}}

\texttt{ASMTMarker::With()} erzeugt einen Marker, der von
\texttt{ASMTSpatialItem} erbt und damit eine lokale Position plus
Orientierung besitzt. Marker liegen typischerweise im Koordinatensystem
einer Assembly oder eines Parts. Gelenke verbinden die Marker über die
geerbten \texttt{ASMTItemIJ}-Funktionen \texttt{setMarkerI(path)} und
\texttt{setMarkerJ(path)}.

\texttt{ASMTRefPoint}, \texttt{ASMTRefCurve} und \texttt{ASMTRefSurface}
bilden geometrische Referenzen; \texttt{ASMTRefItem::addMarker()} hängt
Marker an ein Referenzobjekt. Wichtig für ASMT-Dateien: Im geprüften
Parserstand werden Referenzpunkte eingelesen, aber die Elementparser für
Referenzkurven und -flächen sind noch nicht implementiert. Ein C++-Typ
im Header bedeutet also nicht, dass der Dateipfad fertig ist.

\hypertarget{beispiel-ruxe4umliche-daten-setzen}{%
\subsection*{Beispiel: räumliche Daten
setzen}\label{beispiel-ruxe4umliche-daten-setzen}}
\addcontentsline{toc}{subsection}{Beispiel: räumliche Daten setzen}

\begin{Shaded}
\begin{Highlighting}[]
\KeywordTok{auto}\NormalTok{ part }\OperatorTok{=}\NormalTok{ MbD}\OperatorTok{::}\NormalTok{ASMTPart}\OperatorTok{::}\NormalTok{With}\OperatorTok{();}
\NormalTok{part}\OperatorTok{{-}\textgreater{}}\NormalTok{setName}\OperatorTok{(}\StringTok{"Rotor"}\OperatorTok{);}
\NormalTok{part}\OperatorTok{{-}\textgreater{}}\NormalTok{setPosition3D}\OperatorTok{(}\FloatTok{0.0}\OperatorTok{,} \FloatTok{0.0}\OperatorTok{,} \FloatTok{0.25}\OperatorTok{);}
\NormalTok{part}\OperatorTok{{-}\textgreater{}}\NormalTok{setRotationMatrix}\OperatorTok{(}
    \FloatTok{1.0}\OperatorTok{,} \FloatTok{0.0}\OperatorTok{,} \FloatTok{0.0}\OperatorTok{,}
    \FloatTok{0.0}\OperatorTok{,} \FloatTok{1.0}\OperatorTok{,} \FloatTok{0.0}\OperatorTok{,}
    \FloatTok{0.0}\OperatorTok{,} \FloatTok{0.0}\OperatorTok{,} \FloatTok{1.0}\OperatorTok{);}
\end{Highlighting}
\end{Shaded}

Das erzeugt nur ein Partobjekt mit Lage; es erzeugt noch keine
vollständige Assembly. Für einen lösbaren Mechanismus müssen zusätzlich
Massemarker, Joint-Marker, Constraints und korrekte Owner-Beziehungen
angelegt werden. Zum Einstieg ist deshalb das Laden eines bekannten
ASMT-Fixtures verlässlicher.

Der Part-Aufbau samt \texttt{addPart()} und Owner-Prüfung ist als
\texttt{./run.sh\ part-api} im Ordner
\href{./examples/08-ondselsolver-api/README.md}{\texttt{examples/08-ondselsolver-api}}
ausführbar.

\hypertarget{kap-joints}{%
\chapter{\texorpdfstring{{12} Joints und kinematische
Bindungen}{12 Joints und kinematische Bindungen}}\label{kap-joints}}

\hypertarget{verbindungsmodell}{%
\section{\texorpdfstring{{12.1}
Verbindungsmodell}{12.1 Verbindungsmodell}}\label{verbindungsmodell}}

Joints drücken Beziehungen zwischen zwei Körpern beziehungsweise zwei
Markern aus. Der ASMT-Basistyp \texttt{ASMTItemIJ} bietet
\texttt{setMarkerI(path)} und \texttt{setMarkerJ(path)}. Ein Joint
besteht damit mindestens aus Typ, Name, zwei Markerreferenzen und
gegebenenfalls typabhängigen Parametern.

Die Solverbindung entsteht nicht allein durch das Erzeugen der
C++-Instanz. Die Referenzpfade müssen in der Assembly auf vorhandene
Marker zeigen, und der Joint muss in die Assembly aufgenommen werden,
zum Beispiel durch \texttt{assembly-\textgreater{}addJoint(joint)}.

\hypertarget{im-asmt-parser-erkannte-jointtypen}{%
\section{\texorpdfstring{{12.2} Im ASMT-Parser erkannte
Jointtypen}{12.2 Im ASMT-Parser erkannte Jointtypen}}\label{im-asmt-parser-erkannte-jointtypen}}

Der folgende Katalog entspricht den Typzweigen in
\texttt{ASMTAssembly::readJoints()} des hier eingebundenen Solverstands.
Er bezeichnet Parsererkennung, nicht eine Garantie, dass jede denkbare
Parameterkombination numerisch sinnvoll oder in FreeCAD exportierbar
ist.

\begin{longtable}[]{@{}
  >{\raggedright\arraybackslash}p{(\columnwidth - 4\tabcolsep) * \real{0.3333}}
  >{\raggedright\arraybackslash}p{(\columnwidth - 4\tabcolsep) * \real{0.3333}}
  >{\raggedright\arraybackslash}p{(\columnwidth - 4\tabcolsep) * \real{0.3333}}@{}}
\toprule\noalign{}
\begin{minipage}[b]{\linewidth}\raggedright
Gruppe
\end{minipage} & \begin{minipage}[b]{\linewidth}\raggedright
ASMT-Typen
\end{minipage} & \begin{minipage}[b]{\linewidth}\raggedright
Bedeutung auf hoher Ebene
\end{minipage} \\
\midrule\noalign{}
\endhead
\bottomrule\noalign{}
\endlastfoot
Fixe/rotatorische Punktverbindungen & \texttt{FixedJoint},
\texttt{RevoluteJoint}, \texttt{SphericalJoint},
\texttt{UniversalJoint}, \texttt{ConstantVelocityJoint} & Punkte
zusammenhalten und rotationsabhängige Freiheitsgrade beschränken \\
Axiale und planare Verbindungen & \texttt{CylindricalJoint},
\texttt{TranslationalJoint}, \texttt{PointInLineJoint},
\texttt{PlanarJoint}, \texttt{PointInPlaneJoint},
\texttt{LineInPlaneJoint} & Linie/Achse oder Ebene als Bezug für
erlaubte Bewegung nutzen \\
Orientierungsvorgaben & \texttt{NoRotationJoint},
\texttt{ParallelAxesJoint}, \texttt{PerpendicularJoint},
\texttt{AngleJoint} & relative Orientierung bzw. Winkelbeziehung
festlegen \\
Übersetzungsbeziehungen & \texttt{GearJoint}, \texttt{RackPinionJoint},
\texttt{ScrewJoint} & gekoppelte Rotation/Translation über Verhältnis,
Radius oder Steigung vorgeben \\
Zusammengesetzte Gelenke & \texttt{SphSphJoint}, \texttt{CylSphJoint},
\texttt{RevCylJoint}, \texttt{RevRevJoint} & Kombinationen elementarer
Verbindungstypen modellieren \\
\end{longtable}

Die interne Verteilung auf Basisklassen folgt der Mechanik:
\texttt{ASMTAtPointJoint}, \texttt{ASMTInLineJoint},
\texttt{ASMTInPlaneJoint} und \texttt{ASMTCompoundJoint} bündeln
gemeinsame Erzeugungslogik. Konkrete Klassen wie
\texttt{ASMTRevoluteJoint} oder \texttt{ASMTGearJoint} ergänzen ihre
Parameter über eigene Felder und \texttt{createMbD()}-Implementierungen.

\hypertarget{das-muster-fuxfcr-den-c-aufbau}{%
\section{\texorpdfstring{{12.3} Das Muster für den
C++-Aufbau}{12.3 Das Muster für den C++-Aufbau}}\label{das-muster-fuxfcr-den-c-aufbau}}

\begin{Shaded}
\begin{Highlighting}[]
\KeywordTok{auto}\NormalTok{ joint }\OperatorTok{=}\NormalTok{ MbD}\OperatorTok{::}\NormalTok{ASMTRevoluteJoint}\OperatorTok{::}\NormalTok{With}\OperatorTok{();}
\NormalTok{joint}\OperatorTok{{-}\textgreater{}}\NormalTok{setName}\OperatorTok{(}\StringTok{"Hinge"}\OperatorTok{);}
\NormalTok{joint}\OperatorTok{{-}\textgreater{}}\NormalTok{setMarkerI}\OperatorTok{(}\StringTok{"/Assembly/Base/HingeMarker"}\OperatorTok{);}
\NormalTok{joint}\OperatorTok{{-}\textgreater{}}\NormalTok{setMarkerJ}\OperatorTok{(}\StringTok{"/Assembly/Link/HingeMarker"}\OperatorTok{);}
\NormalTok{assembly}\OperatorTok{{-}\textgreater{}}\NormalTok{addJoint}\OperatorTok{(}\NormalTok{joint}\OperatorTok{);}
\end{Highlighting}
\end{Shaded}

Das Beispiel zeigt die gemeinsame API, nicht ein vollständiges
lauffähiges Mechanikmodell. Namen und Pfade müssen den tatsächlichen
Container- und Marker-Namen entsprechen. Zusätzlich müssen beide Marker
existieren und auf den richtigen lokalen Koordinatensystemen liegen. Für
Sondertypen wie \texttt{AngleJoint}, \texttt{GearJoint} und
\texttt{ScrewJoint} kommen typspezifische Parameter hinzu; ihre
Feldnamen sind in den konkreten Headern (\texttt{ASMTAngleJoint.h} usw.)
definiert und dürfen nicht aus dem Jointnamen erraten werden.

Der ausführbare Modus \texttt{./run.sh\ joint-api} in
\href{./examples/08-ondselsolver-api/README.md}{\texttt{examples/08-ondselsolver-api}}
prüft Fabrik, Marker-Setter, \texttt{addJoint()} und Owner. Er löst den
absichtlich zusätzlich angelegten Joint nicht.

\hypertarget{gemeinsame-ergebnisfunktionen}{%
\section{\texorpdfstring{{12.4} Gemeinsame
Ergebnisfunktionen}{12.4 Gemeinsame Ergebnisfunktionen}}\label{gemeinsame-ergebnisfunktionen}}

Joints erben die Auswertungsfunktionen von \texttt{ASMTConstraintSet}
und \texttt{ASMTItem}:

\begin{itemize}
\tightlist
\item
  \texttt{updateFromMbD()} übernimmt den berechneten Solverzustand.
\item
  \texttt{compareResults(type)} vergleicht Objekt-/Solverdaten für eine
  Analyseart.
\item
  \texttt{outputResults(type)} gibt Ergebnisse aus.
\item
  \texttt{createMbD(system,\ units)} erzeugt die interne
  Jointrepräsentation.
\item
  \texttt{storeOnLevel(...)} serialisiert den Joint;
  \texttt{storeOnTimeSeries(...)} schreibt Joint-Zeitreihen.
\end{itemize}

\texttt{ASMTJoint::readJointSeries()} lädt Zeitreihenwerte. Dies ist
keine Voraussetzung für die Definition eines Eingabe-Joints.

\hypertarget{kontakte-klassen-vorhanden-pfad-pruxfcfen}{%
\section{\texorpdfstring{{12.5} Kontakte: Klassen vorhanden, Pfad
prüfen}{12.5 Kontakte: Klassen vorhanden, Pfad prüfen}}\label{kontakte-klassen-vorhanden-pfad-pruxfcfen}}

Der Quellbaum enthält \texttt{ASMTContact},
\texttt{ASMTConeConeContact}, \texttt{ASMTCylConeContact} und
\texttt{ASMTCylCylContact}. Diese Klassen sind als Modellfamilie
erkennbar, werden aber im gezeigten \texttt{readJoints()}-Dispatch nicht
als Joint-Abschnitte angelegt. Sie dürfen daher nicht mit den oben
bestätigten ASMT-Jointtypen gleichgesetzt werden. Vor Verwendung sind
Parser, \texttt{createMbD()} und Tests im konkreten Commit zu prüfen.

\hypertarget{freiheitsgrade-und-uxfcberbestimmung}{%
\section{\texorpdfstring{{12.6} Freiheitsgrade und
Überbestimmung}{12.6 Freiheitsgrade und Überbestimmung}}\label{freiheitsgrade-und-uxfcberbestimmung}}

Ein Joint fügt Gleichungsbedingungen hinzu und reduziert die
Beweglichkeit. Wird ein Mechanismus überbestimmt oder geometrisch
widersprüchlich, kann die Lösung scheitern. Der richtige Workflow ist:

\begin{enumerate}
\def\labelenumi{\arabic{enumi}.}
\tightlist
\item
  Markerpositionen und lokale Achsen kontrollieren.
\item
  Joints zunächst schrittweise ergänzen.
\item
  Nach jeder Änderung mit einem kleinen bekannten Modell lösen.
\item
  Erst danach Bewegungen und Lasten ergänzen.
\end{enumerate}

Der numerische Solver kann keine inkonsistente Modellabsicht erraten;
eine Fehlermeldung ist deshalb häufig ein Hinweis auf Geometrie,
redundante Constraints oder ungünstige Anfangswerte.

\hypertarget{kap-bewegungen-kraefte}{%
\chapter{\texorpdfstring{{13} Bewegungen, Kräfte und
Analyseparameter}{13 Bewegungen, Kräfte und Analyseparameter}}\label{kap-bewegungen-kraefte}}

\hypertarget{motion-als-vorgeschriebene-bewegung}{%
\section{\texorpdfstring{{13.1} Motion als vorgeschriebene
Bewegung}{13.1 Motion als vorgeschriebene Bewegung}}\label{motion-als-vorgeschriebene-bewegung}}

\texttt{ASMTMotion} ist ein kinematischer ConstraintSet-Typ. Die
konkreten Klassen werden im Parser unter \texttt{ConstraintSets/Motions}
erkannt:

\begin{itemize}
\tightlist
\item
  \texttt{ASMTRotationalMotion}: vorgeschriebene Rotation
\item
  \texttt{ASMTTranslationalMotion}: vorgeschriebene Translation
\item
  \texttt{ASMTGeneralMotion}: allgemeinere Motion-Darstellung
\item
  \texttt{ASMTAllowRotation}: Rotationsfreiheit beziehungsweise
  entsprechende Erlaubnisbedingung
\end{itemize}

Die gemeinsamen Funktionen sind \texttt{initMarkers()},
\texttt{createMbD(system,\ units)}, \texttt{readMotionSeries()},
\texttt{storeOnLevel()} und \texttt{storeOnTimeSeries()}. Für eine
Motion müssen die Marker beziehungsweise Joint-Bezüge und die
Bewegungsgesetze der konkreten Unterklasse konsistent sein. Die konkrete
Parameterbenennung ist klassenabhängig; der jeweilige
\texttt{ASMT*.h}-Header und ein echtes Fixture sind die Referenz.

Ein Kinematikmodell benötigt nicht immer externe Kräfte. Eine
vorgeschriebene Bewegung ist eine Randbedingung und unterscheidet sich
von einem Drehmoment, das eine dynamische Reaktion auslöst.

\hypertarget{limits}{%
\section{\texorpdfstring{{13.2} Limits}{13.2 Limits}}\label{limits}}

\texttt{ASMTLimit} referenziert über \texttt{setmotionJoint(path)} die
Motion oder das Joint, auf das die Grenze wirkt. Der Parser erkennt:

\begin{itemize}
\tightlist
\item
  \texttt{ASMTRotationLimit}
\item
  \texttt{ASMTTranslationLimit}
\end{itemize}

Limits werden über \texttt{assembly-\textgreater{}addLimit(limit)}
eingehängt. Grenzen sind nicht dasselbe wie ein zusätzliches Gelenk: Sie
begrenzen eine Koordinate oder Bewegung und können je nach Zustand
aktiv/inaktiv werden. Welche Wertebereiche und Verhaltenseigenschaften
gelten, stehen in \texttt{ASMTRotationLimit.h} beziehungsweise
\texttt{ASMTTranslationLimit.h} und der MbD-Implementierung.

\hypertarget{kraft-drehmoment-und-schwerkraft}{%
\section{\texorpdfstring{{13.3} Kraft, Drehmoment und
Schwerkraft}{13.3 Kraft, Drehmoment und Schwerkraft}}\label{kraft-drehmoment-und-schwerkraft}}

\texttt{ASMTForceTorque} beschreibt eine Belastung zwischen I/J-Bezügen.
Die Assembly hält solche Objekte im Container \texttt{forcesTorques};
anders als \texttt{addPart()} oder \texttt{addJoint()} ist im sichtbaren
\texttt{ASMTAssembly}-Header kein \texttt{addForceTorque()}-Setter
deklariert. Beim direkten C++-Aufbau muss daher der aktuelle
Owner-/Container-Aufbau in den Quellen beachtet werden.

\texttt{ASMTConstantGravity} repräsentiert eine konstante
Beschleunigung. Das Objekt wird mit \texttt{ASMTConstantGravity::With()}
erzeugt, über \texttt{setg(x,\ y,\ z)} oder \texttt{setg(vector)}
konfiguriert und mit
\texttt{assembly-\textgreater{}setConstantGravity(gravity)} der Assembly
zugewiesen.

\texttt{ASMTForceTorque} und die Klassen \texttt{ASMTConeConeContact},
\texttt{ASMTCylConeContact}, \texttt{ASMTCylCylContact} sind nicht mit
den vier bestätigten Motion-Typen oder dem Joint-Dispatch zu
verwechseln. Die konkrete Parser- und Löseunterstützung ist für diese
Typen getrennt zu prüfen.

\hypertarget{simulationseinstellungen}{%
\section{\texorpdfstring{{13.4}
Simulationseinstellungen}{13.4 Simulationseinstellungen}}\label{simulationseinstellungen}}

\texttt{ASMTSimulationParameters::With()} erzeugt die Laufzeitvorgaben.
Dokumentierte Setter sind:

\begin{longtable}[]{@{}
  >{\raggedright\arraybackslash}p{(\columnwidth - 2\tabcolsep) * \real{0.5000}}
  >{\raggedright\arraybackslash}p{(\columnwidth - 2\tabcolsep) * \real{0.5000}}@{}}
\toprule\noalign{}
\begin{minipage}[b]{\linewidth}\raggedright
Setter
\end{minipage} & \begin{minipage}[b]{\linewidth}\raggedright
Bedeutung
\end{minipage} \\
\midrule\noalign{}
\endhead
\bottomrule\noalign{}
\endlastfoot
\texttt{settstart(t)} / \texttt{settend(t)} & Start- und Endzeit \\
\texttt{sethmin(h)} / \texttt{sethmax(h)} & minimale und maximale
Integrationsschrittweite \\
\texttt{sethout(h)} & gewünschter Ausgabeabstand \\
\texttt{seterrorTol(tol)} & allgemeine Fehlertoleranz \\
\texttt{setmaxIter(n)} & Iterationsmaximum \\
\end{longtable}

Weitere Felder im Header steuern unter anderem
Positions-/Beschleunigungs-Kinematiktoleranzen, absolute/relative
Korrektur- und Integrationstoleranzen, Grenzwerte sowie
\texttt{iterMaxPosKine}, \texttt{iterMaxAccKine}, \texttt{iterMaxDyn}
und \texttt{orderMax}. Nicht jedes Feld besitzt einen eigenen Setter; in
diesem Stand sind einige direkt als öffentliche Member vorhanden.

\begin{Shaded}
\begin{Highlighting}[]
\KeywordTok{auto}\NormalTok{ parameters }\OperatorTok{=}\NormalTok{ MbD}\OperatorTok{::}\NormalTok{ASMTSimulationParameters}\OperatorTok{::}\NormalTok{With}\OperatorTok{();}
\NormalTok{parameters}\OperatorTok{{-}\textgreater{}}\NormalTok{settstart}\OperatorTok{(}\FloatTok{0.0}\OperatorTok{);}
\NormalTok{parameters}\OperatorTok{{-}\textgreater{}}\NormalTok{settend}\OperatorTok{(}\FloatTok{1.0}\OperatorTok{);}
\NormalTok{parameters}\OperatorTok{{-}\textgreater{}}\NormalTok{sethmin}\OperatorTok{(}\FloatTok{1.0e{-}8}\OperatorTok{);}
\NormalTok{parameters}\OperatorTok{{-}\textgreater{}}\NormalTok{sethmax}\OperatorTok{(}\FloatTok{0.02}\OperatorTok{);}
\NormalTok{parameters}\OperatorTok{{-}\textgreater{}}\NormalTok{sethout}\OperatorTok{(}\FloatTok{0.01}\OperatorTok{);}
\NormalTok{parameters}\OperatorTok{{-}\textgreater{}}\NormalTok{seterrorTol}\OperatorTok{(}\FloatTok{1.0e{-}6}\OperatorTok{);}
\NormalTok{assembly}\OperatorTok{{-}\textgreater{}}\NormalTok{setSimulationParameters}\OperatorTok{(}\NormalTok{parameters}\OperatorTok{);}
\NormalTok{assembly}\OperatorTok{{-}\textgreater{}}\NormalTok{runKINEMATIC}\OperatorTok{();}
\end{Highlighting}
\end{Shaded}

Ein \texttt{tstart\ ==\ tend} wird in den eingebauten Beispielroutinen
ausdrücklich als Initial-Conditions-only-Fall verwendet. Für eine echte
Zeitintegration ist ein Intervall mit
\texttt{tend\ \textgreater{}\ tstart} nötig. Wichtig:
\texttt{ASMTAssembly::solve()} setzt die Parameter vor dem Lauf selbst
wieder auf seine eingebauten Werte mit \texttt{tstart\ ==\ tend\ ==\ 0}.
Wer eigene Zeitgrenzen konfiguriert, darf daher nicht erwarten, dass ein
anschließendes \texttt{solve()} sie verwendet; der direkte
\texttt{runKINEMATIC()}-Pfad nutzt die gesetzten Einstellungen, hat im
aktuellen Code aber eine eigene Fehlerbehandlungsgrenze (siehe
\protect\hyperlink{kap-solve-ergebnisse}{Kapitel 14}).

Der Modus \texttt{./run.sh\ settings-api} in
\href{./examples/08-ondselsolver-api/README.md}{\texttt{examples/08-ondselsolver-api}}
prüft die Setter und Owner-Verknüpfungen, ohne danach \texttt{solve()}
aufzurufen.

\hypertarget{animation-zeit-und-einheiten}{%
\section{\texorpdfstring{{13.5} Animation, Zeit und
Einheiten}{13.5 Animation, Zeit und Einheiten}}\label{animation-zeit-und-einheiten}}

\texttt{ASMTAnimationParameters} enthält Einstellungen für die
Darstellung/Frames und ist von den physikalischen Integrationsparametern
zu unterscheiden. \texttt{ASMTTime} bindet die Zeitvariable für
Funktionen ein. \texttt{Units} bündelt Zeit-, Masse-, Längen- und
Winkelskalen sowie abgeleitete Faktoren; die Standardfaktoren stehen in
\texttt{Units.h}.

Alle Komponenten eines Modells müssen dieselbe Einheitenskala benutzen.
Insbesondere Winkelwerte in ASMT sind im Solverkontext typischerweise
Radiant; die Fixture \texttt{\_\_cubes.asmt} zeigt für 10 Grad den Wert
\texttt{0.17453292519943295}. Eine stillschweigende Umrechnung von Grad
nach Radiant sollte man nicht voraussetzen.

\hypertarget{programmatisch-unterstuxfctzte-klassenfamilien-im-uxfcberblick}{%
\section{\texorpdfstring{{13.6} Programmatisch unterstützte
Klassenfamilien im
Überblick}{13.6 Programmatisch unterstützte Klassenfamilien im Überblick}}\label{programmatisch-unterstuxfctzte-klassenfamilien-im-uxfcberblick}}

Der ASMT-Parser erkennt die vier Motiontypen und die zwei Limittypen
oben. Er lädt außerdem ForceTorque- und ConstantGravity-Abschnitte.
\texttt{GeneralConstraintSets} und \texttt{KinematicIJs} sind im
geprüften Parserstand nicht vollständig implementiert. Ein
Dateiformatblock, ein Header und ein erfolgreicher Solverpfad sind drei
verschiedene Unterstützungsstufen; in Tests sollten alle drei separat
bestätigt werden.

\hypertarget{kap-solve-ergebnisse}{%
\chapter{\texorpdfstring{{14} Solve, Ergebnisse und
Funktionskatalog}{14 Solve, Ergebnisse und Funktionskatalog}}\label{kap-solve-ergebnisse}}

\hypertarget{einmaliges-luxf6sen-und-analysearten}{%
\section{\texorpdfstring{{14.1} Einmaliges Lösen und
Analysearten}{14.1 Einmaliges Lösen und Analysearten}}\label{einmaliges-luxf6sen-und-analysearten}}

\texttt{ASMTAssembly::solve()} ist laut Header ein One-shot-Solve der
Assembly. Im aktuellen Quellstand setzt die Methode allerdings zuerst
die Simulationsparameter fest auf \texttt{tstart\ =\ tend\ =\ 0},
\texttt{hmin\ =\ 1e-9}, \texttt{hmax\ =\ 1}, \texttt{hout\ =\ 0.04} und
\texttt{errorTol\ =\ 1e-6} und ruft dann \texttt{runKINEMATIC()} auf.
Das ist eine Lösung der Anfangskonfiguration, keine konfigurierbare
Zeitintegration. Eigene zuvor gesetzte Simulationsparameter werden
dadurch überschrieben. Der praktische Aufruf bleibt daher einfach:

\begin{Shaded}
\begin{Highlighting}[]
\KeywordTok{auto}\NormalTok{ assembly }\OperatorTok{=}\NormalTok{ MbD}\OperatorTok{::}\NormalTok{ASMTAssembly}\OperatorTok{::}\NormalTok{assemblyFromFile}\OperatorTok{(}\NormalTok{inputPath}\OperatorTok{);}
\NormalTok{assembly}\OperatorTok{{-}\textgreater{}}\NormalTok{solve}\OperatorTok{();}
\NormalTok{assembly}\OperatorTok{{-}\textgreater{}}\NormalTok{outputFile}\OperatorTok{(}\NormalTok{outputPath}\OperatorTok{);}
\end{Highlighting}
\end{Shaded}

Die Assembly stellt außerdem \texttt{runKINEMATIC()},
\texttt{runPreDrag()}, \texttt{runDragStep(parts)} und
\texttt{runPostDrag()} bereit. \texttt{runKINEMATIC()} verwendet die
aktuell gesetzten Parameter, fängt aber \texttt{SimulationStoppingError}
intern ab und gibt keinen Erfolgsstatus zurück. Das ist eine
Einschränkung der Fehlerbeobachtbarkeit. Dragging-Aufrufe und
\texttt{run*Test()}-Funktionen sind spezialisierte
Interaktions-/Testpfade, nicht der normale Batch-Solve-Vertrag.
\texttt{outputFor(type)}, \texttt{compareResults(type)} und
\texttt{outputResults(type)} sind Analyse-/Diagnosepfade.

\hypertarget{ergebnisse-abfragen}{%
\section{\texorpdfstring{{14.2} Ergebnisse
abfragen}{14.2 Ergebnisse abfragen}}\label{ergebnisse-abfragen}}

Ergebnisserien liegen abhängig vom Typ in \texttt{times},
\texttt{partSeries}, \texttt{jointSeries} oder \texttt{motionSeries}.
\texttt{numberOfFrames()} liefert die Anzahl vorhandener Frames;
\texttt{updateForFrame(index)} aktualisiert räumliche Objekte auf einen
Frame. Räumliche Getter wie \texttt{getPosition3D(i)},
\texttt{getRotationMatrix(i)}, \texttt{getVelocity3D(i)} und
\texttt{getAcceleration3D(i)} lesen die entsprechenden Frame-Daten.

Der ASMT-Schreibpfad \texttt{outputFile()} ist der einfachste
persistente Ergebniszugriff. Diese Ausgabe ist kein automatischer Export
in ein FreeCAD-Dokument und ersetzt keine CAD-Geometrie. Das
Buchbeispiel überprüft den End-to-End-Pfad, indem es die gelöste Datei
erneut als ASMT-Artefakt schreibt.

Das erneute Einlesen dieser Ausgabe ist im aktuellen GCC-Build nicht
verlässlich: \texttt{ASMTJoint::storeOnLevel()} erzeugt aus dem
RTTI-Namen teilweise ungültige Typzeilen. Beispiel 08 prüft daher die
ASMT-Anwendungsfunktionen im Speicher und behauptet keinen erfolgreichen
Serializer-Roundtrip.

\hypertarget{fehlerbehandlung-und-diagnose}{%
\section{\texorpdfstring{{14.3} Fehlerbehandlung und
Diagnose}{14.3 Fehlerbehandlung und Diagnose}}\label{fehlerbehandlung-und-diagnose}}

Das Clientprogramm sollte Parser-, Datei- und Solverfehler abfangen. Die
vorhandenen Exceptions umfassen beispielsweise
\texttt{SimulationStoppingError}, \texttt{NewtonRaphsonError},
\texttt{MaximumIterationError}, \texttt{SingularMatrixError},
\texttt{TooSmallStepSizeError}, \texttt{TooManyTriesError},
\texttt{DiscontinuityError} und \texttt{NotKinematicError}. Viele leiten
sich von Standard-Exceptions ab;
\texttt{catch\ (const\ std::exception\&)} bildet daher eine brauchbare
äußere Fehlergrenze.

Bei einem Solve-Fehler zuerst prüfen:

\begin{itemize}
\tightlist
\item
  Ist die ASMT-Einrückung und Blockreihenfolge gültig?
\item
  Existieren alle Markerpfade und zeigen sie auf sinnvolle
  Koordinatensysteme?
\item
  Ist die Zahl und Kombination der Constraints konsistent?
\item
  Liegt die Anfangslage nahe einer lösbaren Konfiguration?
\item
  Sind Zeitintervall, Schrittweiten und Einheitenskalen plausibel?
\end{itemize}

\texttt{setDebug(true)} und die vorhandenen
Ergebnis-/Vergleichsfunktionen können bei der Diagnose helfen, sind aber
keine Garantie für eine vollständige Solverdiagnose.

\hypertarget{uxf6ffentliche-asmt-funktionen-nach-klassenfamilie}{%
\section{\texorpdfstring{{14.4} Öffentliche ASMT-Funktionen nach
Klassenfamilie}{14.4 Öffentliche ASMT-Funktionen nach Klassenfamilie}}\label{uxf6ffentliche-asmt-funktionen-nach-klassenfamilie}}

Dieser Katalog beschreibt die fachlich relevanten Methoden des
Anwendungszugangs; alle vererbten Parser-/Serializer-Helfer werden nicht
pro Unterklasse wiederholt.

\begin{longtable}[]{@{}
  >{\raggedright\arraybackslash}p{(\columnwidth - 2\tabcolsep) * \real{0.5000}}
  >{\raggedright\arraybackslash}p{(\columnwidth - 2\tabcolsep) * \real{0.5000}}@{}}
\toprule\noalign{}
\begin{minipage}[b]{\linewidth}\raggedright
Familie
\end{minipage} & \begin{minipage}[b]{\linewidth}\raggedright
Funktionen und Zweck
\end{minipage} \\
\midrule\noalign{}
\endhead
\bottomrule\noalign{}
\endlastfoot
Assembly & \texttt{With}, \texttt{assemblyFromFile}, \texttt{solve},
\texttt{outputFile}; \texttt{addPart}, \texttt{addJoint},
\texttt{addMotion}, \texttt{addLimit}; \texttt{setConstantGravity},
\texttt{setSimulationParameters}; Namens-/Pfadsuche \texttt{partNamed},
\texttt{partPartialNamed} (nur mit garantiertem Treffer),
\texttt{partAt}, \texttt{markerAt}, \texttt{jointAt}, \texttt{motionAt},
\texttt{forceTorqueAt}; Framezugriff \texttt{numberOfFrames},
\texttt{updateForFrame} \\
Gemeinsames Objekt (\texttt{ASMTItem}) & \texttt{setName},
\texttt{classname}, \texttt{root}, \texttt{partOrAssembly},
\texttt{part}, \texttt{fullName}, \texttt{initialize};
Zustands-/MbD-Übergang \texttt{createMbD}, \texttt{deleteMbD},
\texttt{updateFromInputState},
\texttt{updateFromInitiallyAssembledState}, \texttt{updateFromMbD},
\texttt{updateForFrame}; Ausgabe \texttt{compareResults},
\texttt{outputResults} \\
Lage (\texttt{ASMTSpatialItem}) & \texttt{setPosition3D},
\texttt{getPosition3D}, \texttt{setRotationMatrix},
\texttt{getRotationMatrix}, \texttt{setQuarternions},
\texttt{getQuarternions}, \texttt{restorePosRot} \\
Räumlicher Container & \texttt{setPrincipalMassMarker},
\texttt{addRefPoint}, \texttt{addMarker}, \texttt{markerList},
\texttt{generateUniqueMarkerName}; \texttt{setVelocity3D},
\texttt{setOmega3D}, \texttt{getVelocity3D}, \texttt{getOmega3D},
\texttt{getAcceleration3D}, \texttt{getAlpha3D}; \texttt{rOcmO},
\texttt{qEp} \\
Verbindungen (\texttt{ASMTItemIJ}/\texttt{ASMTJoint}) &
\texttt{setMarkerI}, \texttt{setMarkerJ}; geerbte Solve-/Resultmethoden;
\texttt{readJointSeries} und \texttt{storeOnTimeSeries} für
Ergebnisdaten \\
Motion & \texttt{initMarkers}, \texttt{readMotionSeries},
\texttt{storeOnTimeSeries}; konkrete Bewegungsgesetze in den
Unterklassen \\
Limits & \texttt{setmotionJoint}; Typ und Grenzparameter in
\texttt{ASMTRotationLimit}/\texttt{ASMTTranslationLimit} \\
Simulation & \texttt{settstart}, \texttt{settend}, \texttt{sethmin},
\texttt{sethmax}, \texttt{sethout}, \texttt{seterrorTol},
\texttt{setmaxIter}; weitere toleranz-/integratorbezogene Felder in
\texttt{ASMTSimulationParameters} \\
Datei und Text & \texttt{parseASMT}, \texttt{read*},
\texttt{storeOnLevel*}, \texttt{storeOnTimeSeries}, \texttt{logString};
Parser-/Serialisierungsmechanik, kein typischer Clientcode \\
\end{longtable}

Die API ist vererbungsbasiert: Eine konkrete Unterklasse erhält
zahlreiche Funktionen aus \texttt{ASMTItem}, \texttt{ASMTSpatialItem}
oder \texttt{ASMTConstraintSet}. Der Tabellenkatalog fasst diese einmal
pro Familie zusammen, damit dieselbe Funktion nicht dutzendfach als
scheinbar eigene Methode wiederholt wird.

\hypertarget{interner-numerischer-kern-klassifikation}{%
\section{\texorpdfstring{{14.5} Interner numerischer Kern:
Klassifikation}{14.5 Interner numerischer Kern: Klassifikation}}\label{interner-numerischer-kern-klassifikation}}

Der Solverkern bietet wesentlich mehr Methoden als für eine gewöhnliche
ASMT-Anwendung benötigt werden. Für Orientierung lassen sie sich so
einteilen:

\begin{longtable}[]{@{}
  >{\raggedright\arraybackslash}p{(\columnwidth - 4\tabcolsep) * \real{0.3333}}
  >{\raggedright\arraybackslash}p{(\columnwidth - 4\tabcolsep) * \real{0.3333}}
  >{\raggedright\arraybackslash}p{(\columnwidth - 4\tabcolsep) * \real{0.3333}}@{}}
\toprule\noalign{}
\begin{minipage}[b]{\linewidth}\raggedright
Funktionsbereich
\end{minipage} & \begin{minipage}[b]{\linewidth}\raggedright
Beispiele für Klassen
\end{minipage} & \begin{minipage}[b]{\linewidth}\raggedright
Verantwortung
\end{minipage} \\
\midrule\noalign{}
\endhead
\bottomrule\noalign{}
\endlastfoot
System und Zustandsmodell & \texttt{System}, \texttt{StateData},
\texttt{Part}, \texttt{MarkerFrame}, \texttt{CartesianFrame} & globale
Mechanik, Zustandsvariablen und Bezugssysteme verwalten \\
Constraints und Joints & \texttt{Constraint}, \texttt{ConstraintSet},
\texttt{ConstraintIJ}, \texttt{RevoluteJoint}, \texttt{FixedJoint},
Spezial-Constraints & Gleichungen und kinematische Bindungen
aufstellen \\
Position-/Geschwindigkeits-/Beschleunigungslöser &
\texttt{PosNewtonRaphson}, \texttt{VelSolver}, \texttt{AccNewtonRaphson}
und Kinematikvarianten & algebraische Gleichungen auf verschiedenen
Ableitungsstufen lösen \\
Zeitintegration & \texttt{Integrator}, \texttt{KineIntegrator},
\texttt{BasicIntegrator}, \texttt{LinearMultiStepMethod},
\texttt{StableBackwardDifference} & Zustände über die Zeit
integrieren \\
Lineare Algebra & \texttt{FullMatrix}, \texttt{SparseMatrix},
\texttt{FullVector}, \texttt{SparseVector}, \texttt{MatrixSolver},
\texttt{GE*}, \texttt{LDU*} & Gleichungssysteme und Zerlegungen
berechnen \\
Symbolische Ausdrücke/Funktionen & \texttt{Function},
\texttt{ExpressionX}, \texttt{Polynomial}, \texttt{Sine},
\texttt{Cosine}, \texttt{GeneralSpline}, \texttt{PiecewiseFunction} &
zeitabhängige Gesetze und Ableitungen repräsentieren \\
Einheiten und Numerik & \texttt{Units}, \texttt{MbDMath},
\texttt{Numeric}, \texttt{EulerParameters}, \texttt{Orientation} &
Skalen, numerische Hilfen und Rotationsdarstellungen \\
Fehlerzustände & \texttt{NewtonRaphsonError},
\texttt{SingularMatrixError}, \texttt{MaximumIterationError} usw. &
Konvergenz- und Simulationsfehler melden \\
\end{longtable}

Diese Klassen sind keine austauschbaren „Utilities``. Viele Solvertypen
werden intern anhand des Modells ausgewählt und setzen Kenntnisse über
Gleichungsaufbau, Freiheitsgrade und Zustandsverwaltung voraus. Für die
Anwendungsintegration zuerst ASMT nutzen; direktes Arbeiten mit dem
MbD-Kern ist eine bewusste Erweiterung.

\hypertarget{reproduzierbare-validierung}{%
\section{\texorpdfstring{{14.6} Reproduzierbare
Validierung}{14.6 Reproduzierbare Validierung}}\label{reproduzierbare-validierung}}

Das Projekt enthält drei reale Fixtures und den End-to-End-Aufruf unter
\href{./examples/07-ondselsolver-real-asmt/README.md}{\texttt{examples/07-ondselsolver-real-asmt}}.
Die drei Prüfpfade decken unterschiedliche Modelltypen ab:

\begin{itemize}
\tightlist
\item
  \texttt{\_\_cubes.asmt}: zwei Körper mit Fixierung und Winkelbedingung
\item
  \texttt{fourbar.asmt}: Viergelenkmechanismus
\item
  \texttt{piston.asmt}: Kurbel-Schieber
\end{itemize}

Ausführen:

\begin{Shaded}
\begin{Highlighting}[]
\BuiltInTok{cd}\NormalTok{ examples/07{-}ondselsolver{-}real{-}asmt}
\ExtensionTok{./run.sh}
\ExtensionTok{./run.sh}\NormalTok{ fourbar.asmt}
\ExtensionTok{./run.sh}\NormalTok{ piston.asmt}
\end{Highlighting}
\end{Shaded}

Damit wird der echte lokale OndselSolver gebaut/verwendet, die Datei
eingelesen, \texttt{solve()} ausgeführt und ein Ergebnisfile
geschrieben. Wenn der API-Code geändert oder neue ASMT-Typen aufgenommen
werden, sollten Parser-, Solve- und Ausgabeweg jeweils mit einem
gezielten Fixture erneut getestet werden.

\hypertarget{quellenhinweis}{%
\section{\texorpdfstring{{14.7}
Quellenhinweis}{14.7 Quellenhinweis}}\label{quellenhinweis}}

Die methodennahen Aussagen in diesen Kapiteln folgen den Headern und
Implementierungen \texttt{ASMTAssembly.*}, \texttt{ASMTItem.*},
\texttt{ASMTSpatialItem.*}, \texttt{ASMTSpatialContainer.*},
\texttt{ASMTPart.*}, \texttt{ASMTJoint.*}, \texttt{ASMTMotion.*},
\texttt{ASMTLimit.*} und den jeweiligen Unterklassen im eingebundenen
Submodul. Besonders wichtig sind \texttt{ASMTAssembly::readJoints()},
\texttt{readMotions()} und \texttt{readLimits()}: Sie definieren die im
konkreten Parserstand erkannten ASMT-Typen.

\hypertarget{kap-api-funktionsgruppen}{%
\chapter{\texorpdfstring{{15} Funktionsgruppen und
Parameterreferenz}{15 Funktionsgruppen und Parameterreferenz}}\label{kap-api-funktionsgruppen}}

\hypertarget{leseregel-und-vollstuxe4ndigkeitsgrenze}{%
\section{\texorpdfstring{{15.1} Leseregel und
Vollständigkeitsgrenze}{15.1 Leseregel und Vollständigkeitsgrenze}}\label{leseregel-und-vollstuxe4ndigkeitsgrenze}}

Der Solver besteht aus 319 C++-Headern. Seine Methoden sind nicht alle
gleichartige Anwendungsfunktionen: Ein Teil ist eine direkt nutzbare
ASMT-Fassade, ein großer Teil sind virtuelle Hooks, Ableitungen
einzelner Zwangsgleichungen, Matrixoperationen und Integratorschritte.
Dieser Katalog ordnet die öffentliche Headeroberfläche nach Aufgabe und
erklärt die Argumentrollen. Für konkrete Signaturen und Überladungen
bleibt der benannte Header maßgeblich; die numerischen
Implementierungshooks sind versionsgebunden und nicht als stabile ABI
zugesagt.

Notation: \texttt{shared\_ptr\textless{}T\textgreater{}} ist ein
geteilter Zeiger auf ein Solverobjekt; \texttt{FColDsptr},
\texttt{FRowDsptr}, \texttt{FMatDsptr} sind geteilte Zeiger auf
Double-Vektor, Double-Zeile und Double-Matrix. \texttt{size\_t}
bezeichnet eine nichtnegative Größe beziehungsweise einen Index. Die
ASMT-Fassade ist die normale Client-API, die
\texttt{System}-/Matrix-/Newton-Raphson-Schicht ist die
Erweiterungsebene.

\hypertarget{assembly-lebenszyklus-und-dateifunktionen}{%
\section{\texorpdfstring{{15.2} Assembly-Lebenszyklus und
Dateifunktionen}{15.2 Assembly-Lebenszyklus und Dateifunktionen}}\label{assembly-lebenszyklus-und-dateifunktionen}}

\texttt{ASMTAssembly} koordiniert Objektmodell, Übersetzung und Analyse.

\begin{longtable}[]{@{}
  >{\raggedright\arraybackslash}p{(\columnwidth - 2\tabcolsep) * \real{0.5000}}
  >{\raggedright\arraybackslash}p{(\columnwidth - 2\tabcolsep) * \real{0.5000}}@{}}
\toprule\noalign{}
\begin{minipage}[b]{\linewidth}\raggedright
Funktion
\end{minipage} & \begin{minipage}[b]{\linewidth}\raggedright
Parameter und Wirkung
\end{minipage} \\
\midrule\noalign{}
\endhead
\bottomrule\noalign{}
\endlastfoot
\texttt{With()} & Keine Parameter; erzeugt eine leere Assemblywurzel. \\
\texttt{assemblyFromFile(const\ std::string\&\ fileName)} &
\texttt{fileName}: Pfad der ASMT-Eingabe; parst eine Assembly. \\
\texttt{setNotes(const\ std::string\&\ str)} & \texttt{str}: Notiztext
für die Assembly. \\
\texttt{setFilename(const\ std::string\&\ filename)} &
\texttt{filename}: Dateiname/Metadatum der Assembly. \\
\texttt{setDebug(bool\ todebug)} & \texttt{todebug}: Debug-Ausgabe
beziehungsweise Diagnosepfad ein-/ausschalten. \\
\texttt{addPart(shared\_ptr\textless{}ASMTPart\textgreater{}\ part)} &
\texttt{part}: Körperobjekt; wird dem Partcontainer hinzugefügt und
erhält diese Assembly als \texttt{owner}. \\
\texttt{addJoint(shared\_ptr\textless{}ASMTJoint\textgreater{}\ joint)}
& \texttt{joint}: Jointobjekt mit gültigen Markerpfaden; fügt hinzu und
setzt \texttt{owner}. \\
\texttt{addMotion(shared\_ptr\textless{}ASMTMotion\textgreater{}\ motion)}
& \texttt{motion}: Bewegungsvorgabe; fügt hinzu, setzt \texttt{owner}
und initialisiert Motion-Marker. \\
\texttt{addLimit(shared\_ptr\textless{}ASMTLimit\textgreater{}\ limit)}
& \texttt{limit}: Bewegungsgrenze; fügt hinzu, setzt \texttt{owner} und
initialisiert Marker. \\
\texttt{setConstantGravity(shared\_ptr\textless{}ASMTConstantGravity\textgreater{}\ constantGravity)}
& \texttt{constantGravity}: Schwerkraftvektor in konsistenten
Einheiten. \\
\texttt{setSimulationParameters(shared\_ptr\textless{}ASMTSimulationParameters\textgreater{}\ simulationParameters)}
& \texttt{simulationParameters}: Zeit-, Schrittweiten- und
Toleranzkonfiguration. \\
\texttt{solve()} & Keine Parameter; führt im aktuellen Stand den
initialen kinematischen Solve bei \texttt{tstart=tend=0} aus und
überschreibt zuvor gesetzte Simulationsparameter. \\
\texttt{runKINEMATIC()} & Keine Parameter; startet den konfigurierten
Kinematiklauf. Fängt \texttt{SimulationStoppingError} intern ab, liefert
keinen Erfolgsstatus. \\
\texttt{outputFile(const\ std::string\&\ filename)} & \texttt{filename}:
Zielpfad; serialisiert die ASMT-Struktur. Unter GCC erzeugt die aktuelle
Joint-Serialisierung ungültige Typnamen, siehe
\protect\hyperlink{kap-solve-ergebnisse}{Kapitel 14}. \\
\texttt{numberOfFrames()} & Keine Parameter; Anzahl gespeicherter
Zeitframes. \\
\texttt{updateForFrame(size\_t\ index)} & \texttt{index}: Frameindex in
den Ergebnisserien; setzt Objekte auf diesen Frame. \\
\texttt{geoTime()} & Keine Parameter; liefert den ASMT-Zeit-Ausdruck. \\
\end{longtable}

Die Loader-, Parser- und Writer-Familie besitzt zusätzliche
Abschnittsmethoden (\texttt{readNotes}, \texttt{readParts},
\texttt{readJoints}, \texttt{readMotions}, \texttt{readLimits},
\texttt{readForcesTorques}, \texttt{readConstantGravity},
\texttt{readSimulationParameters}, \texttt{readAnimationParameters},
\texttt{readTimeSeries}) und Serializer (\texttt{storeOnLevel*},
\texttt{storeOnTimeSeries}). Ihre Argumente sind jeweils
\texttt{std::vector\textless{}std::string\textgreater{}\&\ lines} für
den veränderbaren Parserrest beziehungsweise
\texttt{std::ofstream\&\ os} plus \texttt{size\_t\ level} für
Ausgabestrom und Einrückungsebene. Sie sind Erweiterungshooks, keine
üblichen Clientaufrufe. Nicht implementierte Parserzweige werfen
\texttt{SimulationStoppingError}.

Die übrigen öffentlichen \texttt{ASMTAssembly}-Methoden nach
Funktionsgruppe:

\begin{itemize}
\tightlist
\item
  \textbf{Demo-/Dateiroutinen:}
  \texttt{runSinglePendulumSuperSimplified()},
  \texttt{runSinglePendulumSuperSimplified2()},
  \texttt{runSinglePendulumSimplified()}, \texttt{runSinglePendulum()},
  \texttt{runDraggingLogTest()}, \texttt{runDraggingLogTest2()},
  \texttt{runDraggingLogTest3()}, \texttt{runDraggingTest()},
  \texttt{runDraggingTest2()}, \texttt{runDraggingTest3()} haben keine
  Argumente und sind eingebaute Test-/Beispielroutinen.
  \texttt{runFile(const\ std::string\&\ chars)} und
  \texttt{readWriteFile(const\ std::string\&\ chars)} erhalten Text
  beziehungsweise Dateiinhalte.
\item
  \textbf{Parserabschnitte:} \texttt{parseASMT(lines)},
  \texttt{readNotes(lines)}, \texttt{readParts(lines)},
  \texttt{readPart(lines)}, \texttt{readKinematicIJs(lines)},
  \texttt{readKinematicIJ(lines)}, \texttt{readConstraintSets(lines)},
  \texttt{readJoints(lines)}, \texttt{readMotions(lines)},
  \texttt{readLimits(lines)}, \texttt{readGeneralConstraintSets(lines)},
  \texttt{readForcesTorques(lines)},
  \texttt{readConstantGravity(lines)},
  \texttt{readSimulationParameters(lines)},
  \texttt{readAnimationParameters(lines)},
  \texttt{readTimeSeries(lines)}, \texttt{readTimes(lines)},
  \texttt{readAssemblySeries(lines)},
  \texttt{readPartSeriesMany(lines)}, \texttt{readPartSeries(lines)},
  \texttt{readJointSeriesMany(lines)}, \texttt{readJointSeries(lines)},
  \texttt{readMotionSeriesMany(lines)},
  \texttt{readMotionSeries(lines)}. \texttt{lines} ist ein mutierbarer
  Parsercursor; die konkrete Methode konsumiert die Zeilen ihres
  Abschnitts. \texttt{readKinematicIJ},
  \texttt{readGeneralConstraintSets} und einige Serien-/Referenzzweige
  sind im aktuellen Stand nicht vollständig implementiert.
\item
  \textbf{Solverlebenszyklus:}
  \texttt{preMbDrun(shared\_ptr\textless{}System\textgreater{}\ mbdSys)}
  erhält das Zielsystem; \texttt{preMbDrunDragStep(mbdSys,\ dragParts)}
  erhält zusätzlich die zu ziehenden MbD-Parts; \texttt{postMbDrun()}
  schließt den Lauf ab. \texttt{deleteMbD()},
  \texttt{createMbD(mbdSys,\ mbdUnits)}, \texttt{initialize()},
  \texttt{root()} und \texttt{runKINEMATIC()} sind parameterlos
  beziehungsweise erhalten System und Einheiten wie angegeben.
\item
  \textbf{Drag-/Zustandsfunktionen:} \texttt{runPreDrag()},
  \texttt{runDragStep(shared\_ptr\textless{}vector\textless{}shared\_ptr\textless{}ASMTPart\textgreater{}\textgreater{}\textgreater{}\ dragParts)},
  \texttt{runPostDrag()} führen Dragphasen aus; \texttt{restorePosRot()}
  stellt gespeicherte Lagen wieder her;
  \texttt{initprincipalMassMarker()} initialisiert den
  Assembly-Massemarker. \texttt{updateFromMbD()},
  \texttt{updateForFrame(index)}, \texttt{compareResults(type)} und
  \texttt{outputResults(type)} übernehmen Zustand, Frame oder
  Analyseart.
\item
  \textbf{Skalierung:} \texttt{calcCharacteristicDimensions()} berechnet
  charakteristische Größen; \texttt{calcCharacteristicTime()},
  \texttt{calcCharacteristicMass()} und
  \texttt{calcCharacteristicLength()} liefern die jeweiligen Skalen ohne
  Argumente.
\item
  \textbf{Lookup/Koordinaten:}
  \texttt{spatialContainerAt(self,\ longname)} erhält
  Assembly-Smartpointer und Pfad als veränderbare Zeichenkette;
  \texttt{partAt}, \texttt{markerAt}, \texttt{jointAt},
  \texttt{motionAt}, \texttt{forceTorqueAt} erhalten Pfadstrings.
  \texttt{vOcmO()} und \texttt{omeOpO()} liefern
  Assembly-Geschwindigkeit beziehungsweise Winkelgeschwindigkeit;
  \texttt{geoTime()} liefert Zeit.
\item
  \textbf{Ausgabecontainer:} \texttt{storeOnLevel(os,\ level)}
  serialisiert Rootobjekt;
  \texttt{storeOnLevelNotes/Parts/KinematicIJs/ConstraintSets/ForceTorques/Joints/Motions/Limits/GeneralConstraintSets(os,\ level)}
  schreiben den jeweiligen Abschnitt. \texttt{storeOnTimeSeries(os)}
  schreibt Zeitreihen. \texttt{connectorList()} und \texttt{markerMap()}
  benötigen keine Parameter und liefern die aufgebauten
  Connector-/Markerübersichten.
\end{itemize}

Suchfunktionen \texttt{partAt(path)}, \texttt{markerAt(path)},
\texttt{jointAt(path)}, \texttt{motionAt(path)} und
\texttt{forceTorqueAt(path)} nehmen einen vollständigen ASMT-Pfad.
\texttt{partAt} liefert bei Nichttreffer \texttt{nullptr};
\texttt{partNamed(name)} und \texttt{partPartialNamed(name)} müssen im
aktuellen Stand mit einem garantiert existierenden Treffer aufgerufen
werden. \texttt{partNamed} verwendet entgegen dem Namen den
vollständigen Pfadvergleich.

\hypertarget{abs-api-gemeinsame-funktionen}{%
\section{\texorpdfstring{{15.3} Gemeinsame Objekt- und
Übersetzungsfunktionen}{15.3 Gemeinsame Objekt- und Übersetzungsfunktionen}}\label{abs-api-gemeinsame-funktionen}}

\texttt{ASMTItem} ist Basisklasse der ASMT-Objekte. Folgende Methoden
sind nach ihrem Argumentmuster gruppiert:

\begin{longtable}[]{@{}
  >{\raggedright\arraybackslash}p{(\columnwidth - 2\tabcolsep) * \real{0.5000}}
  >{\raggedright\arraybackslash}p{(\columnwidth - 2\tabcolsep) * \real{0.5000}}@{}}
\toprule\noalign{}
\begin{minipage}[b]{\linewidth}\raggedright
Funktionsgruppe
\end{minipage} & \begin{minipage}[b]{\linewidth}\raggedright
Signatur-/Parameterrolle
\end{minipage} \\
\midrule\noalign{}
\endhead
\bottomrule\noalign{}
\endlastfoot
Identität & \texttt{setName(const\ std::string\&\ str)} setzt den Namen;
\texttt{classname()} liefert den Laufzeittypnamen;
\texttt{fullName(const\ std::string\&\ partialName)} bildet einen
vollständigen Pfad. \\
Baumzugriff & \texttt{root()}, \texttt{partOrAssembly()},
\texttt{part()} liefern Root, räumlichen Ownercontainer beziehungsweise
Partkontext; keine Eingabeparameter. \\
Lebenszyklus & \texttt{initialize()} initialisiert das ASMT-Objekt;
\texttt{createMbD(shared\_ptr\textless{}System\textgreater{}\ mbdSys,\ shared\_ptr\textless{}Units\textgreater{}\ mbdUnits)}
erhält das Zielsystem und die Einheitenskalen; \texttt{deleteMbD()}
räumt die Solverabbildung auf. \\
Zustand & \texttt{updateForFrame(size\_t\ index)} nimmt einen
Frameindex; \texttt{updateFromInputState()},
\texttt{updateFromInitiallyAssembledState()}, \texttt{updateFromMbD()}
übertragen Zustand zwischen Eingabe, Anfangslage und MbD. \\
Ergebnis & \texttt{compareResults(AnalysisType\ type)} und
\texttt{outputResults(AnalysisType\ type)} nehmen die gewünschte
Analyseart. \\
Zahlen-/Textparser & \texttt{readDouble(line)}, \texttt{readInt(line)},
\texttt{readSize\_t(line)}, \texttt{readBool(line)},
\texttt{readString(line)} interpretieren eine Textzeile als
angeforderten Typ;
\texttt{readRowOfDoubles(line)}/\texttt{readColumnOfDoubles(line)}
erzeugen Wertecontainer. \\
Parsercursor & \texttt{popOffTop(lines)},
\texttt{readStringOffTop(lines)},
\texttt{readRowOfDoublesOffTop(lines)},
\texttt{readColumnOfDoublesOffTop(lines)} konsumieren das erste Element
des veränderbaren Zeilenvektors. \\
Serialisierung & \texttt{storeOnLevel(os,\ level)} schreibt ein Objekt;
die Varianten
\texttt{storeOnLevelString/Double/Int/Size\_t/Bool/Array(os,\ level,\ value)}
erhalten Zielstrom, Einrückungsebene und Nutzwert.
\texttt{storeOnLevelTabs(os,\ level)} schreibt nur Einrückung. \\
Hilfen & \texttt{mbdUnits()} liefert den aktuellen Einheitenzeiger;
\texttt{sptrConstant(double\ value)} erzeugt einen Konstant-Ausdruck;
\texttt{logString(str)} protokolliert einen Text. \\
\end{longtable}

\texttt{owner} ist ein nicht besitzender Rohzeiger zur
Eltern-ASMT-Instanz; \texttt{mbdObject} hält das übersetzte MbD-Objekt.
Man sollte \texttt{owner} nicht selbst setzen, wenn ein passendes
\texttt{add*()} die Ownerzuordnung übernimmt.

Die konkrete öffentliche Methodenmenge von \texttt{ASMTItem} umfasst
außerdem \texttt{noop()}, \texttt{readDoublesInto(str,\ label,\ row)},
\texttt{storeOnLevelName(os,\ level)},
\texttt{storeOnLevelTabs(os,\ level)},
\texttt{storeOnLevelString(os,\ level,\ str)},
\texttt{storeOnLevelDouble(os,\ level,\ value)},
\texttt{storeOnLevelInt(os,\ level,\ value)},
\texttt{storeOnLevelSize\_t(os,\ level,\ value)},
\texttt{storeOnLevelBool(os,\ level,\ value)},
\texttt{storeOnLevelArray(os,\ level,\ values)},
\texttt{storeOnTimeSeries(os)} und \texttt{logString(str)}. Bei
\texttt{readDoublesInto} ist \texttt{str} die Eingabezeile,
\texttt{label} das erwartete Kennwort und \texttt{row} der Zielpuffer.
Bei allen \texttt{storeOnLevel*}-Methoden ist \texttt{os} der
Ausgabestrom und \texttt{level} die Tab-Einrückung; der letzte Parameter
ist der zu serialisierende Wert.

\hypertarget{mbd-basisklasse-item}{%
\subsection*{\texorpdfstring{MbD-Basisklasse
\texttt{Item}}{MbD-Basisklasse Item}}\label{mbd-basisklasse-item}}
\addcontentsline{toc}{subsection}{MbD-Basisklasse \texttt{Item}}

\texttt{Item} ist die Laufzeitbasis des numerischen Systems. Ihre
öffentlichen Funktionen lassen sich vollständig in diese Gruppen
einordnen:

\begin{itemize}
\tightlist
\item
  \textbf{Konstruktion/Baum:} \texttt{Item()}, \texttt{Item(str)},
  \texttt{root()}, \texttt{classname()}, \texttt{initialize()},
  \texttt{initializeLocally()}, \texttt{initializeGlobally()},
  \texttt{noop()}, \texttt{registerName()}. \texttt{str} ist der
  Objektname/-text.
\item
  \textbf{Constraint-/Residuum-/Jacobian-Ausgabe:}
  \texttt{fillConstraints(allConstraints)},
  \texttt{fillDispConstraints(dispConstraints)},
  \texttt{fillEssenConstraints(essenConstraints)},
  \texttt{fillPerpenConstraints(perpenConstraints)},
  \texttt{fillRedundantConstraints(redunConstraints)},
  \texttt{fillAccICIterError(col)}, \texttt{fillAccICIterJacob(mat)},
  \texttt{fillPosICError(col)}, \texttt{fillPosICJacob(mat)},
  \texttt{fillPosKineError(col)}, \texttt{fillPosKineJacob(mat)},
  \texttt{fillVelICError(col)}, \texttt{fillVelICJacob(mat)},
  \texttt{fillDynError(col)}, \texttt{fillStaticError(col)},
  \texttt{fillStaticJacob(mat)}, \texttt{fillCollisionError(col)},
  \texttt{fillCollisionpFpy(mat)}, \texttt{fillCollisionpFpydot(mat)},
  \texttt{fillCollisionDerivativeICError(col)},
  \texttt{fillCollisionDerivativeICJacob(mat)}, \texttt{fillpFpy(mat)},
  \texttt{fillpFpydot(mat)}. \texttt{col} ist der Zielvektor für skalare
  Gleichungsbeiträge; \texttt{mat} die Zielmatrix für
  Jacobian-/Ableitungsbeiträge; \texttt{*Constraints} ist ein
  Zielcontainer.
\item
  \textbf{Generalisierte Zustände/Gewichte:} \texttt{fillpqsumu(col)},
  \texttt{fillpqsumudot(col)}, \texttt{fillqsu(col)},
  \texttt{fillqsudot(col)}, \texttt{fillqsulam(col)},
  \texttt{fillqsudotPlam(col)}, \texttt{fillqsudotPlamDeriv(col)},
  \texttt{fillqsuddotlam(col)}, \texttt{fillqsuWeights(diagMat)},
  \texttt{fillqsudotWeights(diagMat)}. Die Vektoren enthalten
  Koordinaten, Ableitungen und Multiplikatoren; \texttt{diagMat} enthält
  deren Gewichtung.
\item
  \textbf{Zustand setzen:} \texttt{setpqsumu(col)},
  \texttt{setpqsumudot(col)}, \texttt{setpqsumuddot(col)},
  \texttt{setqsu(qsuOld)}, \texttt{setqsudot(col)},
  \texttt{setqsulam(col)}, \texttt{setqsudotlam(col)},
  \texttt{setqsudotPlam(col)}, \texttt{setqsudotPlamDeriv(col)},
  \texttt{setqsuddotlam(col)}. Jeder Vektor muss zur jeweiligen
  Zustands-/Gleichungsreihenfolge passen.
\item
  \textbf{Zeit-/Ereignisdiagnose:}
  \texttt{checkForCollisionDiscontinuityBetweenand(impulsePrevious,\ impulse)},
  \texttt{checkForDynDiscontinuityBetweenand(tprev,\ t)},
  \texttt{discontinuityAtaddTypeTo(t,\ disconTypes)},
  \texttt{discontinuityAtICAddTo(disconTypes)},
  \texttt{normalImpulse(imp)},
  \texttt{suggestSmallerOrAcceptCollisionFirstStepSize(hnew)},
  \texttt{suggestSmallerOrAcceptCollisionStepSize(hnew)},
  \texttt{suggestSmallerOrAcceptDynFirstStepSize(hnew)},
  \texttt{suggestSmallerOrAcceptDynStepSize(hnew)}. Die Argumente sind
  vorheriger/neuer Impuls, alte/neue Zeit, Ereignistypcontainer, Impuls
  beziehungsweise angefragte neue Schrittweite.
\item
  \textbf{Ablaufhooks:} \texttt{preAccIC}, \texttt{preCollision},
  \texttt{preCollisionCorrector},
  \texttt{preCollisionCorrectorIteration},
  \texttt{preCollisionDerivativeIC}, \texttt{preCollisionPredictor},
  \texttt{preCollisionStep}, \texttt{preDyn}, \texttt{preDynCorrector},
  \texttt{preDynCorrectorIteration}, \texttt{preDynFirstStep},
  \texttt{preDynOutput}, \texttt{preDynPredictor}, \texttt{preDynStep},
  \texttt{preICRestart}, \texttt{prePosIC}, \texttt{prePosKine},
  \texttt{preStatic}, \texttt{preVelIC}; passende \texttt{postAccIC},
  \texttt{postAccICIteration}, \texttt{postCollisionCorrector},
  \texttt{postCollisionCorrectorIteration},
  \texttt{postCollisionDerivativeIC}, \texttt{postCollisionPredictor},
  \texttt{postCollisionStep}, \texttt{postDyn},
  \texttt{postDynCorrector}, \texttt{postDynCorrectorIteration},
  \texttt{postDynFirstStep}, \texttt{postDynOutput},
  \texttt{postDynPredictor}, \texttt{postDynStep}, \texttt{postInput},
  \texttt{postPosIC}, \texttt{postPosICIteration}, \texttt{postStatic},
  \texttt{postStaticIteration}, \texttt{postVelIC}. Sie haben keine
  Parameter und markieren Phasenwechsel für Unterklassen.
\item
  \textbf{Weitere Hooks/Protokoll:}
  \texttt{calcPostDynCorrectorIteration()},
  \texttt{constraintsReport()}, \texttt{isJointForce()},
  \texttt{isJointTorque()}, \texttt{isKinedotIJ()}, \texttt{isKineIJ()},
  \texttt{reactivateRedundantConstraints()},
  \texttt{removeRedundantConstraints(redundantEqnNos)},
  \texttt{simUpdateAll()}, \texttt{stateData()},
  \texttt{storeCollisionState()}, \texttt{storeDynState()},
  \texttt{useEquationNumbers()}, \texttt{logString(str)},
  \texttt{logStringwithArgument(chars,\ chars1)},
  \texttt{logStringwithArguments(chars,\ arrayOfChars)},
  \texttt{printOn(stream)}. \texttt{redundantEqnNos} ist der zu
  entfernende Gleichungsindexsatz; Zeichenparameter sind
  Diagnosemeldungen.
\end{itemize}

Diese \texttt{fill*}- und \texttt{pre/post*}-Funktionen sind
Callback-Hooks des Solverzyklus. Ein Clientprogramm ruft sie
normalerweise nicht einzeln auf.

\hypertarget{ruxe4umliche-daten-und-massen}{%
\section{\texorpdfstring{{15.4} Räumliche Daten und
Massen}{15.4 Räumliche Daten und Massen}}\label{ruxe4umliche-daten-und-massen}}

\texttt{ASMTSpatialItem} und \texttt{ASMTSpatialContainer} teilen
räumliche Funktionen mit Part, Marker und Assembly.

\begin{longtable}[]{@{}
  >{\raggedright\arraybackslash}p{(\columnwidth - 2\tabcolsep) * \real{0.5000}}
  >{\raggedright\arraybackslash}p{(\columnwidth - 2\tabcolsep) * \real{0.5000}}@{}}
\toprule\noalign{}
\begin{minipage}[b]{\linewidth}\raggedright
Funktion
\end{minipage} & \begin{minipage}[b]{\linewidth}\raggedright
Parameterbeschreibung
\end{minipage} \\
\midrule\noalign{}
\endhead
\bottomrule\noalign{}
\endlastfoot
\texttt{setPosition3D(FColDsptr\ position3D)} & 3D-Ortsvektor als
Spaltenvektor. \\
\texttt{setPosition3D(double\ x,\ double\ y,\ double\ z)} & Drei
Ortskomponenten in Modelllängeneinheiten. \\
\texttt{getPosition3D(double\&\ x,\ double\&\ y,\ double\&\ z)} & Gibt
drei Ortskomponenten über Referenzen zurück. \\
\texttt{getPosition3D(size\_t\ frame)} & Liefert den Ortsvektor für
einen Ergebnisframe. \\
\texttt{setRotationMatrix(FMatDsptr\ rotationMatrix)} /
\texttt{setRotationMatrix(9\ Skalare)} & Orientierungsmatrix
beziehungsweise ihre zeilenweise neun Komponenten. \\
\texttt{getRotationMatrix(size\_t\ frame)} & Orientierungsmatrix im
Ergebnisframe. \\
\texttt{setQuarternions(q0,\ q1,\ q2,\ q3)} /
\texttt{getQuarternions(...)} & Quaternion setzen/auslesen; \texttt{q0}
ist die skalare Komponente, weitere Parameter bilden den Vektoranteil.
Der Rechtschreibfehler ist API-Name. \\
\texttt{setVelocity3D(vector)} / \texttt{setVelocity3D(x,y,z)} &
Lineargeschwindigkeit als Vektor oder Komponenten. \\
\texttt{setOmega3D(vector)} / \texttt{setOmega3D(x,y,z)} &
Winkelgeschwindigkeit als Vektor oder Komponenten. \\
\texttt{getVelocity3D(frame)}, \texttt{getOmega3D(frame)},
\texttt{getAcceleration3D(frame)}, \texttt{getAlpha3D(frame)} &
Frameindex; liefert lineare/Winkel-Geschwindigkeit beziehungsweise
-Beschleunigung. \\
\texttt{updateMbDFromPosition3D(vector)} / \texttt{(x,y,z)} &
ASMT-Position in die Solverabbildung übertragen. \\
\texttt{updateMbDFromRotationMatrix(matrix)} / \texttt{(9\ Skalare)} &
ASMT-Orientierung in die Solverabbildung übertragen. \\
\texttt{setPrincipalMassMarker(shared\_ptr\textless{}ASMTPrincipalMassMarker\textgreater{}\ marker)}
& Verknüpft Masse-/Trägheitsdaten mit Container. \\
\texttt{addMarker(shared\_ptr\textless{}ASMTMarker\textgreater{}\ marker)},
\texttt{addRefPoint(shared\_ptr\textless{}ASMTRefPoint\textgreater{}\ point)}
& Fügt lokale Referenzobjekte hinzu. \\
\texttt{generateUniqueMarkerName()} / \texttt{markerList()} & Erzeugt
einen eindeutigen Marker-Namen beziehungsweise liefert Marker des
Containers. \\
\end{longtable}

\texttt{ASMTPrincipalMassMarker} stellt \texttt{setMass(double\ mass)},
\texttt{setDensity(double\ density)} sowie
\texttt{setMomentOfInertias(DiagMatDsptr\ values)} oder drei diagonale
Trägheitswerte bereit. Masse, Dichte und Trägheitsmomente müssen
derselben Einheitenkonvention folgen. \texttt{ASMTPart::With()} erzeugt
ein initialisiertes Partobjekt; Position, Orientierung, Referenzpunkte,
Massemarker und \texttt{isFixed} beschreiben seinen Solverzustand, keine
CAD-B-Rep.

\hypertarget{verbindung-constraints-und-gelenke}{%
\section{\texorpdfstring{{15.5} Verbindung, Constraints und
Gelenke}{15.5 Verbindung, Constraints und Gelenke}}\label{verbindung-constraints-und-gelenke}}

\texttt{ASMTItemIJ} ist die gemeinsame I/J-Schnittstelle für Gelenke,
Kräfte und Constraints:

\begin{itemize}
\tightlist
\item
  \texttt{setMarkerI(const\ std::string\&\ path)} und
  \texttt{setMarkerJ(path)}: vollständige Pfade zu den beiden
  Markerbezügen.
\item
  \texttt{createMbD(shared\_ptr\textless{}System\textgreater{}\ system,\ shared\_ptr\textless{}Units\textgreater{}\ units)}:
  interne Übersetzung in das Zielsysten und dessen Einheiten.
\item
  \texttt{parseASMT(lines)}: liest Markerpfade und optionale Zeitreihen
  aus dem Parserrest.
\item
  \texttt{storeOnLevel(os,\ level)} und \texttt{storeOnTimeSeries(os)}:
  serialisiert Struktur beziehungsweise Zeitreihen.
\item
  \texttt{mbdClassNew()}: virtuelle Fabrik für die konkrete
  MbD-Constraint-/Jointklasse.
\end{itemize}

Die konkreten Joint-Fabriken \texttt{ASMTFixedJoint::With()},
\texttt{ASMTRevoluteJoint::With()}, \texttt{ASMTSphericalJoint::With()}
usw. haben keine Parameter; die Beziehung kommt über Marker I/J und
typspezifische Member. Der ASMT-Parser kennt die im
\protect\hyperlink{kap-joints}{Kapitel 12} gelisteten festen
Joint-Typstrings. Gelenkgruppen sind: Punktgelenke,
Achsen-/Liniengelenke, Ebenengelenke, Orientierungsbedingungen,
Getriebe-/Schraubenverhältnisse und zusammengesetzte Gelenke. Winkel,
Übersetzungsverhältnis, Radien und Steigung sind Größen des jeweiligen
Jointtyps und keine gemeinsamen \texttt{ASMTJoint}-Parameter.

\texttt{Constraint} im MbD-Kern verwaltet Gleichungen und Reaktionen.
Öffentliche Funktionsgruppen:

\begin{longtable}[]{@{}
  >{\raggedright\arraybackslash}p{(\columnwidth - 2\tabcolsep) * \real{0.5000}}
  >{\raggedright\arraybackslash}p{(\columnwidth - 2\tabcolsep) * \real{0.5000}}@{}}
\toprule\noalign{}
\begin{minipage}[b]{\linewidth}\raggedright
Methoden
\end{minipage} & \begin{minipage}[b]{\linewidth}\raggedright
Argumente und Rolle
\end{minipage} \\
\midrule\noalign{}
\endhead
\bottomrule\noalign{}
\endlastfoot
\texttt{fillConstraints(list)}, \texttt{fillDispConstraints(list)},
\texttt{fillEssenConstraints(list)},
\texttt{fillPerpenConstraints(list)},
\texttt{fillRedundantConstraints(list)} & Zielvektor, in den dieses
Constraint oder seine Unterconstraints aufgenommen werden. Überladungen
mit \texttt{shared\_ptr\textless{}Constraint\textgreater{}\ self} plus
Zielvektor erlauben rekursive Hierarchien. \\
\texttt{fillPosICError(column)}, \texttt{fillPosKineError(column)},
\texttt{fillAccICIterJacob(matrix)} & Zielspaltenvektor beziehungsweise
Sparse-Matrix, die mit Fehler- oder Jacobianbeiträgen gefüllt wird. \\
\texttt{fillqsulam(column)}, \texttt{fillqsudotlam(column)},
\texttt{fillqsuddotlam(column)} & Zielvektor für generalisierte
Koordinaten-/Ableitungs- und Lagrange-Multiplikatorbeiträge. \\
\texttt{setqsulam(column)}, \texttt{setqsudotlam(column)},
\texttt{setqsuddotlam(column)} & Quellvektor mit
Zustands-/Multiplikatorwerten zum Setzen des Constraint-Zustands. \\
\texttt{addToJointForceI(column)}, \texttt{addToJointTorqueI(column)} &
Zielvektor für Reaktionskraft beziehungsweise Reaktionsmoment auf Seite
I. \\
\texttt{removeRedundantConstraints(shared\_ptr\textless{}vector\textless{}size\_t\textgreater{}\textgreater{}\ equationNumbers)}
& Indizes der als redundant erkannten Gleichungen. \\
\texttt{setConstant(double\ value)} & Konstante der skalaren
Constraint-Gleichung. \\
\texttt{isRedundant()}, \texttt{type()}, \texttt{constraintSpec()} &
Keine Parameter; Redundanzstatus, Constraint-Enum und symbolische
Spezifikation. \\
\end{longtable}

Die einzelnen Klassen \texttt{DistanceConstraint*},
\texttt{DirectionCosineConstraint*}, \texttt{AtPointConstraint*},
\texttt{TranslationConstraint*}, \texttt{AngleZConstraint*},
\texttt{GearConstraint*}, \texttt{ScrewConstraint*},
\texttt{RackPinConstraint*} und verwandte
\texttt{Dist*}/\texttt{DispComp*} bilden Varianten derselben
Funktionsfamilien für verschiedene Bezugssysteme und
Koordinatenrepräsentationen. \texttt{I/J} benennt die zwei
Körper-/Marker-Seiten; \texttt{ec}, \texttt{qc} und \texttt{qct} im
Klassennamen kodieren die jeweilige Endframe-/Orientierungsdarstellung.
Die Varianten liefern Gleichungswerte, Ableitungen und Jacobianbeiträge
und werden überwiegend intern aus dem Joint erzeugt.

\hypertarget{abs-api-motion-limits-lasten}{%
\section{\texorpdfstring{{15.6} Motion, Limits und
Lasten}{15.6 Motion, Limits und Lasten}}\label{abs-api-motion-limits-lasten}}

\begin{longtable}[]{@{}
  >{\raggedright\arraybackslash}p{(\columnwidth - 2\tabcolsep) * \real{0.5000}}
  >{\raggedright\arraybackslash}p{(\columnwidth - 2\tabcolsep) * \real{0.5000}}@{}}
\toprule\noalign{}
\begin{minipage}[b]{\linewidth}\raggedright
Typ/Funktion
\end{minipage} & \begin{minipage}[b]{\linewidth}\raggedright
Parameter und Wirkung
\end{minipage} \\
\midrule\noalign{}
\endhead
\bottomrule\noalign{}
\endlastfoot
\texttt{ASMTMotion::initMarkers()} & Keine expliziten Parameter; setzt
interne Markerreferenzen für die konkrete Motion. \\
\texttt{ASMTMotion::createMbD(system,\ units)} & Zielsystem und
Einheitenskalen für die Übersetzung. \\
\texttt{ASMTRotationalMotion}, \texttt{ASMTTranslationalMotion},
\texttt{ASMTGeneralMotion}, \texttt{ASMTAllowRotation} & Fabriken ohne
Argumente; Bewegungskomponenten/-gesetze liegen in typgebundenen Feldern
und Funktionen. \\
\texttt{ASMTLimit::setmotionJoint(const\ std::string\&\ path)} &
Vollständiger ASMT-Pfad zur begrenzten Bewegung oder zum begrenzten
Joint. \\
\texttt{ASMTRotationLimit}, \texttt{ASMTTranslationLimit} & Grenzen für
Winkel- beziehungsweise Wegkoordinate; konkrete Grenzwerte stehen in der
Unterklasse. \\
\texttt{ASMTConstantGravity::setg(FColDsptr\ g)} / \texttt{setg(x,y,z)}
& Beschleunigungsvektor als Spaltenvektor oder drei Komponenten. \\
\texttt{ASMTForceTorque} & I/J-Pfade plus Zeitreihenkomponenten
\texttt{FX/FY/FZ} und \texttt{TX/TY/TZ}; Kräfte und Momente wirken im
jeweiligen Referenzsystem. \\
\texttt{setConstantGravity(gravity)}, \texttt{addLimit(limit)},
\texttt{addMotion(motion)} & Assembly-Einbindung; \texttt{addLimit} und
\texttt{addMotion} initialisieren zusätzlich die Marker. \\
\end{longtable}

\texttt{ASMTSimulationParameters} besitzt die Setter
\texttt{settstart(double\ tstart)}, \texttt{settend(double\ tend)},
\texttt{sethmin(double\ hmin)}, \texttt{sethmax(double\ hmax)},
\texttt{sethout(double\ hout)}, \texttt{seterrorTol(double\ errorTol)}
und \texttt{setmaxIter(size\_t\ maxIter)}. Zeitwerte/Schrittweiten
verwenden die Zeiteinheit des Modells; \texttt{errorTol} ist eine
numerische Toleranz, \texttt{maxIter} eine Iterationsanzahl. Öffentliche
Zusatzfelder umfassen positions-/beschleunigungsbezogene Toleranzen,
absolute/relative Korrektur- und Integrationstoleranzen,
\texttt{translationLimit}, \texttt{rotationLimit}, Iterationsgrenzen
sowie \texttt{orderMax}. \texttt{solve()} ersetzt wesentliche
Einstellungen mit seinen Defaults; für eigene Parameter ist
\texttt{runKINEMATIC()} nötig.

\hypertarget{system--und-solversteuerung}{%
\section{\texorpdfstring{{15.7} System- und
Solversteuerung}{15.7 System- und Solversteuerung}}\label{system--und-solversteuerung}}

\texttt{System} ist der MbD-Container für Parts, Joints/Motions, Limits,
Kräfte und Zeit. Wichtige Methoden und Parameterrollen:

\begin{itemize}
\tightlist
\item
  \texttt{System()} beziehungsweise
  \texttt{System(const\ std::string\&\ str)}: leeres System oder System
  mit Name/Text.
\item
  \texttt{initialize()}, \texttt{initializeLocally()},
  \texttt{initializeGlobally()}, \texttt{clear()}: Gesamt-, lokale,
  globale Initialisierung beziehungsweise Zurücksetzen.
\item
  \texttt{runKINEMATIC(shared\_ptr\textless{}System\textgreater{}\ self)}:
  Systeminstanz als Selbstreferenz; startet Kinematiklauf.
\item
  \texttt{runPreDrag(self)} /
  \texttt{runDragStep(self,\ shared\_ptr\textless{}vector\textless{}shared\_ptr\textless{}Part\textgreater{}\textgreater{}\textgreater{}\ dragParts)}:
  interaktiver Drag-Vorlauf und Schritt mit den gezogenen Körpern.
\item
  \texttt{addPart(part)}, \texttt{addJoint(joint)},
  \texttt{addMotion(motion)}, \texttt{addLimit(limit)},
  \texttt{addForceTorque(force)}: fügt ein MbD-Objekt der passenden
  Systemfamilie hinzu.
\item
  \texttt{jointsMotionsDo(function\textless{}void(shared\_ptr\textless{}Joint\textgreater{})\textgreater{}\ f)}
  und
  \texttt{partsJointsMotionsDo(function\textless{}void(shared\_ptr\textless{}Item\textgreater{})\textgreater{}\ f)}
  (sowie Forces-/Limits-Varianten): Callback \texttt{f} wird für jedes
  Objekt der benannten Gruppen ausgeführt.
\item
  \texttt{mbdTimeValue()} liest aktuelle Zeit;
  \texttt{mbdTimeValue(double\ t)} setzt Zeit \texttt{t}.
\item
  \texttt{allConstraints()}, \texttt{essentialConstraints()},
  \texttt{displacementConstraints()},
  \texttt{perpendicularConstraints()},
  \texttt{allRedundantConstraints()}, \texttt{allConstraintsLimits()}:
  keine Argumente; liefern klassifizierte Constraint-Listen.
\item
  \texttt{outputFor(AnalysisType\ type)}: schreibt/aktiviert Ausgabe der
  Analyseart \texttt{type}; \texttt{limitsSatisfied()} prüft Grenzen;
  \texttt{deactivateLimits()} deaktiviert aktive Limits.
\item
  \texttt{maximumMass()}, \texttt{maximumMomentOfInertia()},
  \texttt{translationLimit()}, \texttt{rotationLimit()}: keine
  Parameter; liefern charakteristische Größen/Grenzen des Systems.
\end{itemize}

\hypertarget{part-und-joint-kuxf6rperzustand-und-reaktionen}{%
\subsection*{\texorpdfstring{\texttt{Part} und \texttt{Joint}:
Körperzustand und
Reaktionen}{Part und Joint: Körperzustand und Reaktionen}}\label{part-und-joint-kuxf6rperzustand-und-reaktionen}}
\addcontentsline{toc}{subsection}{\texttt{Part} und \texttt{Joint}:
Körperzustand und Reaktionen}

\texttt{Part} repräsentiert einen dynamischen Körper. Seine direkten
Zustandsfunktionen sind:

\begin{itemize}
\tightlist
\item
  \texttt{setqX(FColDsptr\ x)}, \texttt{qX(FColDsptr\ x)},
  \texttt{getqX()}, \texttt{qX()}: Translationskoordinate setzen/lesen.
\item
  \texttt{setqXdot(x)}, \texttt{qXdot(x)}, \texttt{getqXdot()},
  \texttt{qXdot()}: translatorische Geschwindigkeit setzen/lesen.
\item
  \texttt{setqXddot(x)}, \texttt{qXddot(x)}, \texttt{getqXddot()},
  \texttt{qXddot()}: translatorische Beschleunigung setzen/lesen.
\item
  \texttt{setqE(FColDsptr\ x)}, \texttt{getqE()},
  \texttt{qE(shared\_ptr\textless{}EulerParameters\textless{}double\textgreater{}\textgreater{}\ x)},
  \texttt{qE()}: Orientierung setzen/lesen; die \texttt{qE}-Überladung
  nimmt Eulerparameter.
\item
  \texttt{setqEddot(x)}, \texttt{getqEddot()}, \texttt{qEddot(x)},
  \texttt{qEddot()}, \texttt{qEdot()}: Eulerparameter-Ableitungen
  setzen/lesen; \texttt{qEdot()} liefert die erste Ableitung.
\item
  \texttt{setomeOpO(FColDsptr\ x)}, \texttt{getomeOpO()},
  \texttt{omeOpO(x)}, \texttt{omeOpO()}: Winkelgeschwindigkeit
  setzen/lesen.
\item
  \texttt{setaAap(FMatDsptr\ mat)}, \texttt{aAOp()}, \texttt{alpOpO()}:
  Frameorientierung setzen; Orientierungs- und
  Winkelbeschleunigungsgrößen abfragen.
\item
  \texttt{setSystem(System*\ sys)}, \texttt{asFixed()}: Systemreferenz
  setzen beziehungsweise Körper fixieren.
\item
  \texttt{iqX()}, \texttt{iqE()}, \texttt{iqX(size\_t\ eqnNo)},
  \texttt{iqE(size\_t\ eqnNo)}: Gleichungsnummern für Translation und
  Orientierung abfragen/setzen.
\item
  \texttt{calcp()}, \texttt{calcmEdot()}, \texttt{calcpTpE()},
  \texttt{calcppTpEpE()}, \texttt{calcppTpEpEdot()}, \texttt{calcmE()}:
  interne kinematische/Massenmatrizen ohne Argument berechnen.
\end{itemize}

\texttt{Part} überschreibt außerdem Solverhooks aus \texttt{Item}:
\texttt{postInput}, \texttt{calcPostDynCorrectorIteration},
\texttt{prePosIC}, \texttt{prePosKine}, \texttt{fillEssenConstraints},
\texttt{fillRedundantConstraints}, \texttt{fillConstraints},
\texttt{fillqsu}, \texttt{fillqsuWeights}, \texttt{fillqsuddotlam},
\texttt{fillqsulam}, \texttt{fillqsudot}, \texttt{fillqsudotWeights},
\texttt{useEquationNumbers}, \texttt{setqsu}, \texttt{setqsulam},
\texttt{setqsudot}, \texttt{setqsudotlam}, \texttt{postPosICIteration},
\texttt{fillPosICError}, \texttt{fillPosICJacob},
\texttt{removeRedundantConstraints},
\texttt{reactivateRedundantConstraints}, \texttt{constraintsReport},
\texttt{postPosIC}, \texttt{preDyn}, \texttt{storeDynState},
\texttt{fillPosKineError}, \texttt{fillPosKineJacob}, \texttt{preVelIC},
\texttt{postVelIC}, \texttt{fillVelICError}, \texttt{fillVelICJacob},
\texttt{preAccIC}, \texttt{fillAccICIterError},
\texttt{fillAccICIterJacob}, \texttt{setqsuddotlam}, \texttt{stateData},
\texttt{suggestSmallerOrAcceptDynStepSize(hnew)}, \texttt{postDynStep}
und \texttt{postAccIC}. Die Vektor-/Matrixparameter folgen den
Bedeutungen aus
\protect\hyperlink{abs-api-gemeinsame-funktionen}{Abschnitt 15.3};
\texttt{hnew} ist die angefragte Schrittweite.

\texttt{Joint} ergänzt \texttt{ConstraintSet} um
\texttt{initializeLocally()}, \texttt{constraintsReport()},
\texttt{fillRedundantConstraints(redunConstraints)},
\texttt{reactivateRedundantConstraints()},
\texttt{removeRedundantConstraints(equationNumbers)} und
\texttt{stateData()}. Getter \texttt{aFIeJtIe()}, \texttt{aFIeJtO()},
\texttt{aFX()}, \texttt{aTIeJtIe()}, \texttt{aTIeJtO()}, \texttt{aTX()},
\texttt{jointForceI()} und \texttt{jointTorqueI()} benötigen keine
Eingaben und liefern Reaktionskraft-/Reaktionsmomentvektoren in den
durch den Namen bezeichneten Frames. \texttt{Joint()} und
\texttt{Joint(const\ std::string\&\ str)} erzeugen ein unbenanntes
beziehungsweise benanntes MbD-Joint.

\texttt{Solver} hat zusätzlich den Lifecycle \texttt{noop()},
\texttt{initialize()}, \texttt{initializeLocally()},
\texttt{initializeGlobally()}, \texttt{assignEquationNumbers()},
\texttt{preRun()}, \texttt{run()}, \texttt{finalize()},
\texttt{reportStats()}, \texttt{postRun()},
\texttt{handleSingularMatrix()} und
\texttt{logString(const\ std::string\&\ str)}.
\texttt{setSystem(Solver*\ sys)} ist abstrakt und bindet den Solver;
\texttt{sys} ist ein nicht besitzender Zeiger. Die parameterlosen
Methoden initialisieren, nummerieren Gleichungen, starten/abschließen
einen Lauf oder melden Statistik/Matrixfehler.

\texttt{ExternalSystem} verbindet MbD mit CAD/ASMT:
\texttt{preMbDrun(shared\_ptr\textless{}System\textgreater{}\ mbdSys)},
\texttt{preMbDrunDragStep(mbdSys,\ dragParts)},
\texttt{updateFromMbD()}, \texttt{outputFor(AnalysisType\ type)},
\texttt{logString(string\textbar{}double)}, \texttt{runOndselPiston()},
\texttt{runPiston()} und \texttt{postMbDrun()}. \texttt{mbdSys} ist das
aktuelle System; \texttt{dragParts} die Liste bewegter Parts;
\texttt{type} die Analyseart; Logargumente sind Text oder Zahl. Die
Zeiger \texttt{cadSystem}, \texttt{asmtAssembly} und
\texttt{freecadAssemblyObject} verweisen auf die angebundenen
Hostobjekte.

\texttt{StateData()} und \texttt{initialize()} erzeugen/initialisieren
einen Zustandscontainer. \texttt{PrescribedMotion} hat Konstruktoren
ohne Argument oder mit \texttt{str}, \texttt{connectsItoJ(frmI,frmJ)}
mit zwei Endframes und \texttt{initMotions()}; die sechs Ausdrücke
\texttt{xBlk}, \texttt{yBlk}, \texttt{zBlk}, \texttt{phiBlk},
\texttt{theBlk}, \texttt{psiBlk} beschreiben Translation/Rotation.
\texttt{ConstantGravity::fillAccICIterError(col)} fügt den
Gravitationsbeitrag in den Zielvektor \texttt{col} der
Beschleunigungs-Anfangsbedingung ein.

\hypertarget{systemsolver-uxf6ffentliche-ablaufmethoden}{%
\subsection*{\texorpdfstring{\texttt{SystemSolver}: öffentliche
Ablaufmethoden}{SystemSolver: öffentliche Ablaufmethoden}}\label{systemsolver-uxf6ffentliche-ablaufmethoden}}
\addcontentsline{toc}{subsection}{\texttt{SystemSolver}: öffentliche
Ablaufmethoden}

\texttt{SystemSolver(System*\ x)} erhält das MbD-\texttt{System}, dessen
Solver er verwaltet; \texttt{setSystem(Solver*\ sys)} setzt die
Solververknüpfung. Die Laufmethoden haben keine Parameter, außer
Dragvarianten:

\begin{itemize}
\tightlist
\item
  \textbf{Initialbedingungen:} \texttt{runAllIC()}, \texttt{runPosIC()},
  \texttt{runVelIC()}, \texttt{runAccIC()}, \texttt{needToRedoPosIC()}.
\item
  \textbf{Kinematik:} \texttt{runBasicKinematic()},
  \texttt{runQuasiKinematic()}, \texttt{runPosKine()},
  \texttt{runVelKine()}, \texttt{runAccKine()}, \texttt{runPosICKine()},
  \texttt{runVelICKine()}, \texttt{runAccICKine()}.
\item
  \textbf{Kollision/Drag:} \texttt{preCollision()},
  \texttt{runCollisionDerivativeIC()}, \texttt{runBasicCollision()},
  \texttt{runPreDrag()}, \texttt{runDragStep(dragParts)},
  \texttt{runPosICDrag(dragParts)},
  \texttt{runPosICDragLimit(dragParts)}. \texttt{dragParts} ist eine
  Liste der zu ziehenden MbD-Parts.
\item
  \textbf{Constraints/Iteration:} \texttt{postNewtonRaphson()},
  \texttt{discontinuityBlock()}, \texttt{partsJointsMotionsDo(f)},
  \texttt{partsJointsMotionsForcesTorquesDo(f)},
  \texttt{partsJointsMotionsLimitsDo(f)},
  \texttt{partsJointsMotionsLimitsForcesTorquesDo(f)}. \texttt{f} ist
  ein Callback pro Objekt der genannten Kategorien.
\item
  \textbf{Constraintabfragen:} \texttt{parts()}, \texttt{limits()},
  \texttt{essentialConstraints()}, \texttt{displacementConstraints()},
  \texttt{perpendicularConstraints()},
  \texttt{allRedundantConstraints()}, \texttt{allConstraints()},
  \texttt{allConstraintsLimits()} liefern Objektlisten ohne Eingabe.
\item
  \textbf{Zeit/Ausgabe:} \texttt{startTime()},
  \texttt{outputStepSize()}, \texttt{maxStepSize()},
  \texttt{minStepSize()}, \texttt{firstOutputTime()}, \texttt{endTime()}
  lesen Werte; \texttt{settime(double\ tnew)} setzt die neue Systemzeit;
  \texttt{time(double\ t)} aktualisiert Solverzeit;
  \texttt{tstartPastsAddFirst(double\ t)} ergänzt einen Startzeitwert;
  \texttt{output()} schreibt einen Zustand; \texttt{logString(str)}
  protokolliert.
\item
  \textbf{Limits:} \texttt{limitsSatisfied()} liefert den Grenzstatus;
  \texttt{deactivateLimits()} deaktiviert aktive Grenzen.
\end{itemize}

Die öffentlichen Solverfelder \texttt{tstart}, \texttt{tend},
\texttt{toutFirst}, \texttt{hmin}, \texttt{hmax}, \texttt{hout},
\texttt{direction}, \texttt{corAbsTol}, \texttt{corRelTol},
\texttt{intAbsTol}, \texttt{intRelTol}, \texttt{iterMaxPosKine},
\texttt{iterMaxAccKine}, \texttt{iterMaxDyn}, \texttt{orderMax},
\texttt{translationLimit} und \texttt{rotationLimit} sind die
ausgewerteten Zeit-, Toleranz-, Iterations- und Bewegungsgrenzparameter.
C++-Direktzugriff auf diese Felder ist quellstandspezifisch.

\hypertarget{newton-raphson-kinematik-und-dynamik}{%
\section{\texorpdfstring{{15.8} Newton-Raphson, Kinematik und
Dynamik}{15.8 Newton-Raphson, Kinematik und Dynamik}}\label{newton-raphson-kinematik-und-dynamik}}

Newton-Raphson-Klassen lösen iterativ nichtlineare Gleichungssysteme.
Gemeinsame Parameterkonzepte:

\begin{itemize}
\tightlist
\item
  \texttt{x}: aktueller Lösungsvektor;
\item
  \texttt{deltaX}: Korrektur pro Iteration;
\item
  Residuum/Funktionsvektor: Gleichungsfehler am aktuellen Zustand;
\item
  Jacobian: partielle Ableitungen nach unbekannten Koordinaten;
\item
  \texttt{errorTol}: Abbruchschwelle;
\item
  \texttt{maxIter}/\texttt{iterMax*}: maximale Iterationszahl.
\end{itemize}

\texttt{NewtonRaphson}, \texttt{ScalarNewtonRaphson},
\texttt{VectorNewtonRaphson}, \texttt{SystemNewtonRaphson},
\texttt{PosNewtonRaphson}, \texttt{PosKineNewtonRaphson},
\texttt{PosICKineNewtonRaphson}, \texttt{PosICNewtonRaphson},
\texttt{PosICDragNewtonRaphson}, \texttt{PosICDragLimitNewtonRaphson},
\texttt{AccNewtonRaphson}, \texttt{AccKineNewtonRaphson},
\texttt{AccICNewtonRaphson} und \texttt{AccICKineNewtonRaphson} teilen
diese Iterationsaufgabe nach Skalar-/Vektorgleichung und Positions-,
Geschwindigkeits-, Beschleunigungs-, Anfangsbedingungs- oder
Drag-Zustand auf.

Typische Parameter in ihren Methoden sind Matrix-/Vektorpuffer (werden
vom Solver befüllt), Iterationsindex \texttt{i}/\texttt{iter},
Solver-/Systemzeiger und Zustandsspalten. Die einzelnen \texttt{pre*},
\texttt{post*}, \texttt{fill*}, \texttt{update*} und
\texttt{solve*}-Funktionen sind Lifecycle-Hooks: \texttt{pre} bereitet
einen Schritt vor, \texttt{fill} schreibt Residuum/Jacobian,
\texttt{solve} berechnet die Korrektur, \texttt{post} übernimmt den
Zustand. Sie sind keine Client-API für eigene Gleichungen ohne passende
Unterklasse.

Die konkrete Basisschnittstelle \texttt{NewtonRaphson} umfasst
\texttt{initialize()}, \texttt{initializeLocally()}, \texttt{run()},
\texttt{setSystem(Solver*\ sys)}, \texttt{iterate()}, \texttt{fillY()},
\texttt{fillPyPx()}, \texttt{calcyNorm()}, \texttt{calcdxNorm()},
\texttt{solveEquations()}, \texttt{incrementIterNo()},
\texttt{updatexold()}, \texttt{xEqualxoldPlusdx()},
\texttt{isConverged()}, \texttt{askSystemToUpdate()},
\texttt{passRootToSystem()}, \texttt{isConvergedToNumericalLimit()},
\texttt{calcDXNormImproveRootCalcYNorm()} und \texttt{postRun()}. Die
parameterlosen Funktionen sind algorithmische Hooks; \texttt{sys} ist
die Solververknüpfung. Öffentliche Zustandsfelder \texttt{dxNorm},
\texttt{yNorm}, \texttt{yNormOld}, \texttt{yNormTol}, \texttt{dxTol},
\texttt{twoAlp}, \texttt{lam}, \texttt{iterNo}, \texttt{iterMax},
\texttt{nDivergence}, \texttt{nBackTracking}, \texttt{dxNorms} und
\texttt{yNorms} enthalten Korrektur-/Residuumgrößen, Abbruchschwellen,
Suchparameter, Zähler und Konvergenzhistorie. \texttt{iterMax} begrenzt
die Iteration; \texttt{nBackTracking} begrenzt Rückwärtssuche.

\hypertarget{abs-api-zeitintegration}{%
\section{\texorpdfstring{{15.9}
Zeitintegration}{15.9 Zeitintegration}}\label{abs-api-zeitintegration}}

\texttt{Integrator} ist die abstrakte Zeitschrittsteuerung:

\begin{longtable}[]{@{}
  >{\raggedright\arraybackslash}p{(\columnwidth - 4\tabcolsep) * \real{0.3333}}
  >{\raggedright\arraybackslash}p{(\columnwidth - 4\tabcolsep) * \real{0.3333}}
  >{\raggedright\arraybackslash}p{(\columnwidth - 4\tabcolsep) * \real{0.3333}}@{}}
\toprule\noalign{}
\begin{minipage}[b]{\linewidth}\raggedright
Hook
\end{minipage} & \begin{minipage}[b]{\linewidth}\raggedright
Argument
\end{minipage} & \begin{minipage}[b]{\linewidth}\raggedright
Zweck
\end{minipage} \\
\midrule\noalign{}
\endhead
\bottomrule\noalign{}
\endlastfoot
\texttt{firstStep()}, \texttt{preFirstStep()}, \texttt{postFirstStep()}
& keine & erster Integrationsschritt und dessen Vor-/Nachbereitung \\
\texttt{subsequentSteps()}, \texttt{nextStep()}, \texttt{preStep()},
\texttt{postStep()} & keine & Folge-, Auswahl-, Vor- und
Nachbereitungsschritte \\
\texttt{runInitialConditionTypeSolution()} & keine &
Anfangsbedingungslösung statt normaler Zeitintegration \\
\texttt{iStep(size\_t\ i)} & \texttt{i}: aktuelle Schrittzahl &
Verarbeitung des angeforderten Integrationsschritts \\
\texttt{selectOrder()} & keine & Integrationsordnung auswählen \\
\texttt{setSystem(Solver*\ sys)} & \texttt{sys}: zu integrierendes
Solver-System & verknüpft Integrator und System \\
\end{longtable}

\texttt{BasicIntegrator}, \texttt{BasicQuasiIntegrator},
\texttt{KineIntegrator}, \texttt{ICKineIntegrator},
\texttt{QuasiIntegrator}, \texttt{LinearMultiStepMethod} und
\texttt{StableBackwardDifference} spezialisieren Schrittverfahren.
\texttt{hmin}, \texttt{hmax}, \texttt{hout}, \texttt{tstart},
\texttt{tend}, \texttt{direction} und \texttt{orderMax} sind zentrale
Steuergrößen: minimale/maximale/gewünschte Schrittweite, Zeitintervall,
Vorwärts-/Rückwärtsrichtung und maximale Methodeordnung. Zeiteinheiten
werden mit \texttt{Units::time} skaliert.

\texttt{IntegratorInterface} ergänzt \texttt{initializeGlobally()},
\texttt{preRun()}, \texttt{checkForDiscontinuity()},
\texttt{setSystem(sys)}, \texttt{logString(str)}, \texttt{run()},
\texttt{orderMax()}, \texttt{incrementTime(tnew)},
\texttt{postFirstStep()},
\texttt{suggestSmallerOrAcceptFirstStepSize(hnew)},
\texttt{suggestSmallerOrAcceptStepSize(hnew)},
\texttt{checkForOutputThrough(t)} und \texttt{interpolateAt(t)}.
\texttt{hnew} ist Kandidat für die nächste Schrittweite, \texttt{tnew}
die neue Systemzeit, \texttt{t} Abfrage-/Ausgabezeit. \texttt{preRun},
Discontinuityprüfung und Schrittweitenwahl sind virtuelle
Integratorpflichten; \texttt{run()} steuert die Schleife.

\hypertarget{abs-api-matrizen}{%
\section{\texorpdfstring{{15.10} Matrizen und
Vektoren}{15.10 Matrizen und Vektoren}}\label{abs-api-matrizen}}

Die lineare Algebra ist in dichte (\texttt{Full*}) und dünn besetzte
(\texttt{Sparse*}) Strukturen unterteilt.
\texttt{FullMatrix\textless{}T\textgreater{}} bietet Konstruktoren für
Matrixgröße \texttt{(m,n)}, Zeilenlisten oder verschachtelte Werte;
Parameter \texttt{m}/\texttt{n} sind Zeilen-/Spaltenzahl. Wichtige
Operationen:

\begin{itemize}
\tightlist
\item
  \texttt{rotatex(angle)}, \texttt{rotatey(angle)},
  \texttt{rotatez(angle)}: Winkel im Bogenmaß; liefern
  3×3-Rotationsmatrix.
\item
  \texttt{rotate*rotDot(angle,\ angleDot)} und
  \texttt{rotate*rotDotrotDDot(angle,\ angleDot,\ angleDDot)}: Winkel,
  Winkelgeschwindigkeit und optionale Winkelbeschleunigung für
  Matrixableitungen.
\item
  \texttt{identitysptr(size\_t\ n)}: Identitätsmatrix \texttt{n×n};
  \texttt{tildeMatrix(FColDsptr\ col)}: schiefsymmetrische
  Kreuzproduktmatrix des 3D-Vektors.
\item
  \texttt{column(size\_t\ j)}: Spalte \texttt{j};
  \texttt{atijput(i,j,value)} setzt Matrixelement;
  \texttt{atiput(i,row)} ersetzt Zeile \texttt{i}.
\item
  \texttt{timesFullColumn(vector)}, \texttt{timesFullMatrix(matrix)},
  \texttt{timesTransposeFullMatrix(matrix)},
  \texttt{transposeTimesFullMatrix(matrix)}: Produktoperationen;
  Eingabedimensionen müssen kompatibel sein.
\item
  \texttt{times(a)}, \texttt{plusFullMatrix(matrix)},
  \texttt{minusFullMatrix(matrix)}, \texttt{transpose()},
  \texttt{negated()}: Skalierung, Summen/Differenz, Transponieren,
  Negation.
\item
  \texttt{atijplusNumber(i,j,value)},
  \texttt{atijminusNumber(i,j,value)}, \texttt{magnifySelf(factor)},
  \texttt{zeroSelf()}: Matrix in-place ändern; \texttt{sumOfSquares()},
  \texttt{trace()}, \texttt{maxMagnitude()}, \texttt{isDiagonal()}
  liefern Kennwerte.
\item
  \texttt{isDiagonalToWithin(ratio)}, \texttt{equaltol(matrix,\ ratio)},
  \texttt{conditionSelfWithTol(tol)}: numerische
  Verhältnis-/Toleranzparameter für Prüfung/Konditionierung.
\end{itemize}

\texttt{FullRow\textless{}T\textgreater{}}/\texttt{FullColumn\textless{}T\textgreater{}}
implementieren Zeilen-/Spaltenvektoren; \texttt{SparseMatrix},
\texttt{SparseRow}, \texttt{SparseColumn} und \texttt{SparseVector}
speichern nur Nichtnullwerte. Ihre \texttt{at}, \texttt{atiput},
\texttt{atijput}, \texttt{plus}, \texttt{minus}, \texttt{times}-Familien
erhalten Indizes, Werte oder Vektoren; Matrix-/Vektordimensionen sind
Vorbedingung. \texttt{DiagonalMatrix} spezialisiert eine Matrix auf
Diagonaleinträge.

\texttt{MatrixSolver} löst \texttt{A\ x\ =\ b}.
\texttt{solvewithsaveOriginal(matrix,\ rightHandSide,\ saveOriginal)}
nimmt die Koeffizientenmatrix \texttt{A}, rechte Seite \texttt{b} und
ein boolesches Flag, ob A vor der Faktorisierung gesichert wird. Dichte
und Sparse-Überladungen verwenden \texttt{FMatDsptr} beziehungsweise
\texttt{SpMatDsptr}. \texttt{timedSolvewithsaveOriginal} misst
zusätzlich Laufzeit. \texttt{doPivoting(pivotIndex)},
\texttt{forwardEliminateWithPivot(pivotIndex)},
\texttt{backSubstituteIntoDU()} und \texttt{postSolve()} sind
Eliminations-Hooks; \texttt{pivotIndex} ist die aktuelle
Pivotspalte/-zeile. \texttt{singularPivotTolerance} setzt die
Singularitätsschwelle.

Die öffentliche
\texttt{FullMatrix\textless{}T\textgreater{}}-Operationsfamilie umfasst
außerdem vollständig: \texttt{identity()}, \texttt{symLowerWithUpper()},
\texttt{atijputFullColumn(i,j,column)},
\texttt{atijplusFullRow(i,j,row)}, \texttt{atijput(i,row)},
\texttt{atijplusNumber(i,j,value)}, \texttt{atijminusNumber(i,j,value)},
\texttt{copy()}, \texttt{operator+(matrix)},
\texttt{transposeTimesFullColumn(column)}, \texttt{asEulerParameters()},
\texttt{bryantAngles()}, \texttt{eulerAngleszxz()},
\texttt{asDiagonalMatrix()}, \texttt{printOn(stream)} und
Kennwertabfragen \texttt{trace()}, \texttt{sumOfSquares()},
\texttt{maxMagnitude()}, \texttt{isDiagonal()},
\texttt{isDiagonalToWithin(ratio)}, \texttt{equaltol(matrix,ratio)}.
\texttt{rotatex/rotatey/rotatez} erhalten Winkel; \texttt{rotate*rotDot}
zusätzlich Winkelgeschwindigkeit; \texttt{rotate*rotDotrotDDot}
zusätzlich Winkelbeschleunigung. Elementänderungen nehmen
Zeilen-/Spaltenindizes \texttt{i},\texttt{j} und den Wert oder
Zeilen-/Spaltenblock; \texttt{ratio}/\texttt{tol} sind relative
beziehungsweise absolute Vergleichsschwellen.

\texttt{SparseMatrix\textless{}T\textgreater{}} besitzt Konstruktoren
\texttt{(m)}, \texttt{(m,n)} und Initialisierer; \texttt{m/n} sind
Dimensionen. Ihre Methoden: \texttt{atiput(i,row)},
\texttt{atijput(i,j,value)}, \texttt{atijplusNumber(i,j,value)},
\texttt{atijminusNumber(i,j,value)},
\texttt{atijplusDiagonalMatrix(i,j,diag)},
\texttt{atijminusDiagonalMatrix(i,j,diag)},
\texttt{atijplusFullRow(i,j,row)},
\texttt{atijminusFullColumn(i,j,column)},
\texttt{atijplusFullColumn(i,j,column)},
\texttt{atijplusFullMatrix(i,j,matrix)},
\texttt{atijminusFullMatrix(i,j,matrix)},
\texttt{atijplusTransposeFullMatrix(i,j,matrix)},
\texttt{atijminusTransposeFullMatrix(i,j,matrix)},
\texttt{atijplusFullMatrixtimes(i,j,matrix,factor)},
\texttt{timesFullColumn(column)}, \texttt{plusSparseMatrix(matrix)},
\texttt{clonesptr()}, \texttt{magnifySelf(factor)}, \texttt{zeroSelf()},
\texttt{sumOfSquares()}, \texttt{maxMagnitude()},
\texttt{printOn(stream)}. \texttt{i,j} bestimmen die Einfügeposition,
\texttt{factor} skaliert den einzufügenden Matrixblock, die übrigen
Parameter sind der jeweilige Vektor-/Matrixblock. Dimensionen und
Blockposition müssen kompatibel sein.

Die konkreten Zerlegungen heißen \texttt{GEFullMat*}, \texttt{GESpMat*}
(Gauß-Elimination) und \texttt{LDUFullMat*}/\texttt{LDUSpMat*} (LDU);
Suffixe kennzeichnen Voll-/Sparse-Matrix, Pivoting und
Positions-IC-Varianten. Diese Klassen sind Solverintern.

\hypertarget{vektor--array--und-diagonaloperationen}{%
\subsection*{Vektor-, Array- und
Diagonaloperationen}\label{vektor--array--und-diagonaloperationen}}
\addcontentsline{toc}{subsection}{Vektor-, Array- und
Diagonaloperationen}

\texttt{Array\textless{}T\textgreater{}}-Konstruktoren erhalten einen
Vektor, Elementanzahl \texttt{count}, Füllwert oder Wertebereich. Die
allgemeinen Methoden sind \texttt{initialize()},
\texttt{equaltol(x,xx,tol)}, \texttt{copyFrom(array)},
\texttt{zeroSelf()}, \texttt{sumOfSquares()}, \texttt{rootMeanSquare()},
\texttt{numberOfElements()}, \texttt{swapElems(i,ii)},
\texttt{maxMagnitude()}, \texttt{maxMagnitudeOfVector()},
\texttt{equalArrayAt(array,i)}, \texttt{atiput(i,value)},
\texttt{magnifySelf(factor)}, \texttt{negateSelf()},
\texttt{atitimes(i,factor)} und \texttt{printOn(stream)}. \texttt{i/ii}
sind Elementindizes, \texttt{factor} skaliert Werte, \texttt{tol} ist
Vergleichstoleranz. Abstrakte
\texttt{sumOfSquares}/\texttt{maxMagnitude} werden in konkreten
Untertypen bereitgestellt.

\texttt{FullVector\textless{}T\textgreater{}} fügt \texttt{dot(vector)},
\texttt{atiplusNumber(i,value)}, \texttt{atiminusNumber(i,value)},
\texttt{atiplusFullVector(i,vector)},
\texttt{atiplusFullVectortimes(i,vector,factor)},
\texttt{equalSelfPlusFullVectortimes(vector,factor)},
\texttt{normalizeSelf()}, \texttt{length()}, \texttt{conditionSelf()},
\texttt{conditionSelfWithTol(tol)}, \texttt{clonesptr()},
\texttt{isIncreasing()},
\texttt{isIncreasingIfExceptionsAreLessThan(tol)},
\texttt{isDecreasingIfExceptionsAreLessThan(tol)} hinzu. \texttt{dot}
ist Skalarprodukt; \texttt{i} startet das Blockupdate; \texttt{tol} ist
erlaubte Monotonieabweichung beziehungsweise Nullschwelle. Im Header
sind generische \texttt{T}-Spezialisierungen von \texttt{sumOfSquares},
\texttt{zeroSelf}, \texttt{maxMagnitude}, \texttt{conditionSelfWithTol}
und \texttt{clonesptr} teils nicht implementiert; getestete Spezialfälle
gelten vor allem für \texttt{double}.

\texttt{FullColumn\textless{}T\textgreater{}}-Methoden:
\texttt{plusFullColumn(column)}, \texttt{minusFullColumn(column)},
\texttt{times(scalar)}, \texttt{negated()},
\texttt{atiputFullColumn(i,column)},
\texttt{atiplusFullColumn(i,column)},
\texttt{equalSelfPlusFullColumnAt(column,i)},
\texttt{atiminusFullColumn(i,column)},
\texttt{equalFullColumnAt(column,i)}, \texttt{copy()},
\texttt{transpose()}, \texttt{atiplusFullColumntimes(i,column,factor)},
\texttt{transposeTimesFullColumn(column)},
\texttt{equalSelfPlusFullColumntimes(column,factor)},
\texttt{cross(column)}, \texttt{simplified()}, \texttt{dot(vector)} und
\texttt{dot(vectorOfColumns)}. \texttt{cross} verlangt 3D-Vektoren;
Blockoperationen nehmen Startindex \texttt{i}; \texttt{factor} skaliert
den addierten Block. \texttt{simplified()} und einige generische
Skalarüberladungen werfen in der aktuellen Implementierung „To be
implemented``.

\texttt{FullRow\textless{}T\textgreater{}}-Methoden:
\texttt{times(scalar)}, \texttt{negated()}, \texttt{plusFullRow(row)},
\texttt{minusFullRow(row)}, \texttt{timesFullColumn(column)} (shared/raw
pointer), \texttt{timesFullMatrix(matrix)},
\texttt{timesTransposeFullMatrix(matrix)},
\texttt{equalSelfPlusFullRowTimes(row,factor)},
\texttt{equalFullRow(row)}, \texttt{transpose()}, \texttt{copy()},
\texttt{atiplusFullRow(j,row)}, \texttt{transposeTimesFullRow(row)},
\texttt{dot(vector)}, \texttt{dot(vectorOfColumns)},
\texttt{printOn(stream)}. \texttt{j} ist Startspalte des Blockupdates;
Produkte verlangen kompatible Dimensionen.

\texttt{SparseVector\textless{}T\textgreater{}} speichert indexierte
Nichtnullwerte in einer Map. \texttt{SparseVector(n)} setzt die logische
Länge; weitere Konstruktoren erhalten Index/Wert-Paare. Methoden:
\texttt{rootMeanSquare()}, \texttt{numberOfElements()},
\texttt{sumOfSquares()}, \texttt{atiput(i,value)},
\texttt{atiplusNumber(i,value)}, \texttt{atiminusNumber(i,value)},
\texttt{zeroSelf()}, \texttt{maxMagnitude()},
\texttt{magnifySelf(factor)}, \texttt{printOn(stream)}. \texttt{i} ist
Vektorindex, \texttt{n} logische Vektorlänge; Nichtgespeicherte Elemente
gelten als null.

\texttt{DiagonalMatrix\textless{}T\textgreater{}} erhält Dimension
\texttt{count}, Füllwert oder Diagonaleinträge.
\texttt{atiputDiagonalMatrix(i,matrix)} kopiert ab Diagonalindex
\texttt{i}; \texttt{times(factor)}, \texttt{timesFullColumn(column)},
\texttt{timesFullMatrix(matrix)}, \texttt{nrow()}, \texttt{ncol()},
\texttt{sumOfSquares()}, \texttt{numberOfElements()},
\texttt{zeroSelf()}, \texttt{maxMagnitude()}, \texttt{printOn(stream)}
skalieren, multiplizieren oder liefern Kennwerte.
\texttt{numberOfElements()} zählt die äquivalente volle Matrixgröße
\texttt{n*n}; einige generische Typüberladungen sind nicht
implementiert.

\hypertarget{symbolische-funktionen-und-zeitgesetze}{%
\section{\texorpdfstring{{15.11} Symbolische Funktionen und
Zeitgesetze}{15.11 Symbolische Funktionen und Zeitgesetze}}\label{symbolische-funktionen-und-zeitgesetze}}

\texttt{Symbolic}, \texttt{Function}, \texttt{ExpressionX} und
\texttt{FunctionX} bilden Ausdrucksbäume und zeitabhängige Ausdrücke.
\texttt{Function::arguments(Symsptr\ args)} erhält die Argumentliste des
Ausdrucks. \texttt{FunctionFromData}, \texttt{FunctionXY},
\texttt{FunctionWithManyArgs}, \texttt{FunctionXcParameter},
\texttt{UserFunction} und \texttt{BasicUserFunction} repräsentieren
interpolierte, mehrargumentige, parameterisierte oder benutzerdefinierte
Funktionen.

Mathematische Klassen \texttt{Abs}, \texttt{Negative},
\texttt{Reciprocal}, \texttt{Power}, \texttt{Product}, \texttt{Sum},
\texttt{Ln}, \texttt{Log10}, \texttt{LogN}, \texttt{Exponential},
\texttt{Sine}, \texttt{Cosine}, \texttt{ArcSine}, \texttt{ArcTan},
\texttt{ArcTan2} und \texttt{DifferenceOperator} implementieren
Operationen auf symbolischen Argumenten. Parameterrollen folgen dem
Operator: \texttt{Power(base,\ exponent)}, \texttt{LogN(value,\ base)},
\texttt{ArcTan2(y,\ x)}, Summand-/Faktorlisten für
\texttt{Sum}/\texttt{Product}. \texttt{Polynomial},
\texttt{GeneralSpline}, \texttt{DifferentiatedGeneralSpline},
\texttt{PiecewiseFunction}, \texttt{StepFunction} und
\texttt{RampStepFunction} kodieren zeitabhängige Gesetze;
Knoten-/Datenvektoren sind Definitionsdaten, die nicht mit Solverzeit
oder Zeitschritt verwechselt werden dürfen.

\texttt{IndependentVariable}/\texttt{ASMTTime} liefern die Zeitvariable;
\texttt{Integral} und Differenzoperatoren bilden integrierte
beziehungsweise differenzierte Ausdrücke. Funktionswerte, erste und
zweite Ableitungen werden für Motion, Geschwindigkeits- und
Beschleunigungsgleichungen verwendet.

Die vollständige Grundschnittstelle \texttt{Symbolic} umfasst
\texttt{times(arg,arg1)}, \texttt{sum(arg,arg1)}, \texttt{raisedTo(x,y)}
(statische Binäroperatoren mit zwei Ausdruckszeigern),
\texttt{initialize()}, \texttt{differentiateWRT(var)},
\texttt{integrateWRT(var)}, \texttt{simplified()},
\texttt{simplified(sptr)}, \texttt{expandUntil(set)},
\texttt{expandUntil(sptr,set)}, \texttt{simplifyUntil(sptr,set)},
\texttt{isZero()}, \texttt{isOne()}, \texttt{isSum()},
\texttt{isProduct()}, \texttt{isConstant()}, \texttt{isVariable()},
\texttt{getTerms()}, \texttt{addTerm(term)}, \texttt{getValue()},
\texttt{getValue(double\ arg)}, \texttt{setValue(double\ val)},
\texttt{createMbD(system,units)}, \texttt{clonesptr()},
\texttt{sptrConstant(value)},
\texttt{setIntegrationConstant(integConstant)} und
\texttt{printOn(stream)}. \texttt{var} ist die symbolische Variable der
Ableitung/Integration; \texttt{set} ist eine Menge bereits
besuchter/erweiterter Ausdrucksknoten; \texttt{arg} in
\texttt{getValue(arg)} ist der numerische Eingabewert der Funktion;
\texttt{val} setzt einen aktuellen Skalarwert.

\texttt{SymbolicParser} zerlegt ASMT-Ausdrücke.
\texttt{parseString(const\ std::string\&\ expr)} parst einen
Ausdruckstext; \texttt{parseUserFunction(Symsptr\ userFunc)} bindet ihn
an eine Benutzerfunktion; \texttt{initialize()} setzt den Parser zurück.
Parserzustände/Teilregeln ohne Parameter sind \texttt{commaExpression},
\texttt{plusTerm}, \texttt{minusTerm}, \texttt{plainTerm},
\texttt{term}, \texttt{plainFunction}, \texttt{timesFunction},
\texttt{divideByFunction}, \texttt{xLetter}, \texttt{xDigit},
\texttt{xDoubleQuote}, \texttt{symfunction}, \texttt{expression},
\texttt{expressionInParentheses}, \texttt{constant},
\texttt{namedFunction}, \texttt{intrinsic}, \texttt{variable},
\texttt{raisedTo} und \texttt{signedNumber}.
\texttt{peekForTypeNoPush(c)} prüft ein Zeichen ohne Stackänderung;
\texttt{peekForTypevalue(type,symbol)} prüft Tokenart/-wert;
\texttt{scanToken()} liefert das nächste Token; \texttt{expected(msg)}
prüft eine erwartete Syntax;
\texttt{notify(msg)}/\texttt{notifyat(msg,mark)} melden Text
beziehungsweise Zeichenposition; \texttt{combineStackTo(pos)} fasst den
Ausdruckstack bis zum Index \texttt{pos} zusammen;
\texttt{isNextLineTag(char\ c)} erkennt ein Zeilenabschlusszeichen.
\texttt{expr} ist Originaltext, \texttt{Symsptr} der resultierende
Ausdrucksbaum.

\texttt{Arguments::arguments(Symsptr\ args)} setzt die symbolische
Argumentliste; \texttt{Functions} erbt diesen Vertrag ohne zusätzliche
Methoden. \texttt{Function::arguments(args)} ist dieselbe abstrakte
Schnittstelle der Funktionstypen.

\hypertarget{koordinaten-orientierung-und-zustandsdaten}{%
\section{\texorpdfstring{{15.12} Koordinaten, Orientierung und
Zustandsdaten}{15.12 Koordinaten, Orientierung und Zustandsdaten}}\label{koordinaten-orientierung-und-zustandsdaten}}

\texttt{CartesianFrame}, \texttt{PartFrame}, \texttt{MarkerFrame},
\texttt{EndFramec}, \texttt{EndFrameqc}, \texttt{EndFrameqct} und
\texttt{EndFrameqct2} repräsentieren räumliche Bezugssysteme.
\texttt{EulerParameters}, deren Dot/DDot-Varianten,
\texttt{EulerAngles}, \texttt{EulerAngleszxz}, \texttt{Orientation},
\texttt{EulerConstraint}, \texttt{Translation} und \texttt{ZRotation}
verwalten Orientierungs- und Lagekoordinaten.

Parameter in diesen Klassen sind typischerweise ein Frame-/Partzeiger,
Ortsvektor \texttt{r}, Rotationsmatrix \texttt{A},
Eulerparameter/Quaternion \texttt{q}, Winkel \texttt{theta}, deren
zeitliche Ableitungen oder ein Koordinatenindex. \texttt{PosVelAccData}
bündelt Lage, Geschwindigkeit und Beschleunigung für einen Zeitpunkt.
\texttt{Units} enthält Skalenfaktoren für Zeit, Länge, Masse, Winkel und
abgeleitete Größen; Vektoren/Matrizen müssen dieselbe Basis und
Skalierung verwenden.

\hypertarget{fehler--und-diagnosefunktionen}{%
\section{\texorpdfstring{{15.13} Fehler- und
Diagnosefunktionen}{15.13 Fehler- und Diagnosefunktionen}}\label{fehler--und-diagnosefunktionen}}

Fehlerklassen melden unterschiedliche Abbruchursachen:
\texttt{SyntaxError}/\texttt{SimulationStoppingError} (Eingabe oder
Laufabbruch), \texttt{NewtonRaphsonError}/\texttt{MaximumIterationError}
(Konvergenz), \texttt{SingularMatrixError} (lineares System),
\texttt{TooSmallStepSizeError}/\texttt{TooManyTriesError} (Integration),
\texttt{DiscontinuityError} (Sprungstelle) und
\texttt{NotKinematicError} (ungeeigneter Analysemodus).
Konstruktorargumente \texttt{std::string} enthalten in der Regel den
Diagnosekontext; bei \texttt{SingularMatrixError} können zusätzlich
redundante Gleichungsnummern übergeben werden.

Für Clientcode gilt: C++-Exceptions mit \texttt{std::exception}
abfangen. Das reicht nicht, wenn ein interner Pfad den Fehler wie
\texttt{ASMTAssembly::runKINEMATIC()} abfängt und keinen Status
zurückliefert; solche Aufrufe benötigen Ergebnis-/Zeitreihenprüfungen
oder einen anderen Diagnosepfad.

\hypertarget{asmt-unterklassen-zusuxe4tzliche-funktionen}{%
\section{\texorpdfstring{{15.14} ASMT-Unterklassen: zusätzliche
Funktionen}{15.14 ASMT-Unterklassen: zusätzliche Funktionen}}\label{asmt-unterklassen-zusuxe4tzliche-funktionen}}

Die folgenden Methoden ergänzen die geerbten Funktionen der
ASMT-Basisklassen. Parsermethoden konsumieren jeweils den Textrest
\texttt{std::vector\textless{}std::string\textgreater{}\&\ lines};
Writer erhalten \texttt{std::ofstream\&\ os} und \texttt{size\_t\ level}
(Einrückungsebene).

\begin{longtable}[]{@{}
  >{\raggedright\arraybackslash}p{(\columnwidth - 2\tabcolsep) * \real{0.5000}}
  >{\raggedright\arraybackslash}p{(\columnwidth - 2\tabcolsep) * \real{0.5000}}@{}}
\toprule\noalign{}
\begin{minipage}[b]{\linewidth}\raggedright
Klasse/Familie
\end{minipage} & \begin{minipage}[b]{\linewidth}\raggedright
Zusätzliche Methoden und Parameter
\end{minipage} \\
\midrule\noalign{}
\endhead
\bottomrule\noalign{}
\endlastfoot
\texttt{ASMTSpatialItem} & \texttt{readPosition3D(lines)},
\texttt{readRotationMatrix(lines)} parsen Lage und Orientierung;
\texttt{storeOnLevel(os,level)},
\texttt{storeOnLevelPosition(os,level)},
\texttt{storeOnLevelPositionRaw(os,level)},
\texttt{storeOnLevelRotationMatrix(os,level)},
\texttt{storeOnLevelRotationMatrixRaw(os,level)} serialisieren sie;
\texttt{restorePosRot()} setzt alte Lagewerte zurück. \\
\texttt{ASMTSpatialContainer} &
\texttt{readRefPoints/RefPoint/RefCurves/RefCurve/RefSurfaces/RefSurface(lines)}
lesen Referenzobjekte. \texttt{readXs/Ys/Zs},
\texttt{readBryantxs/ys/zs}, \texttt{readVXs/VYs/VZs},
\texttt{readOmegaXs/ys/zs}, \texttt{readAXs/ys/zs},
\texttt{readAlphaXs/ys/zs(lines)} lesen Zeitreihenkomponenten für Lage,
Bryant-Winkel, Linear-/Winkelgeschwindigkeit und -beschleunigung.
\texttt{setVelocity3D(vector)}/\texttt{setOmega3D(vector)} und die
3-Skalar-Überladungen setzen Anfangswerte;
\texttt{readVelocity3D(lines)}/\texttt{readOmega3D(lines)} parsen sie.
\texttt{getVelocity3D/getOmega3D/getAcceleration3D/getAlpha3D(frame)}
erwarten einen Frameindex.
\texttt{updateMbDFromPosition3D(vector\textbar{}x,y,z)} und
\texttt{updateMbDFromRotationMatrix(matrix\textbar{}9\ Werte)}
übertragen Lage/Orientierung ins MbD-Modell. \\
\texttt{ASMTPart} & \texttt{With()} erzeugt ein initialisiertes Part.
\texttt{parseASMT(lines)}, \texttt{readFeatureOrder(lines)},
\texttt{readPrincipalMassMarker(lines)}, \texttt{readPartSeries(lines)}
parsen Daten. \texttt{vOcmO()}/\texttt{omeOpO()} lesen
lineare/Winkelgeschwindigkeit; \texttt{part()} liefert den Partkontext.
\texttt{createMbD(system,units)},
\texttt{preMbDrunDragStep(system,units)},
\texttt{storeOnLevel(os,level)},
\texttt{storeOnLevelMassMarker(os,level)},
\texttt{storeOnTimeSeries(os)} übersetzen, bereiten Drag vor und
serialisieren. \texttt{system} ist Zielsystem, \texttt{units} die
Einheitenskala. \\
\texttt{ASMTRefItem} &
\texttt{addMarker(shared\_ptr\textless{}ASMTMarker\textgreater{}\ marker)}
hängt einen Marker an;
\texttt{readMarkers(lines)}/\texttt{readMarker(lines)} parsen Marker;
\texttt{storeOnLevel(os,level)} schreibt sie. \texttt{ASMTRefPoint},
\texttt{ASMTRefCurve}, \texttt{ASMTRefSurface} erben diese
Funktionen. \\
\texttt{ASMTMarker} & \texttt{With()}, \texttt{parseASMT(lines)},
\texttt{rpmp()}, \texttt{aApm()}, \texttt{createMbD(system,units)},
\texttt{storeOnLevel(os,level)}. \texttt{rpmp} liefert Relativposition,
\texttt{aApm} die Orientierungsabbildung. \\
\texttt{ASMTPrincipalMassMarker} & \texttt{With()},
\texttt{parseASMT(lines)}, \texttt{setMass(double\ mass)},
\texttt{setDensity(double\ density)},
\texttt{setMomentOfInertias(DiagMatDsptr\ matrix)} oder drei
Diagonalwerte, \texttt{storeOnLevel(os,level)}. Masse/Dichte sind
Skalare, Trägheit ist Diagonalmatrix oder drei Hauptträgheitsmomente. \\
\texttt{ASMTItemIJ} & \texttt{initialize()},
\texttt{setMarkerI/J(path)}, \texttt{readMarkerI/J(lines)},
\texttt{readFXonIs/FYonIs/FZonIs/TXonIs/TYonIs/TZonIs(lines)},
\texttt{parseASMT(lines)}, \texttt{createMbD(system,units)},
\texttt{mbdClassNew()}, \texttt{storeOnLevel(os,level)},
\texttt{storeOnTimeSeries(os)}. Kraft-/Momentleser füllen Zeitreihen auf
Seite I; Pfade identifizieren I/J. \\
\texttt{ASMTJoint} & \texttt{With()}, \texttt{readJointSeries(lines)},
\texttt{storeOnLevel(os,level)}, \texttt{storeOnTimeSeries(os)},
\texttt{createMbD(system,units)}. Die Zeitreihenmethode
liest/serialisiert Jointreaktionen. \\
Joint-Fabriken & Parameterlose \texttt{With()}-Fabriken existieren für
\texttt{ASMTAngleJoint}, \texttt{ASMTGearJoint},
\texttt{ASMTNoRotationJoint}, \texttt{ASMTParallelAxesJoint},
\texttt{ASMTPerpendicularJoint}, \texttt{ASMTRackPinionJoint},
\texttt{ASMTScrewJoint}, \texttt{ASMTConstantVelocityJoint},
\texttt{ASMTFixedJoint}, \texttt{ASMTRevoluteJoint},
\texttt{ASMTSphericalJoint}, \texttt{ASMTUniversalJoint},
\texttt{ASMTSphSphJoint}, \texttt{ASMTCylSphJoint},
\texttt{ASMTRevCylJoint}, \texttt{ASMTRevRevJoint},
\texttt{ASMTCylindricalJoint}, \texttt{ASMTPointInLineJoint},
\texttt{ASMTTranslationalJoint}, \texttt{ASMTLineInPlaneJoint},
\texttt{ASMTPlanarJoint} und \texttt{ASMTPointInPlaneJoint}. Sie erben
\texttt{setMarkerI/J}; spezielle Übersetzungsfabriken
\texttt{mbdClassNew()} haben keine Argumente. Typspezifische Werte
liegen in den jeweiligen Klassenfeldern. \\
\texttt{ASMTAngleJoint} & \texttt{With()}, \texttt{mbdClassNew()},
\texttt{parseASMT(lines)}, \texttt{readTheIzJz(lines)},
\texttt{createMbD(system,units)}, \texttt{storeOnLevel(os,level)}.
\texttt{theIzJz} ist der relative Winkel um z (Radiant). \\
\texttt{ASMTGearJoint} & \texttt{With()}, \texttt{mbdClassNew()},
\texttt{parseASMT(lines)}, \texttt{readMarkerK(lines)},
\texttt{readRadiusI(lines)}, \texttt{readRadiusJ(lines)},
\texttt{setMarkerK(path)}, \texttt{createMbD(system,units)},
\texttt{storeOnLevel(os,level)}. \texttt{markerK} ist zusätzlicher
Markerpfad; \texttt{radiusI}/\texttt{radiusJ} sind die beiden wirksamen
Radien; \texttt{aConstant} ist der gespeicherte Kopplungsoffset. \\
\texttt{ASMTScrewJoint} & \texttt{With()}, \texttt{mbdClassNew()},
\texttt{parseASMT(lines)}, \texttt{readPitch(lines)},
\texttt{createMbD(system,units)}, \texttt{storeOnLevel(os,level)}.
\texttt{pitch} ist die Schraubensteigung; \texttt{aConstant} der
Kopplungsoffset. \\
\texttt{ASMTRackPinionJoint} & \texttt{With()}, \texttt{mbdClassNew()},
\texttt{parseASMT(lines)}, \texttt{readPitchRadius(lines)},
\texttt{createMbD(system,units)}, \texttt{storeOnLevel(os,level)}.
\texttt{pitchRadius} ist der wirksame Ritzelradius; \texttt{aConstant}
der Kopplungsoffset. \\
\texttt{ASMTMotion} & \texttt{readMotionSeries(lines)},
\texttt{initMarkers()}, \texttt{storeOnLevel(os,level)},
\texttt{storeOnTimeSeries(os)}, \texttt{createMbD(system,units)}. \\
\texttt{ASMTRotationalMotion} & \texttt{With()},
\texttt{parseASMT(lines)}, \texttt{readMotionJoint/RotationZ(lines)},
\texttt{initMarkers()}, \texttt{createMbD(system,units)},
\texttt{mbdClassNew()},
\texttt{setMotionJoint(std::string\ motionJoint)},
\texttt{setRotationZ(std::string\ rotZ)},
\texttt{storeOnLevel(os,level)}, \texttt{storeOnTimeSeries(os)}.
\texttt{motionJoint} ist der Jointpfad; \texttt{rotZ} ist ein Ausdruck
für die Rotation um z. \\
\texttt{ASMTAllowRotation} & \texttt{With()}, \texttt{parseASMT(lines)},
\texttt{readMotionJoint(lines)}, \texttt{initMarkers()},
\texttt{mbdClassNew()},
\texttt{setMotionJoint(std::string\ motionJoint)},
\texttt{storeOnLevel(os,level)}, \texttt{storeOnTimeSeries(os)}. Der
Jointpfad benennt den Bezug; \texttt{rotationZ} ist der zugehörige
gespeicherte Z-Rotationsausdruck. \\
\texttt{ASMTTranslationalMotion} & \texttt{parseASMT(lines)},
\texttt{initMarkers()}, \texttt{createMbD(system,units)},
\texttt{mbdClassNew()}, \texttt{readMotionJoint(lines)},
\texttt{setTranslationZ(std::string\ tranZ)},
\texttt{readTranslationZ(lines)}, \texttt{storeOnLevel(os,level)},
\texttt{storeOnTimeSeries(os)}. \texttt{tranZ} ist der ASMT-Ausdruck für
die Z-Translation. \\
\texttt{ASMTGeneralMotion} & \texttt{With()}, \texttt{parseASMT(lines)},
\texttt{readrIJI/angIJJ/RotationOrder(lines)}, \texttt{mbdClassNew()},
\texttt{createMbD(system,units)}, \texttt{storeOnLevel(os,level)},
\texttt{storeOnTimeSeries(os)}. \texttt{rIJI} enthält drei
translatorische Ausdrücke, \texttt{angIJJ} drei Winkelgesetze;
\texttt{rotationOrder} gibt deren Reihenfolge an. \\
\texttt{ASMTLimit} & \texttt{initMarkers()},
\texttt{readMotionJoint/Limit/Type/Tol(lines)},
\texttt{parseASMT(lines)}, \texttt{createMbD(system,units)},
\texttt{setmotionJoint(string)}, \texttt{settype(string)},
\texttt{setlimit(string)}, \texttt{settol(string)},
\texttt{storeOnLevel(os,level)}. Diese Setterparameter sind Strings,
nicht \texttt{double}. \\
\texttt{ASMTForceTorque} & \texttt{updateFromMbD()},
\texttt{compareResults(AnalysisType\ type)},
\texttt{outputResults(type)} ergänzen die geerbten
I/J-Kraft-/Momentdaten; \texttt{type} bestimmt die Analyseart. \\
\texttt{ASMTConstantGravity} & \texttt{With()},
\texttt{parseASMT(lines)}, \texttt{createMbD(system,units)},
\texttt{getg()}, \texttt{setg(vector)}/\texttt{setg(x,y,z)},
\texttt{storeOnLevel(os,level)}. \\
\texttt{ASMTSimulationParameters} & \texttt{With()},
\texttt{parseASMT(lines)}, die in
\protect\hyperlink{abs-api-motion-limits-lasten}{Abschnitt 15.6}
gelisteten Setter und \texttt{storeOnLevel(os,level)}. \\
\texttt{ASMTAnimationParameters} & \texttt{With()},
\texttt{parseASMT(lines)}, \texttt{storeOnLevel(os,level)}. Es gibt
keine speziellen Setter; \texttt{nframe}, \texttt{icurrent},
\texttt{istart}, \texttt{iend}, \texttt{framesPerSecond} sind
Frameanzahl/Indizes beziehungsweise Bildrate, \texttt{isForward} ist ein
Richtungsflag. \\
\texttt{ASMTTime} & \texttt{With()}, Konstruktoren
\texttt{ASMTTime()}/\texttt{ASMTTime(Symsptr\ arg)},
\texttt{deleteMbD()}, \texttt{createMbD(system,units)},
\texttt{expandUntil(expr,set)}, \texttt{simplifyUntil(expr,set)},
\texttt{copyWith(arg)}, \texttt{isVariable()},
\texttt{setValue(double\ val)}. \texttt{arg}/\texttt{expr} sind
symbolische Ausdrücke, \texttt{set} besuchte Knoten, \texttt{val} der
Zeitwert. \\
\end{longtable}

\hypertarget{vollstuxe4ndiger-klassenindex-nach-funktionsfamilie}{%
\section{\texorpdfstring{{15.15} Vollständiger Klassenindex nach
Funktionsfamilie}{15.15 Vollständiger Klassenindex nach Funktionsfamilie}}\label{vollstuxe4ndiger-klassenindex-nach-funktionsfamilie}}

Dieser Index ordnet alle fachlichen Klassenfamilien des eingebundenen
\texttt{OndselSolver/OndselSolver}-Headerverzeichnisses ein. Die
konkreten Gleichungsunterklassen erben ihre Aufrufverträge; ihre
zusätzlichen Aufgaben bestehen überwiegend darin, Constraintwert,
Ableitungen und Jacobian für eine andere geometrische
Koordinatenkombination zu liefern.

\begin{longtable}[]{@{}
  >{\raggedright\arraybackslash}p{(\columnwidth - 4\tabcolsep) * \real{0.3333}}
  >{\raggedright\arraybackslash}p{(\columnwidth - 4\tabcolsep) * \real{0.3333}}
  >{\raggedright\arraybackslash}p{(\columnwidth - 4\tabcolsep) * \real{0.3333}}@{}}
\toprule\noalign{}
\begin{minipage}[b]{\linewidth}\raggedright
Funktionsfamilie
\end{minipage} & \begin{minipage}[b]{\linewidth}\raggedright
Zugeordnete Klassen/Varianten
\end{minipage} & \begin{minipage}[b]{\linewidth}\raggedright
Gemeinsame Parameter und Funktion
\end{minipage} \\
\midrule\noalign{}
\endhead
\bottomrule\noalign{}
\endlastfoot
ASMT-Root und Parser & \texttt{ASMTAssembly}, \texttt{ASMTItem},
\texttt{ASMTItemIJ}, \texttt{ASMTSpatialItem},
\texttt{ASMTSpatialContainer}, \texttt{ASMTPart}, \texttt{ASMTMarker},
\texttt{ASMTRefItem}, \texttt{ASMTRefPoint}, \texttt{ASMTRefCurve},
\texttt{ASMTRefSurface}, \texttt{ASMTPrincipalMassMarker},
\texttt{ASMTConstraintSet}, \texttt{ASMTKinematicIJ},
\texttt{ASMTExtrusion} & Parserzeilen, Objekt-/Pfadnamen, Writerstream
und Ebene; ausführliche Signaturen in 16.2--16.5 und 16.14.
\texttt{ASMTExtrusion} ist deklariert, aber die
Feature-Order-Parserbehandlung ist noch auskommentiert. \\
ASMT-Joints & \texttt{ASMTAngleJoint}, \texttt{ASMTAtPointJoint},
\texttt{ASMTCompoundJoint}, \texttt{ASMTConstantVelocityJoint},
\texttt{ASMTCylindricalJoint}, \texttt{ASMTCylSphJoint},
\texttt{ASMTFixedJoint}, \texttt{ASMTGearJoint},
\texttt{ASMTInLineJoint}, \texttt{ASMTInPlaneJoint}, \texttt{ASMTJoint},
\texttt{ASMTLineInPlaneJoint}, \texttt{ASMTNoRotationJoint},
\texttt{ASMTParallelAxesJoint}, \texttt{ASMTPerpendicularJoint},
\texttt{ASMTPlanarJoint}, \texttt{ASMTPointInLineJoint},
\texttt{ASMTPointInPlaneJoint}, \texttt{ASMTRackPinionJoint},
\texttt{ASMTRevCylJoint}, \texttt{ASMTRevRevJoint},
\texttt{ASMTRevoluteJoint}, \texttt{ASMTScrewJoint},
\texttt{ASMTSphSphJoint}, \texttt{ASMTSphericalJoint},
\texttt{ASMTTranslationalJoint}, \texttt{ASMTUniversalJoint} & Fabriken
\texttt{With()}, Markerpfade I/J und ggf. K,
\texttt{createMbD(system,units)}, Parser-/Writerhooks; Sonderparameter
Winkel, Radien, Steigung. \\
ASMT-Motion/Limit/Last & \texttt{ASMTAllowRotation},
\texttt{ASMTRotationalMotion}, \texttt{ASMTTranslationalMotion},
\texttt{ASMTGeneralMotion}, \texttt{ASMTMotion}, \texttt{ASMTLimit},
\texttt{ASMTRotationLimit}, \texttt{ASMTTranslationLimit},
\texttt{ASMTForceTorque}, \texttt{ASMTConstantGravity},
\texttt{ASMTSimulationParameters}, \texttt{ASMTAnimationParameters},
\texttt{ASMTTime}, \texttt{ASMTContact}, \texttt{ASMTConeConeContact},
\texttt{ASMTCylConeContact}, \texttt{ASMTCylCylContact} &
Motion-/Jointpfade, Bewegungsausdrücke, Grenzstrings,
Kraft-/Momentreihen, Zeit-/Schrittparameter und Einheiten; Details in
16.6/16.14. Kontaktklassen stehen im Headerbaum, sind aber nicht Teil
des gezeigten Joint-Parserdispatchs. \\
MbD-Baum/Zustand & \texttt{Item}, \texttt{System},
\texttt{ExternalSystem}, \texttt{StateData}, \texttt{PosVelAccData},
\texttt{Part}, \texttt{PartFrame}, \texttt{MarkerFrame},
\texttt{CartesianFrame}, \texttt{KinematicIeJe}, \texttt{EndFramec},
\texttt{EndFrameqc}, \texttt{EndFrameqct}, \texttt{EndFrameqct2},
\texttt{Constraint}, \texttt{ConstraintIJ}, \texttt{ConstraintSet},
\texttt{RedundantConstraint}, \texttt{AbsConstraint}, \texttt{LimitIJ},
\texttt{ItemIJ} & System-/Objektzeiger, Zustandsvektoren, Framebezug,
Gleichungsnummern und Ereigniszustände; \texttt{fill*} füllt
Ergebnisvektoren/-matrizen. \\
Elementare/zusammengesetzte Joints & \texttt{Joint},
\texttt{AtPointJoint}, \texttt{FixedJoint}, \texttt{RevoluteJoint},
\texttt{SphericalJoint}, \texttt{CylindricalJoint},
\texttt{TranslationalJoint}, \texttt{UniversalJoint},
\texttt{PlanarJoint}, \texttt{InLineJoint}, \texttt{InPlaneJoint},
\texttt{AngleJoint}, \texttt{CompoundJoint}, \texttt{SphSphJoint},
\texttt{CylSphJoint}, \texttt{RevCylJoint}, \texttt{RevRevJoint},
\texttt{NoRotationJoint}, \texttt{ParallelAxesJoint},
\texttt{PerpendicularJoint}, \texttt{GearJoint}, \texttt{RackPinJoint},
\texttt{ScrewJoint}, \texttt{ConstantVelocityJoint} & Endframes I/J;
Constraints werden erzeugt und Reaktionskräfte/-momente aggregiert.
Gear-/Screw-Parameter sind Verhältnis/Radien/Steigung. \\
Punkt-/Abstandsbedingungen & \texttt{AtPointConstraintIJ},
\texttt{AtPointConstraintIqcJc}, \texttt{AtPointConstraintIqcJqc},
\texttt{AtPointConstraintIqctJqc}, \texttt{DistanceConstraintIJ},
\texttt{DistanceConstraintIqcJc}, \texttt{DistanceConstraintIqcJqc},
\texttt{DistanceConstraintIqctJqc}, \texttt{DistancexyConstraintIJ},
\texttt{DistancexyConstraintIqcJc}, \texttt{DistancexyConstraintIqcJqc},
\texttt{DistancexyConstraintIqctJqc}, \texttt{DistIecJec},
\texttt{DistIeqcJec}, \texttt{DistIeqcJeqc}, \texttt{DistIeqctJeqc},
\texttt{DistxyIecJec}, \texttt{DistxyIeqcJec}, \texttt{DistxyIeqcJeqc},
\texttt{DistxyIeqctJeqc} & Körper-/Framekoordinaten I/J; Zielabstand
oder Komponenten; Funktionen füllen Residuum, Zeitableitung und
Jacobian. \texttt{xy} beschränkt eine Projektion statt vollen
3D-Abstand. \\
Orientierung/Winkel & \texttt{DirectionCosineConstraintIJ},
\texttt{DirectionCosineConstraintIqcJc},
\texttt{DirectionCosineConstraintIqcJqc},
\texttt{DirectionCosineConstraintIqctJqc},
\texttt{DirectionCosineIecJec}, \texttt{DirectionCosineIeqcJec},
\texttt{DirectionCosineIeqcJeqc}, \texttt{DirectionCosineIeqctJeqc},
\texttt{AngleZConstraintIJ}, \texttt{AngleZConstraintIqcJc},
\texttt{AngleZConstraintIqcJqc}, \texttt{AngleZIecJec},
\texttt{AngleZIeqcJec}, \texttt{AngleZIeqcJeqc},
\texttt{OrbitAngleZIecJec}, \texttt{OrbitAngleZIeqcJec},
\texttt{OrbitAngleZIeqcJeqc}, \texttt{EulerConstraint},
\texttt{AllowZRotation}, \texttt{AllowZRotationConstraintIqctJqc} &
Orientierungs-/Richtungsvektoren, Komponentenachsen und gegebenenfalls
Zielwinkel; Winkel sind im Solvermaß Radiant. Ergebnis ist
Winkel-/Richtungsgleichung samt Ableitungen. \\
Translation, Getriebe und Schraube & \texttt{TranslationConstraintIJ},
\texttt{TranslationConstraintIqcJc},
\texttt{TranslationConstraintIqcJqc},
\texttt{TranslationConstraintIqctJqc}, \texttt{TranslationLimitIJ},
\texttt{RotationLimitIJ}, \texttt{LimitIJ}, \texttt{GearConstraintIJ},
\texttt{GearConstraintIqcJc}, \texttt{GearConstraintIqcJqc},
\texttt{CoaxialGearConstraintIJ}, \texttt{RackPinConstraintIJ},
\texttt{RackPinConstraintIqcJc}, \texttt{RackPinConstraintIqcJqc},
\texttt{ScrewConstraintIJ}, \texttt{ScrewConstraintIqcJc},
\texttt{ScrewConstraintIqcJqc}, \texttt{ConstVelConstraintIJ},
\texttt{ConstVelConstraintIqcJc}, \texttt{ConstVelConstraintIqcJqc} &
Weg-/Winkelkoordinate, Grenzwert, Radius/Steigung beziehungsweise
Übersetzungsverhältnis; konkrete Constraints liefern die jeweilige
gekoppelte Gleichung. \\
Komponenten-/Verschiebungsbedingungen & \texttt{DispCompIecJecIe},
\texttt{DispCompIecJecKec}, \texttt{DispCompIecJecKeqc},
\texttt{DispCompIecJecO}, \texttt{DispCompIeqcJecIe},
\texttt{DispCompIeqcJecKeqc}, \texttt{DispCompIeqcJecO},
\texttt{DispCompIeqcJeqcIe}, \texttt{DispCompIeqcJeqcKeqc},
\texttt{DispCompIeqcJeqcKeqct}, \texttt{DispCompIeqcJeqcO},
\texttt{DispCompIeqctJeqcIe}, \texttt{DispCompIeqctJeqcKeqct},
\texttt{DispCompIeqctJeqcO} & Frame-/Achspunkte I/J/K und
Projektionsrichtung; Parameter bestimmen, welche Verschiebungskomponente
als skalares Residuum ausgewertet wird. \\
Solversteuerung & \texttt{Solver}, \texttt{SystemSolver},
\texttt{SystemNewtonRaphson}, \texttt{NewtonRaphson},
\texttt{ScalarNewtonRaphson}, \texttt{VectorNewtonRaphson},
\texttt{PosNewtonRaphson}, \texttt{PosKineNewtonRaphson},
\texttt{PosICNewtonRaphson}, \texttt{PosICKineNewtonRaphson},
\texttt{PosICDragNewtonRaphson}, \texttt{PosICDragLimitNewtonRaphson},
\texttt{AnyPosICNewtonRaphson}, \texttt{VelSolver},
\texttt{VelKineSolver}, \texttt{VelICSolver}, \texttt{VelICKineSolver},
\texttt{AccNewtonRaphson}, \texttt{AccKineNewtonRaphson},
\texttt{AccICNewtonRaphson}, \texttt{AccICKineNewtonRaphson} &
Systemzeiger, Zustands-/Korrekturvektor, Residuum/Jacobian, Toleranz und
Iterations-/Backtrackingzähler. Die Unterschiede sind Analyseart und
Anfangs-/Dragzustand. \\
Integrationsverfahren & \texttt{Integrator},
\texttt{IntegratorInterface}, \texttt{BasicIntegrator},
\texttt{BasicQuasiIntegrator}, \texttt{KineIntegrator},
\texttt{ICKineIntegrator}, \texttt{QuasiIntegrator},
\texttt{LinearMultiStepMethod}, \texttt{StableBackwardDifference} &
System, Zeitrichtung, Zeitintervall, Schrittweiten, Ordnung und
Fehlertoleranzen; Hooks sind im
\protect\hyperlink{abs-api-zeitintegration}{Abschnitt 15.9}
aufgeführt. \\
Dichte/Sparse-Algebra und Zerlegung & \texttt{Array},
\texttt{FullVector}, \texttt{FullRow}, \texttt{FullColumn},
\texttt{FullMatrix}, \texttt{SparseVector}, \texttt{SparseRow},
\texttt{SparseColumn}, \texttt{SparseMatrix}, \texttt{DiagonalMatrix},
\texttt{RowTypeMatrix}, \texttt{MatrixSolver},
\texttt{MatrixDecomposition}, \texttt{MatrixGaussElimination},
\texttt{MatrixLDU}, \texttt{GEFullMat}, \texttt{GEFullMatFullPv},
\texttt{GEFullMatParPv}, \texttt{GESpMat}, \texttt{GESpMatFullPv},
\texttt{GESpMatFullPvPosIC}, \texttt{GESpMatParPv},
\texttt{GESpMatParPvMarko}, \texttt{GESpMatParPvMarkoFast},
\texttt{GESpMatParPvPrecise}, \texttt{LDUFullMat},
\texttt{LDUFullMatParPv}, \texttt{LDUSpMat}, \texttt{LDUSpMatParPv},
\texttt{LDUSpMatParPvMarko}, \texttt{LDUSpMatParPvPrecise},
\texttt{EigenDecomposition} & Matrix-/Vektordimensionen,
Zeilen-/Spalten-/Pivotindizes, RHS, Toleranz und optionale
Blockfaktoren; Details zu Kernoperationen in 16.10. \\
Symbolik, Funktionen und Zeitgesetze & \texttt{MbDMath},
\texttt{Numeric}, \texttt{Symbolic}, \texttt{SymbolicParser},
\texttt{ExpressionX}, \texttt{Function}, \texttt{FunctionX},
\texttt{FunctionXY}, \texttt{FunctionFromData},
\texttt{FunctionWithManyArgs}, \texttt{FunctionXcParameter},
\texttt{UserFunction}, \texttt{BasicUserFunction}, \texttt{Constant},
\texttt{Variable}, \texttt{IndependentVariable}, \texttt{Functions},
\texttt{Abs}, \texttt{Negative}, \texttt{Reciprocal}, \texttt{Power},
\texttt{Product}, \texttt{Sum}, \texttt{DifferenceOperator},
\texttt{Exponential}, \texttt{Ln}, \texttt{Log10}, \texttt{LogN},
\texttt{Sine}, \texttt{Cosine}, \texttt{ArcSine}, \texttt{ArcTan},
\texttt{ArcTan2}, \texttt{Integral}, \texttt{Polynomial},
\texttt{GeneralSpline}, \texttt{AnyGeneralSpline},
\texttt{DifferentiatedGeneralSpline}, \texttt{PiecewiseFunction},
\texttt{StepFunction}, \texttt{RampStepFunction}, \texttt{Transitions},
\texttt{SymTime} & Symbolische Operanden, Variable, Argumentliste,
Zeitwert, Daten-/Knotenvektoren und Integrationskonstante;
Auswertung/Differentiation/Integration wie in 16.11. \\
Orientierungsdarstellungen & \texttt{EulerArray}, \texttt{EulerAngles},
\texttt{EulerAnglesDot}, \texttt{EulerAnglesDDot},
\texttt{EulerAngleszxz}, \texttt{EulerAngleszxzDot},
\texttt{EulerAngleszxzDDot}, \texttt{EulerParameters},
\texttt{EulerParametersDot}, \texttt{EulerParametersDDot},
\texttt{Orientation}, \texttt{Translation}, \texttt{ZRotation},
\texttt{ZTranslation} & Winkel/Quaternionen, erste/zweite Zeitableitung,
Rotationsfolge und Frame; Quaternionen benötigen konsistente
Normierung. \\
Fehler, Belastungs- und Nebenklassen & \texttt{NewtonRaphsonError},
\texttt{MaximumIterationError}, \texttt{SingularMatrixError},
\texttt{TooManyTriesError}, \texttt{TooSmallStepSizeError},
\texttt{DiscontinuityError}, \texttt{SimulationStoppingError},
\texttt{SyntaxError}, \texttt{NotKinematicError},
\texttt{ConstantGravity}, \texttt{ForceTorqueItem},
\texttt{ForceTorqueData}, \texttt{PrescribedMotion},
\texttt{FullMotion}, \texttt{Arguments}, \texttt{CREATE},
\texttt{CADSystem}, \texttt{MomentOfInertiaSolver}, \texttt{ZRotation},
\texttt{ZTranslation} & Fehlertext/Fehlerkontext,
System-/Partreferenzen, Motionfunktionen oder numerische
Argumentcontainer; Detailparameter ergeben sich aus den jeweiligen
Headern. \\
\end{longtable}

Für alle Constraint-Unterklassen gilt der in 16.5 beschriebene
gemeinsame Vertrag aus \texttt{fill*}, \texttt{set*},
\texttt{isRedundant}, \texttt{type} und \texttt{constraintSpec}; für
Solverunterklassen gelten die Lifecycle-Hooks aus 16.8/16.9. Der
Klassenindex stellt sicher, dass auch die vielen Variantenfamilien
zugeordnet sind, ohne identische geerbte Signaturen dutzendfach zu
wiederholen.

\hypertarget{parameterkurzreferenz}{%
\section{\texorpdfstring{{15.16}
Parameterkurzreferenz}{15.16 Parameterkurzreferenz}}\label{parameterkurzreferenz}}

\begin{longtable}[]{@{}
  >{\raggedright\arraybackslash}p{(\columnwidth - 4\tabcolsep) * \real{0.3333}}
  >{\raggedright\arraybackslash}p{(\columnwidth - 4\tabcolsep) * \real{0.3333}}
  >{\raggedright\arraybackslash}p{(\columnwidth - 4\tabcolsep) * \real{0.3333}}@{}}
\toprule\noalign{}
\begin{minipage}[b]{\linewidth}\raggedright
Parameterklasse
\end{minipage} & \begin{minipage}[b]{\linewidth}\raggedright
Typische Namen
\end{minipage} & \begin{minipage}[b]{\linewidth}\raggedright
Bedeutung
\end{minipage} \\
\midrule\noalign{}
\endhead
\bottomrule\noalign{}
\endlastfoot
Datei/Parser & \texttt{fileName}, \texttt{filename}, \texttt{lines} &
Eingabe-/Ausgabepfad; \texttt{lines} ist der noch nicht konsumierte
ASMT-Text. \\
Baum/Objekt & \texttt{name}, \texttt{partialName}, \texttt{longname},
\texttt{path}, \texttt{owner} & Objektidentität oder voller ASMT-Pfad;
Owner ist Elternknoten. \\
Solververknüpfung & \texttt{mbdSys}, \texttt{system}, \texttt{sys},
\texttt{mbdUnits} & Zielsystem/-solver und Einheitenkonfiguration für
Erzeugung/Integration. \\
Geometrie & \texttt{position3D}, \texttt{rotationMatrix},
\texttt{markerI}, \texttt{markerJ} & Ortsvektor, Orientierung und die
zwei Verbindungsreferenzen. \\
Zeit & \texttt{t}, \texttt{tstart}, \texttt{tend}, \texttt{hmin},
\texttt{hmax}, \texttt{hout}, \texttt{direction} & Zeit, Intervall,
Integrationsschrittgrenzen, Ausgabeabstand und Laufrichtung. \\
Numerik & \texttt{errorTol}, \texttt{tol}, \texttt{ratio},
\texttt{maxIter}, \texttt{orderMax} & Konvergenz-/Vergleichstoleranz,
Iterationslimit und Integrationsordnung. \\
Algebra & \texttt{matrix}, \texttt{mat}, \texttt{fullMat},
\texttt{spMat}, \texttt{vector}, \texttt{column}, \texttt{row},
\texttt{rhs} & Koeffizienten-/Sparse-Matrix, Zustands-/Fehlervektor oder
rechte Seite des Gleichungssystems. \\
Ablauf/Index & \texttt{index}, \texttt{i}, \texttt{j}, \texttt{frame},
\texttt{level}, \texttt{pivot} & Objekt-/Zeitframeindex,
Matrixzeile/-spalte, ASMT-Einrückung oder Pivotposition. Vor
Indexzugriff Dimension und gültigen Bereich prüfen. \\
Steuerung & \texttt{bool}-Parameter wie \texttt{saveOriginal},
\texttt{todebug} & Option ein-/ausschalten; Bedeutung steht beim
jeweiligen Aufruf. \\
\end{longtable}

Diese Referenz beschreibt die Parameter nach ihrer Verwendung im Solver.
Wo OndselSolver einen generischen Namen wie \texttt{a}, \texttt{b} oder
\texttt{c} verwendet, ergibt sich die genaue Bedeutung aus der konkreten
Methode und nicht allein aus dem Parameternamen.

\hypertarget{fabriken-einheiten-und-numerische-hilfen}{%
\section{\texorpdfstring{{15.17} Fabriken, Einheiten und numerische
Hilfen}{15.17 Fabriken, Einheiten und numerische Hilfen}}\label{fabriken-einheiten-und-numerische-hilfen}}

\hypertarget{generische-fabrik-createt}{%
\subsection*{\texorpdfstring{Generische Fabrik
\texttt{CREATE\textless{}T\textgreater{}}}{Generische Fabrik CREATE\textless T\textgreater{}}}\label{generische-fabrik-createt}}
\addcontentsline{toc}{subsection}{Generische Fabrik
\texttt{CREATE\textless{}T\textgreater{}}}

\texttt{CREATE\textless{}T\textgreater{}::With(...)} erzeugt eine
Instanz, ruft \texttt{initialize()} auf und liefert
\texttt{shared\_ptr\textless{}T\textgreater{}}. Die öffentlichen
Überladungen sind:

\begin{longtable}[]{@{}
  >{\raggedright\arraybackslash}p{(\columnwidth - 2\tabcolsep) * \real{0.5000}}
  >{\raggedright\arraybackslash}p{(\columnwidth - 2\tabcolsep) * \real{0.5000}}@{}}
\toprule\noalign{}
\begin{minipage}[b]{\linewidth}\raggedright
Überladung
\end{minipage} & \begin{minipage}[b]{\linewidth}\raggedright
Parameterbedeutung
\end{minipage} \\
\midrule\noalign{}
\endhead
\bottomrule\noalign{}
\endlastfoot
\texttt{With()} & Standardkonstruktor. \\
\texttt{With(const\ std::string\&\ name)} & Objektname/Text für den
Konstruktor. \\
\texttt{With(const\ std::string\&\ expr,\ double\ unit)} & Ausdruckstext
und zugehörige Einheitenskala. \\
\texttt{With(size\_t\ n)} & Ein-Dimension-Größe, z. B. Vektorlänge. \\
\texttt{With(size\_t\ m,\ size\_t\ n)} & Zeilen-/Spaltenzahl einer
Matrix. \\
\texttt{With(std::initializer\_list\textless{}double\textgreater{}\ values)}
& Anfangswerte eines numerischen Containers. \\
\texttt{With(frameI,\ frameJ)} & Zwei Endframes für eine Verbindung. \\
\texttt{With(frameI,\ frameJ,\ size\_t\ axis)} & Zwei Endframes und eine
Achse. \\
\texttt{With(frameI,\ frameJ,\ frameK,\ size\_t\ axisK)} & Drei
Endframes und Achse am dritten Frame. \\
\texttt{With(frameI,\ frameJ,\ size\_t\ axisI,\ size\_t\ axisJ)} & Zwei
Endframes und je eine lokale Achse. \\
\texttt{ConstraintWith(frameI,frameJ,axis)} /
\texttt{ConstraintWith(frameI,frameJ,axisI,axisJ)} & Erzeugt ein
passendes Constraint für Frames und Achse(n); Auswahl hängt im aktuellen
Code vom konkreten Template-Typ und Frame-Untertyp ab. \\
\end{longtable}

\texttt{frameI/J/K} sind
\texttt{shared\_ptr\textless{}EndFramec\textgreater{}}; Achsindizes
wählen lokale Komponenten. Da \texttt{CREATE} für die konkrete
Unterklasse die passende Konstruktorüberladung benötigt, kann nicht
jeder Typ mit jeder \texttt{With}-Signatur erzeugt werden.

\hypertarget{einheiten-und-skalarhilfen}{%
\subsection*{Einheiten und
Skalarhilfen}\label{einheiten-und-skalarhilfen}}
\addcontentsline{toc}{subsection}{Einheiten und Skalarhilfen}

\texttt{Units()} setzt Einheitsfaktoren auf Standardwerte.
\texttt{Units(double\ unitTime,\ double\ unitMass,\ double\ unitLength,\ double\ unitAngle)}
erhält Zeit-, Masse-, Längen- und Winkelskala und ruft \texttt{calc()}
auf. \texttt{initialize()} ist parameterlos. \texttt{calc()} leitet
daraus ab: \texttt{aJ\ =\ mass*length\^{}2},
\texttt{velocity\ =\ length/time}, \texttt{omega\ =\ angle/time},
\texttt{acceleration\ =\ velocity/time}, \texttt{alpha\ =\ omega/time},
\texttt{force\ =\ mass*acceleration}, \texttt{torque\ =\ force*length}.
Diese Felder sind Umrechnungsfaktoren, nicht physikalische Messwerte.

\texttt{Numeric::arcTan0to2piYoverX(double\ y,\ double\ x)} berechnet
einen Winkel im Bereich \([0,2\pi)\) aus den Komponenten \texttt{y/x};
\texttt{equaltol(double\ x,\ double\ xx,\ double\ tol)} prüft absolute
Differenz gegen \texttt{tol}; \texttt{isIncreasingVector(vector*)} prüft
Monotonie eines Wertevektors.
\texttt{MbDMath::xyzFromInt(int\ i)}/\texttt{XYZFromInt(int\ i)} bilden
einen Index auf Klein-/Großbuchstabenachsen ab; \texttt{noop()} ist
leer.

\texttt{SymbolicParser}, \texttt{Numeric}, \texttt{Units},
\texttt{CREATE} und die Matrix-/Vektorklassen sind Bibliotheksbausteine.
Die ASMT-Integration verwendet sie intern; direkte Clientnutzung setzt
konsistente Einheiten, dimensionstreue Vektoren und gültige Frameindizes
voraus.

\hypertarget{constraintset--frame--und-ableitungsfunktionen}{%
\section{\texorpdfstring{{15.18} ConstraintSet-, Frame- und
Ableitungsfunktionen}{15.18 ConstraintSet-, Frame- und Ableitungsfunktionen}}\label{constraintset--frame--und-ableitungsfunktionen}}

\hypertarget{itemij-constraintij-und-constraintset}{%
\subsection*{\texorpdfstring{\texttt{ItemIJ}, \texttt{ConstraintIJ} und
\texttt{ConstraintSet}}{ItemIJ, ConstraintIJ und ConstraintSet}}\label{itemij-constraintij-und-constraintset}}
\addcontentsline{toc}{subsection}{\texttt{ItemIJ}, \texttt{ConstraintIJ}
und \texttt{ConstraintSet}}

\texttt{ItemIJ()} und \texttt{ItemIJ(const\ std::string\&\ str)}
erzeugen ein unverbundenes beziehungsweise benanntes I/J-Objekt;
\texttt{ItemIJ(frmI,frmJ)} erhält zwei
\texttt{shared\_ptr\textless{}EndFramec\textgreater{}}.
\texttt{connectsItoJ(frmI,frmJ)} setzt diese beiden Framebezüge.
\texttt{ConstraintIJ(frmI,frmJ)} nimmt dieselben Frames und ergänzt
Constraintzustand.

\texttt{ConstraintSet()}/\texttt{ConstraintSet(str)} erzeugen den
Constraintcontainer;
\texttt{addConstraint(shared\_ptr\textless{}Constraint\textgreater{}\ con)}
fügt eine Gleichung hinzu; \texttt{constraintsDo(callback)} führt den
Callback \texttt{void(shared\_ptr\textless{}Constraint\textgreater{})}
für enthaltene Gleichungen aus; \texttt{constraintSpecs()} liefert deren
Spezifikationen. Die übrigen öffentlichen Methoden überschreiben
\texttt{Item}-Hooks und aggregieren ihre Kinder:

\begin{itemize}
\tightlist
\item
  \textbf{Constraintlisten:} \texttt{fillConstraints(all)},
  \texttt{fillDispConstraints(displacement)},
  \texttt{fillEssenConstraints(essential)},
  \texttt{fillPerpenConstraints(perpendicular)}.
\item
  \textbf{Fehler/Jacobians:} \texttt{fillAccICIterError(col)},
  \texttt{fillAccICIterJacob(mat)}, \texttt{fillPosICError(col)},
  \texttt{fillPosICJacob(mat)}, \texttt{fillPosKineError(col)},
  \texttt{fillPosKineJacob(mat)}, \texttt{fillVelICError(col)},
  \texttt{fillVelICJacob(mat)}.
\item
  \textbf{Zustand und Gewichte:} \texttt{fillqsuddotlam(col)},
  \texttt{fillqsulam(col)}, \texttt{fillqsudot(col)},
  \texttt{fillqsudotWeights(diagMat)}, \texttt{setqsuddotlam(col)},
  \texttt{setqsudotlam(col)}, \texttt{setqsulam(col)},
  \texttt{useEquationNumbers()}.
\item
  \textbf{Phasenhooks:} \texttt{initialize()},
  \texttt{initializeLocally()}, \texttt{initializeGlobally()},
  \texttt{postDynStep()}, \texttt{postInput()}, \texttt{postPosIC()},
  \texttt{postPosICIteration()}, \texttt{preAccIC()}, \texttt{preDyn()},
  \texttt{prePosIC()}, \texttt{prePosKine()}, \texttt{preVelIC()}.
\end{itemize}

\texttt{col} ist ein Ziel-/Quellvektor für Zustands-, Multiplikator-
oder Fehlerbeiträge; \texttt{mat} ist der Jacobianpuffer;
\texttt{diagMat} enthält Gewichtungen;
\texttt{all/displacement/essential/perpendicular} sind Zielvektoren der
jeweiligen Constraint-Unterklasse. Die Methoden ohne Parameter sind
Hookpunkte der Solvephasen.

\hypertarget{endframes-und-partframes}{%
\subsection*{Endframes und Partframes}\label{endframes-und-partframes}}
\addcontentsline{toc}{subsection}{Endframes und Partframes}

\texttt{EndFramec} konstruiert ein kartesisches Endframe ohne Namen oder
aus \texttt{str}; \texttt{setMarkerFrame(MarkerFrame*\ markerFrm)}
bindet es an einen Markerframe, \texttt{getMarkerFrame()} liest die
Zuordnung. \texttt{aAeO()} liefert die Orientierung in
Assemblykoordinaten; \texttt{aAjOe(size\_t\ j)} liefert
Spalten-/Achsenkomponente \texttt{j}; \texttt{riOeO(size\_t\ i)} liest
Positionskomponente \texttt{i}. \texttt{initEndFrameqct()},
\texttt{initEndFrameqct2()}, \texttt{initialize()},
\texttt{initializeLocally()}, \texttt{calcPostDynCorrectorIteration()}
und \texttt{root()} sind parameterlose Hooks. \texttt{rmeO()},
\texttt{rpep()}, \texttt{pAOppE()}, \texttt{aBOp()},
\texttt{newCopyEndFrameqc()} und \texttt{isEndFrameqc()} liefern
Positions-/Orientierungsableitungen oder einen abgeleiteten Frame.

\texttt{EndFrameqc} ergänzt \texttt{EndFramec} durch
\texttt{initializeGlobally()}, \texttt{initEndFrameqct()},
\texttt{initEndFrameqct2()}, \texttt{calcPostDynCorrectorIteration()},
\texttt{ppAjOepEpE(j)}, \texttt{pAjOepET(j)}, \texttt{ppriOeOpEpE(i)},
\texttt{iqX()}, \texttt{iqE()}, \texttt{priOeOpE(i)}, \texttt{qXdot()},
\texttt{qEdot()}, \texttt{qXddot()}, \texttt{qEddot()}, \texttt{rpep()},
\texttt{pAOppE()}, \texttt{aBOp()} und \texttt{isEndFrameqc()}.
\texttt{i/j} wählen Koordinaten- oder Orientierungsindex; die übrigen
Methoden liefern Ableitungsvektoren/-matrizen oder Gleichungsnummern.

\texttt{PartFrame} setzt/liest Partlage und -ableitungen über
\texttt{setqX/getqX}, \texttt{setqE/getqE}, \texttt{setqXdot/getqXdot},
\texttt{setomeOpO/getomeOpO/omeOpO}, \texttt{setqXddot/getqXddot},
\texttt{setqEddot/getqEddot}, \texttt{setaAap(mat)}. Es bindet über
\texttt{setPart(Part*)}/\texttt{getPart()},
\texttt{addMarkerFrame(frame)} und \texttt{endFrame(name)} Körper,
Marker und Endframe; \texttt{aGabsDo(callback)} und
\texttt{markerFramesDo(callback)} iterieren Constraints/Marker mit
passenden Callbacks. Lagegetter \texttt{rOpO}, \texttt{aAOp},
\texttt{aC}, \texttt{aCdot}, \texttt{alpOpO}, \texttt{pAOppE},
\texttt{aBOp} liefern Orts-/Orientierungs-/Ableitungsgrößen. Die übrigen
\texttt{fill*}, \texttt{set*}, \texttt{pre/post*}, Constraintreport- und
Zustandsfunktionen übernehmen geerbte \texttt{Item}-Pufferverträge;
\texttt{hnew} ist die vorgeschlagene Dynamikschrittweite.

\hypertarget{kinematicieje-ableitungen-einer-kinematischen-gruxf6uxdfe}{%
\subsection*{\texorpdfstring{\texttt{KinematicIeJe}: Ableitungen einer
kinematischen
Größe}{KinematicIeJe: Ableitungen einer kinematischen Größe}}\label{kinematicieje-ableitungen-einer-kinematischen-gruxf6uxdfe}}
\addcontentsline{toc}{subsection}{\texttt{KinematicIeJe}: Ableitungen
einer kinematischen Größe}

\texttt{KinematicIeJe(EndFrmsptr\ frmI,\ EndFrmsptr\ frmJ)} bindet zwei
Endframes; der parameterlose Konstruktor erzeugt die leere Größe.
\texttt{calc\_value()} berechnet den skalaren kinematischen Wert. Die
parameterlosen
\texttt{calc\_pvaluepX\{I,J,K\}}/\texttt{calc\_pvaluepE\{I,J,K\}} und
\texttt{calc\_ppvaluepX\{I,J,K\}pX\{I,J,K\}}-Familien berechnen erste
und zweite Ableitungen nach Translationskoordinaten \texttt{X} und
Eulerkoordinaten \texttt{E} der Frame-Seiten I/J/K. Die Varianten
\texttt{...pt()} und \texttt{...ptpt()} ergänzen Zeit- und gemischte
Zeit-/Koordinatenableitungen. Getter \texttt{pvaluepX*},
\texttt{pvaluepE*}, \texttt{ppvaluep*}, \texttt{pvaluept()},
\texttt{ppvalueptpt()} und \texttt{value()} liefern vorberechnete
Ableitungen/Wert. Diese Methoden sind Ableitungshooks für konkrete
Constraintgleichungen, keine direkten Eingaben für Clientcode.

\hypertarget{quaternionen-zeitabhuxe4ngige-frames-und-truxe4gheit}{%
\subsection*{Quaternionen, zeitabhängige Frames und
Trägheit}\label{quaternionen-zeitabhuxe4ngige-frames-und-truxe4gheit}}
\addcontentsline{toc}{subsection}{Quaternionen, zeitabhängige Frames und
Trägheit}

\texttt{EulerParameters\textless{}T\textgreater{}}-Konstruktoren
akzeptieren keine Eingabe (4 Komponenten), \texttt{count},
\texttt{(count,value)}, eine Werteliste oder
\texttt{(FColDsptr\ axis,double\ theta)}. Beim Achsenwinkelkonstruktor
ist \texttt{axis} eine 3D-Achse, \texttt{theta} ein Winkel in Radiant;
die Achse wird normiert. Die Klasse bietet \texttt{initialize()},
\texttt{calc()}, \texttt{calcABC()}, \texttt{calcpApE()},
\texttt{conditionSelf()}, \texttt{times(a)},
\texttt{plusFullColumn(eulerPars)}, \texttt{copy()} sowie
Ableitungsmatrizen \texttt{ppApEpEtimesColumn(col)},
\texttt{pCpEtimesColumn(col)}, \texttt{pCTpEtimesColumn(col)},
\texttt{ppApEpEtimesMatrix(mat)}. \texttt{col}/\texttt{mat} sind die zu
transformierenden Koordinaten-/Ableitungsdaten; die Methode liefert die
erste/zweite Rotationsableitung. Komponentenreihenfolge ist laut Header
\texttt{\{qE1,qE2,qE3,qE4\}}, mit skalarer Quaternionkomponente an Index
3.

\texttt{EulerAngles\textless{}T\textgreater{}} wird mit leerer
Dreierliste oder drei Winkelwerten erzeugt. \texttt{setRotOrder(i,j,k)}
setzt die drei Rotationsachsen in Reihenfolge; zulässig sind Achsen 1,
2, 3 ohne gleiche benachbarte Achsen. \texttt{initialize()},
\texttt{calc()} berechnen Drehmatrizen; \texttt{differentiateWRT(var)}
liefert die Winkelableitung nach der angegebenen
symbolischen/numerischen Variable. Nicht unterstützte generische
\texttt{T}-Instanziierungen werfen \texttt{SimulationStoppingError}.

\texttt{EndFrameqct} stellt zeitabhängige
Orientierungs-/Positionsfunktionen bereit: \texttt{initprmemptBlks()},
\texttt{initpprmemptptBlks()}, \texttt{initpPhiThePsiptBlks()},
\texttt{initppPhiThePsiptptBlks()}, \texttt{evalrmem()},
\texttt{evalAme()}, \texttt{evalprmempt()}, \texttt{evalpAmept()},
\texttt{evalpprmemptpt()}, \texttt{evalppAmeptpt()},
\texttt{pAjOept(j)}, \texttt{ppAjOepETpt(j)}, \texttt{ppAjOeptpt(j)},
\texttt{priOeOpt(i)}, \texttt{ppriOeOpEpt(i)}, \texttt{ppriOeOptpt(i)},
sowie \texttt{initialize*}, \texttt{prePosIC}, \texttt{preVelIC},
\texttt{postVelIC}, \texttt{preAccIC}, \texttt{postInput},
\texttt{calcPostDynCorrectorIteration}, \texttt{rmeO}, \texttt{rpep},
\texttt{isEndFrameqc}. \texttt{i/j} sind Frame-/Achsenindizes;
\texttt{time} ist aktuelle Solverzeit. \texttt{EndFrameqct2}
überschreibt die sechs
\texttt{init*PhiThePsi*}-/\texttt{eval*Ame*}-Hooks für eine alternative
Winkelparameterisierung; seine Konstruktoren erhalten optional einen
Namen.

\texttt{MomentOfInertiaSolver} spezialisiert Matrixzerlegung für
Trägheitstensoren. \texttt{setm(double\ mass)} setzt Körpermasse;
\texttt{setJPP(FMatDsptr\ mat)} Haupt-/Bezugs-Trägheitstensor;
\texttt{setrPoP(FColDsptr\ col)} und \texttt{setrPcmP(col)}
Vektorverschiebungen; \texttt{setAPo(FMatDsptr\ mat)} Orientierung.
\texttt{calc()}, \texttt{calcJoo()}, \texttt{calcJpp()},
\texttt{calcAPp()}, \texttt{calcJppFromDiagJcmP()},
\texttt{calcJppFromFullJcmP()} berechnen Tensoren/Frameabbildung;
\texttt{getJoo()}, \texttt{getJpp()}, \texttt{getAPp()} geben Ergebnisse
zurück; \texttt{eigenvectorFor(double\ lam)} liefert Eigenvektor zum
Eigenwert \texttt{lam}; \texttt{modifiedArcCos(val)} verarbeitet
numerisch einen Kosinuswert. Die überschriebenen
\texttt{preSolvewithsaveOriginal(matrix,rhs,saveOriginal)},
\texttt{basicSolvewithsaveOriginal(...)}, \texttt{doPivoting(p)},
\texttt{doFullPivoting(p)}, \texttt{forwardEliminateWithPivot(p)},
\texttt{backSubstituteIntoDU()}, \texttt{postSolve()} und
\texttt{getmatrixArowimaxMagnitude(i)} erben die
MatrixSolver-Parameterrollen aus
\protect\hyperlink{abs-api-matrizen}{Abschnitt 15.10}.

\texttt{PosVelAccData} bündelt \texttt{rFfF} (Ort), \texttt{aAFf}
(Orientierung), \texttt{vFfF}/\texttt{omeFfF}
(lineare/Winkelgeschwindigkeit) und \texttt{aFfF}/\texttt{alpFfF}
(lineare/Winkelbeschleunigung); \texttt{refData} verknüpft Vorzustand.
\texttt{printOn(stream)} schreibt den Container.
\texttt{CartesianFrame()}/\texttt{CartesianFrame(str)} erzeugt ein
kartesisches Frame, \texttt{initialize()} richtet es ein;
\texttt{Orientation()}/\texttt{Orientation(str)} erzeugt eine
Prescribed-Motion-Orientierung und \texttt{initializeGlobally()} richtet
sie für das System ein.

\hypertarget{anhang}{%
\chapter{Anhang}\label{anhang}}

\hypertarget{anh-asmt-format}{%
\chapter{\texorpdfstring{{Anhang A:} ASMT-Format: dokumentierter Stand
am
30.09.2026}{Anhang A: ASMT-Format: dokumentierter Stand am 30.09.2026}}\label{anh-asmt-format}}

\hypertarget{kurzantwort-wo-ist-das-format-beschrieben}{%
\section{\texorpdfstring{{A.1} Kurzantwort: Wo ist das Format
beschrieben?}{A.1 Kurzantwort: Wo ist das Format beschrieben?}}\label{kurzantwort-wo-ist-das-format-beschrieben}}

Eine aktuelle, normative und versionsgebundene ASMT-Spezifikation konnte
ich nicht finden. Die beste Beschreibung des derzeitigen Formats ergibt
sich aus drei Ebenen:

\begin{enumerate}
\def\labelenumi{\arabic{enumi}.}
\tightlist
\item
  \textbf{Historische Primärbeschreibung von Ondsel:} Im
  \href{https://www.ondsel.com/blog/assembly-wb-prerelease/}{Ondsel-Artikel
  zur frühen Assembly-Workbench} steht im Abschnitt „ASMT files``, dass
  ASMT ein einfaches Plain-Text-Format zur Darstellung der
  Assembly-Struktur und damals vor allem ein internes Debugging-Werkzeug
  war. Der Artikel sagt ausdrücklich, dass die Funktion später entfernt
  oder reduziert werden könnte.
\item
  \textbf{Aktuelle Implementierung und Beispiele:} Maßgeblich für den
  Stand dieses Buches sind der Parser und Serializer im Submodul
  \href{vendor/OndselSolver/OndselSolver/ASMTAssembly.cpp}{ASMTAssembly.cpp},
  die Basishilfen in
  \href{vendor/OndselSolver/OndselSolver/ASMTItem.cpp}{ASMTItem.cpp}
  sowie die echten Dateien in
  \href{./examples/07-ondselsolver-real-asmt/README.md}{\texttt{examples/07-ondselsolver-real-asmt}}.
  Der historische Upstream-Stand ist zusätzlich im
  \href{https://github.com/Ondsel-Development/OndselSolver/tree/09d6175a}{archivierten
  OndselSolver-Repository} einsehbar.
\item
  \textbf{FreeCAD-Integrationshinweis:} FreeCAD registriert
  \texttt{.asmt} als Importtyp. Der Kommentar in
  \href{https://github.com/FreeCAD/FreeCAD/blob/main/src/Mod/Assembly/Init.py}{Assembly/Init.py}
  sagt aber ausdrücklich, dass das korrekte Interoperabilitätsformat
  noch Forschungsgegenstand und ASMT ein Platzhalter ist. Der aktuelle
  \href{https://github.com/FreeCAD/FreeCAD/blob/main/src/Mod/Assembly/AssemblyImport.py}{AssemblyImport.py}
  erstellt noch keine Assembly aus der Datei.
\end{enumerate}

Die
\href{https://deepwiki.com/FreeCAD/OndselSolver/4.4-asmt-file-format}{DeepWiki-Seite
zum ASMT-Format} und die
\href{https://deepwiki.com/Ondsel-Development/OndselSolver/4-file-formats-and-exchange}{DeepWiki-Übersicht
zu Dateiformaten} sind nützliche, quellenverlinkte
Sekundärzusammenfassungen. Sie sind keine Herausgeber-Spezifikation;
ihre Aussagen müssen gegen den jeweiligen Solver-Quellstand geprüft
werden. Auch das
\href{https://github.com/deepsaia/pyondsel}{pyondsel-Projekt} beschreibt
die Nutzung des nativen OndselSolver-ASMT, definiert aber keinen
unabhängigen Standard.

\textbf{Fazit:} ASMT ist ein tatsächlich verwendetes, lesbares
Solver-Datenformat, aber nach dem hier geprüften Stand kein
abgeschlossener, herstellerneutraler Dateistandard mit veröffentlichter
normativer Grammatik.

\hypertarget{wozu-die-datei-dient}{%
\section{\texorpdfstring{{A.2} Wozu die Datei
dient}{A.2 Wozu die Datei dient}}\label{wozu-die-datei-dient}}

ASMT serialisiert eine Solver-Assembly: Körper, ihre Anfangslagen und
Bezugssysteme, Marker, Gelenke, Bewegungen, Grenzen, Kräfte sowie
Simulationsparameter. Solver-Ausgaben können zusätzlich Zeitreihen und
berechnete Zustände enthalten.

Es ist nicht automatisch ein vollständiger CAD-Projektcontainer. Die
Fixtures in Beispiel 07 enthalten Solver-Geometrie und Kinematikdaten,
aber keine vollständigen FreeCAD-B-Reps. Außerdem registriert FreeCAD
zwar den Dateityp, der derzeitige Import-Stub erzeugt daraus noch keine
FreeCAD-Objekte.

\hypertarget{aufbau-im-aktuellen-parser}{%
\section{\texorpdfstring{{A.3} Aufbau im aktuellen
Parser}{A.3 Aufbau im aktuellen Parser}}\label{aufbau-im-aktuellen-parser}}

\hypertarget{dateikopf-und-hierarchie}{%
\subsection*{Dateikopf und Hierarchie}\label{dateikopf-und-hierarchie}}
\addcontentsline{toc}{subsection}{Dateikopf und Hierarchie}

Der Parser akzeptiert zwei Kopfzeilen:

\begin{Shaded}
\begin{Highlighting}[]
\NormalTok{OndselSolver}
\NormalTok{Assembly}
\end{Highlighting}
\end{Shaded}

oder den älteren Kopf:

\begin{Shaded}
\begin{Highlighting}[]
\NormalTok{freeCAD: 3D CAD with Motion Simulation  by  askoh.com}
\NormalTok{Assembly}
\end{Highlighting}
\end{Shaded}

Darauf folgt ein hierarchischer Textbaum. Im aktuellen Parser markieren
\textbf{Tabulatoren} die Verschachtelung. Schlüsselwörter und Werte
stehen auf getrennten Zeilen; numerische Vektoren werden als mehrere
durch Whitespace getrennte Zahlen gelesen. Leerzeichen anstelle von Tabs
sind kein verlässlich unterstützter Ersatz: Der Parser vergleicht an
vielen Stellen exakte, eingerückte Schlüsselzeilen.

Ein reales minimales Ausschnittmuster aus \texttt{\_\_cubes.asmt}:

\begin{Shaded}
\begin{Highlighting}[]
\NormalTok{OndselSolver}
\NormalTok{Assembly}
\NormalTok{    Notes}
\NormalTok{        (Text string: \textquotesingle{}\textquotesingle{} runs: (Core.RunArray runs: \#() values: \#()))}
\NormalTok{    Name}
\NormalTok{        OndselAssembly}
\NormalTok{    Position3D}
\NormalTok{        0 0 0}
\NormalTok{    RotationMatrix}
\NormalTok{        1 0 0}
\NormalTok{        0 1 0}
\NormalTok{        0 0 1}
\end{Highlighting}
\end{Shaded}

Das ist ein gekürzter Leseausschnitt, keine vollständige gültige Datei.
Eine vollständige, vom lokalen Solver getestete Eingabe liegt unter
\href{./examples/07-ondselsolver-real-asmt/__cubes.asmt}{\texttt{\_\_cubes.asmt}}.

\hypertarget{root-assembly}{%
\subsection*{Root-Assembly}\label{root-assembly}}
\addcontentsline{toc}{subsection}{Root-Assembly}

\texttt{ASMTAssembly::parseASMT()} liest die Root-Daten in fester
Reihenfolge ein:

\begin{longtable}[]{@{}
  >{\raggedright\arraybackslash}p{(\columnwidth - 2\tabcolsep) * \real{0.5000}}
  >{\raggedright\arraybackslash}p{(\columnwidth - 2\tabcolsep) * \real{0.5000}}@{}}
\toprule\noalign{}
\begin{minipage}[b]{\linewidth}\raggedright
Abschnitt
\end{minipage} & \begin{minipage}[b]{\linewidth}\raggedright
Bedeutung
\end{minipage} \\
\midrule\noalign{}
\endhead
\bottomrule\noalign{}
\endlastfoot
\texttt{Notes}, \texttt{Name} & Freitextnotiz und Assembly-Name \\
\texttt{Position3D}, \texttt{RotationMatrix} & Lage des
Assembly-Bezugssystems \\
\texttt{Velocity3D}, \texttt{Omega3D} & lineare und
Winkelgeschwindigkeit \\
\texttt{RefPoints}, \texttt{RefCurves}, \texttt{RefSurfaces} &
geometrische Referenzobjekte und Markercontainer \\
\texttt{Parts} & enthaltene Körper und deren Eigenschaften \\
\texttt{KinematicIJs} & kinematische Datenobjekte \\
\texttt{ConstraintSets} & \texttt{Joints}, \texttt{Motions},
\texttt{Limits}, \texttt{GeneralConstraintSets} \\
\texttt{ForceTorques}, \texttt{ConstantGravity} & Belastung und
Schwerkraft \\
\texttt{SimulationParameters}, \texttt{AnimationParameters} &
Solve-/Integrations- und Animationssteuerung \\
\texttt{TimeSeries} und nachfolgende Serien & optionale berechnete
Zeitverläufe \\
\end{longtable}

Die Reihenfolge ist eine Implementierungserwartung, keine frei
sortierbare Schlüssel-Wert-Sammlung. Das sieht man unmittelbar an der
fest sequenzierten Aufrufkette in \texttt{ASMTAssembly::parseASMT()}.

\hypertarget{kuxf6rper-marker-und-gelenke}{%
\subsection*{Körper, Marker und
Gelenke}\label{kuxf6rper-marker-und-gelenke}}
\addcontentsline{toc}{subsection}{Körper, Marker und Gelenke}

Ein \texttt{Part} enthält unter anderem Name, Lage, Rotation,
Geschwindigkeiten, \texttt{FeatureOrder}, \texttt{PrincipalMassMarker}
und Referenzpunkte/-kurven/-flächen. Ein \texttt{Marker} speichert Name,
lokale Position und Orientierung; Gelenke verbinden zwei Marker über
Referenzpfade, zum Beispiel \texttt{MarkerI} und \texttt{MarkerJ}.

Der Parser der räumlichen Daten erwartet für \texttt{Position3D} in der
Praxis drei Zahlen. \texttt{RotationMatrix} wird als drei Zeilen
eingelesen, typischerweise mit je drei Zahlen. Die aktuelle
Implementierung prüft die Dimensionen jedoch nicht überall robust;
beschädigte oder falsch eingerückte Dateien können deshalb Assertions,
Exceptions oder im ungünstigen Fall Speicherfehler auslösen.

Der Joint-Parser kennt im aktuellen Quellstand unter anderem
\texttt{FixedJoint}, \texttt{RevoluteJoint}, \texttt{SphericalJoint},
\texttt{CylindricalJoint}, \texttt{TranslationalJoint},
\texttt{UniversalJoint}, \texttt{PlanarJoint}, \texttt{AngleJoint},
\texttt{GearJoint}, \texttt{ScrewJoint} und mehrere zusammengesetzte
Gelenke. Bewegungen umfassen \texttt{RotationalMotion},
\texttt{TranslationalMotion}, \texttt{GeneralMotion} und
\texttt{AllowRotation}; Grenzen umfassen \texttt{RotationLimit} und
\texttt{TranslationLimit}. Welche Parameter je Typ nötig sind, ergibt
sich aus der jeweiligen \texttt{ASMT*}-Klasse und ihren Fixtures.

\hypertarget{einheiten-und-zahlen}{%
\subsection*{Einheiten und Zahlen}\label{einheiten-und-zahlen}}
\addcontentsline{toc}{subsection}{Einheiten und Zahlen}

Die Werte werden als numerische Textwerte eingelesen. In den üblichen
Abschnitten steht kein explizites Einheiten-Token neben jeder Zahl.
OndselSolver besitzt eine separate \texttt{Units}-Struktur für Zeit,
Masse, Länge, Winkel und abgeleitete Größen; deren Standardfaktoren sind
1.0. Bei selbst erzeugten oder extern konvertierten Dateien muss daher
die verwendete Einheitenskala bekannt und konsistent sein.

\hypertarget{ausgabe-und-solver-ergebnis}{%
\subsection*{Ausgabe und
Solver-Ergebnis}\label{ausgabe-und-solver-ergebnis}}
\addcontentsline{toc}{subsection}{Ausgabe und Solver-Ergebnis}

\texttt{ASMTAssembly::outputFile()} schreibt im aktuellen Quellstand
stets den Kopf \texttt{OndselSolver} und serialisiert danach den
Objektbaum. Die Serialisierung verwendet Tabs für Ebenen und schreibt
Fließkommazahlen mit hoher Genauigkeit (\texttt{max\_digits10}). Nach
einer Simulation können Zeitreihenabschnitte hinzukommen.

\hypertarget{parsergrenzen-nicht-mit-format-unterstuxfctzt-verwechseln}{%
\section{\texorpdfstring{{A.4} Parsergrenzen: nicht mit „Format
unterstützt``
verwechseln}{A.4 Parsergrenzen: nicht mit „Format unterstützt`` verwechseln}}\label{parsergrenzen-nicht-mit-format-unterstuxfctzt-verwechseln}}

Die Existenz eines Abschnittsnamens in einer Datei bedeutet nicht, dass
jede Variante unterstützt wird. Im untersuchten Quellstand gilt unter
anderem:

\begin{itemize}
\tightlist
\item
  Nichtleere \texttt{KinematicIJs} enden in \texttt{readKinematicIJ()}
  mit „To be implemented``.
\item
  \texttt{RefCurves} und \texttt{RefSurfaces} werden als Container
  erkannt, ihre Einzelelement-Parser werfen aber ebenfalls „To be
  implemented``.
\item
  Nichtleere \texttt{GeneralConstraintSets} sind noch nicht
  implementiert.
\item
  \texttt{ASMTItem::parseASMT()} ist für nicht speziell behandelte
  Objekttypen ein Fehlerpfad.
\item
  Viele Parserpfade verlassen sich auf Assertions und exakte Zeilenfolge
  statt auf eine fehlertolerante Validierung.
\end{itemize}

Damit beschreibt dieser Appendix den gegenwärtigen
Parser/Serializer-Umfang, nicht eine Garantie, dass alle denkbaren
ASMT-Dateien oder alle historischen Varianten unterstützt werden.

\hypertarget{reproduzierbare-referenzdateien}{%
\section{\texorpdfstring{{A.5} Reproduzierbare
Referenzdateien}{A.5 Reproduzierbare Referenzdateien}}\label{reproduzierbare-referenzdateien}}

Beispiel 07 enthält drei lokale Eingabedateien, die mit dem echten
OndselSolver getestet werden:

\begin{itemize}
\tightlist
\item
  \href{./examples/07-ondselsolver-real-asmt/__cubes.asmt}{\texttt{\_\_cubes.asmt}}:
  zwei Körper, Fixed- und Angle-Joint
\item
  \href{./examples/07-ondselsolver-real-asmt/fourbar.asmt}{\texttt{fourbar.asmt}}:
  Viergelenk
\item
  \href{./examples/07-ondselsolver-real-asmt/piston.asmt}{\texttt{piston.asmt}}:
  Kurbel-Schieber
\end{itemize}

Diese Dateien sind besonders nützlich als praktische Formatbeispiele,
aber nicht als vollständige Schema-Spezifikation. Der Parser selbst und
seine Testfixtures bleiben der genaueste Nachweis dafür, was der lokale
Solverstand akzeptiert.

\hypertarget{wie-freecad-die-asmt-struktur-erzeugt}{%
\section{\texorpdfstring{{A.6} Wie FreeCAD die ASMT-Struktur
erzeugt}{A.6 Wie FreeCAD die ASMT-Struktur erzeugt}}\label{wie-freecad-die-asmt-struktur-erzeugt}}

Der FreeCAD-Export ist keine generische Dokument-Serialisierung. Der
Aufruf aus
\href{https://github.com/FreeCAD/FreeCAD/blob/main/src/Mod/Assembly/CommandExportASMT.py}{\texttt{CommandExportASMT.py}}
endet in \texttt{AssemblyObject::exportAsASMT()}. Der relevante Ablauf
in
\href{https://github.com/FreeCAD/FreeCAD/blob/main/src/Mod/Assembly/App/AssemblyObject.cpp}{\texttt{AssemblyObject.cpp}}
lautet:

\begin{enumerate}
\def\labelenumi{\arabic{enumi}.}
\tightlist
\item
  \texttt{makeMbdAssembly()} erzeugt ein \texttt{ASMTAssembly}, setzt
  dessen Solver-Rootnamen auf \texttt{OndselAssembly} und verknüpft es
  mit dem FreeCAD-Assembly-Objekt.
\item
  \texttt{rebuildRigidClusters()} und \texttt{fixGroundedParts()}
  bereiten die Solver-Körper vor. Für Grounding erzeugt FreeCAD einen
  Marker auf der Assembly an der Zielplatzierung, einen lokalen
  \texttt{FixingMarker} am Part und ein \texttt{ASMTFixedJoint} zwischen
  beiden.
\item
  \texttt{getJoints()} sammelt die Baugruppen-Joints.
  \texttt{jointParts()} ruft für jeden Joint \texttt{makeMbdJoint()} auf
  und fügt das erzeugte Solver-Objekt in die Assembly ein.
\item
  \texttt{getMbDData()} ordnet FreeCAD-Objekte den
  \texttt{ASMTPart}-Objekten zu. Bei Links und verschachtelten
  Baugruppen werden Identität und Container-Placements aufgelöst;
  mehrere FreeCAD-Objekte können über \texttt{objectPartMap} auf
  denselben Solver-Körper zeigen.
\item
  \texttt{handleOneSideOfJoint()} nimmt die Referenz und Placement jeder
  Joint-Seite, berechnet daraus einen part-lokalen Marker und gibt
  dessen ASMT-Pfad zurück. Der Joint verweist anschließend auf
  \texttt{MarkerI} und \texttt{MarkerJ}.
\item
  \texttt{makeMbdJointOfType()} übersetzt FreeCAD-Jointtypen in konkrete
  ASMT-Klassen, zum Beispiel \texttt{Fixed} → \texttt{ASMTFixedJoint},
  \texttt{Revolute} → \texttt{ASMTRevoluteJoint}, \texttt{Cylindrical} →
  \texttt{ASMTCylindricalJoint}, \texttt{Slider} →
  \texttt{ASMTTranslationalJoint} und \texttt{Ball} →
  \texttt{ASMTSphericalJoint}. Typabhängige Parameter wie Winkel,
  Gear-Radien oder Schraubensteigung werden dabei ebenfalls übertragen.
\item
  \texttt{outputFile()} serialisiert die fertige Datenstruktur als
  Textdatei.
\end{enumerate}

Als Ablaufdiagramm:

\begin{Shaded}
\begin{Highlighting}[]
\NormalTok{FreeCAD AssemblyObject}
\NormalTok{    {-}\textgreater{} ASMTAssembly("OndselAssembly")}
\NormalTok{    {-}\textgreater{} ASMTPart + MassMarker + Placement/Rotation}
\NormalTok{    {-}\textgreater{} Ground{-}Marker + FixingMarker + FixedJoint}
\NormalTok{    {-}\textgreater{} Joint{-}Marker aus Reference1/Reference2 und Placement1/Placement2}
\NormalTok{    {-}\textgreater{} ASMT{-}Jointklasse samt Parametern}
\NormalTok{    {-}\textgreater{} outputFile(.asmt)}
\end{Highlighting}
\end{Shaded}

Die Zuordnung zwischen FreeCAD-Code und Datei lässt sich so lesen:

\begin{longtable}[]{@{}
  >{\raggedright\arraybackslash}p{(\columnwidth - 4\tabcolsep) * \real{0.3333}}
  >{\raggedright\arraybackslash}p{(\columnwidth - 4\tabcolsep) * \real{0.3333}}
  >{\raggedright\arraybackslash}p{(\columnwidth - 4\tabcolsep) * \real{0.3333}}@{}}
\toprule\noalign{}
\begin{minipage}[b]{\linewidth}\raggedright
FreeCAD-Methode
\end{minipage} & \begin{minipage}[b]{\linewidth}\raggedright
ASMT-Objekt/Block
\end{minipage} & \begin{minipage}[b]{\linewidth}\raggedright
Konkrete Wirkung
\end{minipage} \\
\midrule\noalign{}
\endhead
\bottomrule\noalign{}
\endlastfoot
\texttt{makeMbdAssembly()} & \texttt{Assembly} & legt die
Solver-Rootassembly \texttt{OndselAssembly} an \\
\texttt{fixGroundedPart()} & Root-Marker, Part-Marker,
\texttt{FixedJoint} & verankert ein Teil über einen Assembly-Marker und
den lokalen \texttt{FixingMarker} \\
\texttt{getMbdData()} / \texttt{makeMbdPart()} & \texttt{Part},
\texttt{Position3D}, \texttt{RotationMatrix},
\texttt{PrincipalMassMarker} & mappt ein FreeCAD-Part/Link auf einen
Solver-Körper; Placement wird in Translation und Rotationsmatrix
zerlegt \\
\texttt{handleOneSideOfJoint()} / \texttt{makeMbdMarker()} &
\texttt{RefPoint} / \texttt{Marker} & berechnet die joint-lokale
Markerposition und hängt sie an den zugehörigen Solver-Körper \\
\texttt{makeMbdJoint()} / \texttt{makeMbdJointOfType()} &
\texttt{ConstraintSets/Joints/\textless{}JointType\textgreater{}} &
wählt die Solver-Jointklasse, setzt Parameter und referenziert beide
Markerpfade \\
\texttt{outputFile()} / \texttt{storeOnLevel()} & Textdatei & schreibt
den Solver-Objektbaum mit Tab-Einrückung aus \\
\end{longtable}

\hypertarget{reales-joint-beispiel-aus-freecadondselsolver}{%
\subsection*{Reales Joint-Beispiel aus
FreeCAD/OndselSolver}\label{reales-joint-beispiel-aus-freecadondselsolver}}
\addcontentsline{toc}{subsection}{Reales Joint-Beispiel aus
FreeCAD/OndselSolver}

Dieser unveränderte Ausschnitt stammt aus
\href{./examples/07-ondselsolver-real-asmt/__cubes.asmt}{\texttt{\_\_cubes.asmt}}.
Er zeigt sowohl die erzeugte Erdung als auch ein Winkelgelenk. Die Pfade
unter \texttt{MarkerI} und \texttt{MarkerJ} verweisen auf die zuvor
serialisierten Marker:

\begin{Shaded}
\begin{Highlighting}[]
\NormalTok{    ConstraintSets}
\NormalTok{        Joints}
\NormalTok{            FixedJoint}
\NormalTok{                Name}
\NormalTok{                    \_\_cubes\#GroundedJoint}
\NormalTok{                MarkerI}
\NormalTok{                    /OndselAssembly/marker{-}\_\_cubes\#Box001}
\NormalTok{                MarkerJ}
\NormalTok{                    /OndselAssembly/\_\_cubes\#Box001/FixingMarker}
\NormalTok{            AngleJoint}
\NormalTok{                Name}
\NormalTok{                    \_\_cubes\#Angle}
\NormalTok{                MarkerI}
\NormalTok{                    /OndselAssembly/\_\_cubes\#Box001/\_\_cubes\#Angle}
\NormalTok{                MarkerJ}
\NormalTok{                    /OndselAssembly/\_\_cubes\#Box002/\_\_cubes\#Angle}
\NormalTok{                theIzJz}
\NormalTok{                    0.17453292519943295}
\end{Highlighting}
\end{Shaded}

\texttt{theIzJz} ist hier der Winkelparameter in Radiant. Der Wert
\texttt{0.17453292519943295} entspricht 10 Grad. Entscheidend ist die
Datenkette: FreeCAD-Joint → ASMT-Jointklasse → zwei serialisierte
Markerpfade plus Gelenkparameter.

\hypertarget{was-wird-dabei-nicht-exportiert}{%
\subsection*{Was wird dabei nicht
exportiert?}\label{was-wird-dabei-nicht-exportiert}}
\addcontentsline{toc}{subsection}{Was wird dabei nicht exportiert?}

Der Pfad \texttt{makeMbdPart()} schreibt Position und
3x3-Rotationsmatrix, erzeugt einen PrincipalMassMarker und setzt im
normalen Aufruf Standardwerte für Masse (\texttt{1.0}), Dichte
(\texttt{1.0}) und Trägheitsmomente (\texttt{1.0,\ 1.0,\ 1.0}). Er
serialisiert damit nicht automatisch CAD-B-Reps oder aus der Form
berechnete reale Masseneigenschaften. ASMT ist an dieser Stelle ein
Solver-Modell, nicht ein Ersatz für die vollständige FreeCAD-Datei.

\texttt{exportAsASMT()} baut und schreibt die ASMT-Struktur; der
Exportpfad ruft selbst nicht \texttt{solve()} auf. Lösen und ASMT-Export
sind getrennte Operationen. Ebenso registriert FreeCAD \texttt{.asmt}
zwar als Dateityp, aber der aktuelle
\href{https://github.com/FreeCAD/FreeCAD/blob/main/src/Mod/Assembly/AssemblyImport.py}{AssemblyImport.py}-Stub
importiert noch keine Solver-Struktur in ein FreeCAD-Dokument.

\hypertarget{quellen-und-stand}{%
\section{\texorpdfstring{{A.7} Quellen und
Stand}{A.7 Quellen und Stand}}\label{quellen-und-stand}}

\begin{itemize}
\tightlist
\item
  Ondsel,
  \href{https://www.ondsel.com/blog/assembly-wb-prerelease/}{„Christmas
  comes early! A pre-release of the integrated assembly workbench``},
  01.12.2025, Abschnitt „ASMT files``: historische Primärbeschreibung
  als internes, einfaches Plain-Text-/Debugformat.
\item
  \href{https://github.com/Ondsel-Development/OndselSolver/blob/09d6175a/OndselSolver/ASMTAssembly.cpp}{OndselSolver,
  ASMTAssembly.cpp} und
  \href{https://github.com/Ondsel-Development/OndselSolver/blob/09d6175a/OndselSolver/ASMTItem.cpp}{ASMTItem.cpp}:
  Parser- und Serializer-Implementierung im archivierten Upstream-Stand.
\item
  FreeCAD,
  \href{https://github.com/FreeCAD/FreeCAD/blob/main/src/Mod/Assembly/Init.py}{Assembly/Init.py}:
  Kommentar, dass ein geeignetes Interoperabilitätsformat noch
  Forschungsgegenstand ist und ASMT als Platzhalter behandelt wird.
\item
  DeepWiki,
  \href{https://deepwiki.com/FreeCAD/OndselSolver/4.4-asmt-file-format}{ASMT
  File Format} und
  \href{https://deepwiki.com/Ondsel-Development/OndselSolver/4-file-formats-and-exchange}{File
  Formats and Exchange}: Sekundärzusammenfassungen mit Quellverweisen;
  keine normative Spezifikation.
\item
  Solver-Quellstand am \textbf{30.09.2026}:
  \href{vendor/OndselSolver/OndselSolver/ASMTAssembly.cpp}{ASMTAssembly.cpp},
  \href{vendor/OndselSolver/OndselSolver/ASMTPart.cpp}{ASMTPart.cpp},
  \href{vendor/OndselSolver/OndselSolver/ASMTSpatialItem.cpp}{ASMTSpatialItem.cpp},
  \href{vendor/OndselSolver/OndselSolver/ASMTAssembly.h}{ASMTAssembly.h}
  und \href{vendor/OndselSolver/OndselSolver/Units.h}{Units.h} im auf
  Commit \texttt{4be80eef} fixierten Submodul.
\item
  FreeCAD-Referenzquellen:
  \href{https://github.com/FreeCAD/FreeCAD/blob/main/src/Mod/Assembly/App/AssemblyObject.cpp}{AssemblyObject.cpp},
  \href{https://github.com/FreeCAD/FreeCAD/blob/main/src/Mod/Assembly/CommandExportASMT.py}{CommandExportASMT.py}
  und
  \href{https://github.com/FreeCAD/FreeCAD/blob/main/src/Mod/Assembly/AssemblyImport.py}{AssemblyImport.py}.
\end{itemize}

\backmatter
\end{document}
